% La Chine : de la révolution culturelle à l’émergence économique — Histoire Tle Tech — Cours signé OnBuch+
% Programme officiel MINESEC. Compiler avec : tectonic N02-chine-revolution-culturelle-emergence-economique.tex
\newcommand{\DOCMATIERE}{Histoire}
\newcommand{\DOCNIVEAU}{Tle Tech}
\newcommand{\DOCMODULE}{Histoire – classe de Terminale technique}
\newcommand{\DOCLECON}{N02}
\newcommand{\DOCTITRE}{La Chine : de la révolution culturelle à l’émergence économique}
% OnBuch+ — préambule « cours signés » ENRICHI (compilé avec tectonic / XeLaTeX)
% Système visuel « Pop » d'OnBuch+ : fond crème (jamais blanc), bordures encre
% épaisses, ombres « dures » décalées, titres Archivo Black en capitales.
% Version MATHÉMATIQUES : graphes (pgfplots/tikz), boîtes pédagogiques
% (définition, propriété, méthode, attention, le savais-tu, exercice résolu,
% exercice, TP, bilan), page de garde et signature OnBuch+ ancrée en bas.
%
% Macros attendues dans main.tex AVANT \input{preamble} :
%   \DOCMATIERE  \DOCNIVEAU  \DOCMODULE  \DOCLECON (numéro)  \DOCTITRE
\documentclass[11pt]{article}

\usepackage{fontspec}
\usepackage[french]{babel}
\usepackage[a4paper,margin=2cm,top=3.4cm,bottom=2.8cm,headheight=1.4cm,footskip=1.3cm]{geometry}
\usepackage{amsmath,amssymb,amsfonts}
\usepackage{mathtools} % \xmapsto, \coloneqq... souvent utilisés par les modèles
\usepackage{cancel} % simplifications barrées (bilans de forces)
% \abs{z} (module, valeur absolue) et \norm{u} (norme), utilisés par les modèles
\providecommand{\abs}{}\let\abs\relax\DeclarePairedDelimiter{\abs}{\lvert}{\rvert}
\providecommand{\norm}{}\let\norm\relax\DeclarePairedDelimiter{\norm}{\lVert}{\rVert}
\usepackage[table]{xcolor} % [table] : fournit aussi \rowcolors (alternance de lignes)
\usepackage{pagecolor}
\usepackage{tikz}
\usetikzlibrary{arrows.meta,positioning,calc,shapes.geometric,shapes.misc,shapes.arrows,decorations.pathmorphing,decorations.pathreplacing,decorations.markings,patterns,shadows,mindmap,trees,fit,backgrounds,matrix,intersections,angles,quotes}
\usepackage{pgfplots}
\usepackage[version=4]{mhchem} % équations nucléaires \ce{^{235}_{92}U}
\usepackage[european,siunitx]{circuitikz} % schémas électriques
\pgfplotsset{compat=1.18}
\usepgfplotslibrary{fillbetween,groupplots} % aires entre courbes (calcul intégral)
\usepackage{enumitem}
\usepackage{fancyhdr}
% [most] charge toutes les bibliothèques tcolorbox (listings, minted, raster,
% documentation...) et gonfle inutilement la mémoire TeX sur un document long
% (~300 boîtes) : on ne charge que ce qui est réellement utilisé.
\usepackage[skins,breakable]{tcolorbox}
\usepackage{graphicx}
\usepackage{colortbl}
\usepackage{booktabs}
\usepackage{tabularx}
\usepackage{multirow}
\usepackage{diagbox} % case d'en-tête diagonale des tableaux à double entrée
% Colonnes extensibles centrée / gauche / droite (C, L, R), souvent utilisées par les modèles
\newcolumntype{C}{>{\centering\arraybackslash}X}
\newcolumntype{L}{>{\raggedright\arraybackslash}X}
\newcolumntype{R}{>{\raggedleft\arraybackslash}X}
\usepackage{array}
\usepackage{multicol}
\usepackage{siunitx}
% parse-numbers=false : accepte aussi \SI{2,5 \times 10^{-3}}{...}
\sisetup{locale=FR,output-decimal-marker={,},group-minimum-digits=4,parse-numbers=false}
\DeclareSIUnit\ohm{\ensuremath{\symup{\Omega}}} % U+2126 absent de la police
\DeclareSIUnit\division{div} % réglages d'oscilloscope (ms/div, V/div)
\DeclareSIUnit\tr{tr} % tours (vitesse de rotation, ex. \SI{1500}{\tr\per\minute})
\DeclareSIUnit\molar{M} % concentration molaire (ex. \SI{3}{\molar}, \SI{1}{\micro\molar})
\DeclareSIUnit\an{an} % année (le modèle confond parfois avec \year, un compteur TeX) — ex. \SI{5}{\kilo\meter\per\an}
\DeclareSIUnit\FCFA{FCFA} % franc CFA (ex. \SI{500}{\FCFA\per\kilo\gram})
\DeclareSIUnit\degreeBrix{\textdegree Brix} % teneur en sucre d'un sirop/jus (réfractométrie)
\DeclareSIUnit\month{mois} % le modèle utilise parfois \month, un compteur TeX, comme unité de durée
\DeclareSIUnit\microliter{\micro\liter} % le modèle écrit parfois \microliter en un seul token
\DeclareSIUnit\million{millions} % dénombrement (ex. \SI{15}{\million\per\mL} spermatozoïdes)
\usepackage{qrcode}
% \degree : commande fréquemment utilisée par les modèles pour un angle ou une
% température en dehors de \SI{}{} (non fournie par les paquets chargés).
\providecommand{\degree}{^{\circ}}
% \qed (fin de démonstration), utilisé par les modèles sans amsthm
\providecommand{\qed}{\hfill$\square$}
% \barre : double barre des valeurs interdites dans un tableau de variations (tkz-tab)
\providecommand{\barre}{\ensuremath{\|}}
% environnement proof (amsthm non chargé) utilisé par les modèles
\ifdefined\proof\else\newenvironment{proof}[1][Démonstration]{\par\noindent\textit{#1.}\ }{\qed\par}\fi
\usepackage{lastpage}
\usepackage{hyperref}
\hypersetup{hidelinks}

% ============================================================
% Palette « Pop » OnBuch+ — mêmes hex que la classe P (plus_ui.dart)
% ============================================================
\definecolor{poporange}{HTML}{F59321}
\definecolor{popdark}{HTML}{E07A0C}
\definecolor{popcream}{HTML}{FFF6E8}
\definecolor{popink}{HTML}{1C1714}
\definecolor{poppurple}{HTML}{7A5AE0}
\definecolor{popgreen}{HTML}{2E9E6A}
\definecolor{poppink}{HTML}{E0457B}
\definecolor{popblue}{HTML}{2F7DE1}
\definecolor{popgold}{HTML}{C98A12}
\definecolor{popmuted}{HTML}{8A7F76}
\definecolor{popline}{HTML}{EADFD2}
\definecolor{popgray}{HTML}{8A7F76}
\colorlet{popgrey}{popgray}
% Alias tolérants (le modèle nomme parfois les couleurs différemment)
\colorlet{popwhite}{white}
\colorlet{popblack}{popink}
\colorlet{popred}{poppink}
\colorlet{popyellow}{popgold}
% Teintes claires (fonds de tableaux/encadrés) — le modèle nomme parfois
% les couleurs avec un suffixe L (ex. popblueL) sans les avoir définies.
\colorlet{poporangeL}{poporange!20}
\colorlet{popdarkL}{popdark!20}
\colorlet{poppurpleL}{poppurple!20}
\colorlet{popgreenL}{popgreen!20}
\colorlet{poppinkL}{poppink!20}
\colorlet{popblueL}{popblue!20}
\colorlet{popgoldL}{popgold!20}
\colorlet{popredL}{popred!20}
\colorlet{popgrayL}{popgray!20}
% Teintes claires pour fonds de tableaux / zones
\colorlet{popblueL}{popblue!10!white}
\colorlet{poporangeL}{poporange!14!white}
\colorlet{popgreenL}{popgreen!12!white}
\colorlet{poppurpleL}{poppurple!12!white}
\colorlet{poppinkL}{poppink!12!white}
\colorlet{popgoldL}{popgold!14!white}

\pagecolor{popcream}

% ============================================================
% Typographie
% ============================================================
\defaultfontfeatures{Ligatures=TeX}
\setmainfont{PlusJakartaSans-Medium.ttf}[Path=../../../fonts/,BoldFont=PlusJakartaSans-Bold.ttf]
% Maths en sans-serif (Fira Math) avec chiffres et lettres droites de Plus
% Jakarta, pour que les formules se fondent dans le texte.
\usepackage[math-style=ISO,bold-style=ISO]{unicode-math}
\setmathfont{FiraMath-Regular.otf}[Path=../../../fonts/]
\setmathfont{PlusJakartaSans-Medium.ttf}[Path=../../../fonts/,range={up/num,up/{Latin,latin},\mathup}]
\usepackage{icomma} % virgule décimale sans espace en mode math
% Caractères mathématiques Unicode absents des polices de texte (Plus Jakarta,
% Archivo Black) : rendus par leur équivalent mathématique, en texte comme en
% formule — sinon ils s'affichent en carré vide (ex. « Arithmétique dans ℤ »).
\usepackage{newunicodechar}
\newunicodechar{ℝ}{\ensuremath{\mathbb{R}}}
\newunicodechar{ℤ}{\ensuremath{\mathbb{Z}}}
\newunicodechar{ℂ}{\ensuremath{\mathbb{C}}}
\newunicodechar{ℕ}{\ensuremath{\mathbb{N}}}
\newunicodechar{ℚ}{\ensuremath{\mathbb{Q}}}
\newunicodechar{𝔻}{\ensuremath{\mathbb{D}}}
\newunicodechar{α}{\ensuremath{\alpha}}
\newunicodechar{β}{\ensuremath{\beta}}
\newunicodechar{γ}{\ensuremath{\gamma}}
\newunicodechar{δ}{\ensuremath{\delta}}
\newunicodechar{ε}{\ensuremath{\varepsilon}}
\newunicodechar{θ}{\ensuremath{\theta}}
\newunicodechar{λ}{\ensuremath{\lambda}}
\newunicodechar{μ}{\ensuremath{\mu}}
\newunicodechar{σ}{\ensuremath{\sigma}}
\newunicodechar{τ}{\ensuremath{\tau}}
\newunicodechar{φ}{\ensuremath{\varphi}}
\newunicodechar{ψ}{\ensuremath{\psi}}
\newunicodechar{ω}{\ensuremath{\omega}}
\newunicodechar{Σ}{\ensuremath{\Sigma}}
% majuscules produites par \MakeUppercase dans les titres (α-aminés -> Α)
\newunicodechar{Α}{\ensuremath{\alpha}}
\newunicodechar{Β}{\ensuremath{\beta}}
\newunicodechar{Γ}{\ensuremath{\Gamma}}
\newunicodechar{Δ}{\ensuremath{\Delta}}
\newunicodechar{Ε}{\ensuremath{\varepsilon}}
\newunicodechar{Θ}{\ensuremath{\Theta}}
\newunicodechar{Λ}{\ensuremath{\Lambda}}
\newunicodechar{Μ}{\ensuremath{\mu}}
\newunicodechar{Τ}{\ensuremath{\tau}}
\newunicodechar{Φ}{\ensuremath{\Phi}}
\newunicodechar{Ψ}{\ensuremath{\Psi}}
\newunicodechar{Ω}{\ensuremath{\Omega}}
\newunicodechar{↦}{\ensuremath{\mapsto}}
\newunicodechar{∈}{\ensuremath{\in}}
\newunicodechar{∉}{\ensuremath{\notin}}
\newunicodechar{∧}{\ensuremath{\wedge}}
\newunicodechar{⇔}{\ensuremath{\Leftrightarrow}}
\newunicodechar{⟺}{\ensuremath{\Longleftrightarrow}}
\newunicodechar{⇒}{\ensuremath{\Rightarrow}}
\newunicodechar{⟹}{\ensuremath{\Longrightarrow}}
\newunicodechar{∘}{\ensuremath{\circ}}
\newunicodechar{∪}{\ensuremath{\cup}}
\newunicodechar{∩}{\ensuremath{\cap}}
\newunicodechar{≡}{\ensuremath{\equiv}}
\newunicodechar{⊂}{\ensuremath{\subset}}
\newunicodechar{⊥}{\ensuremath{\perp}}
\newunicodechar{∃}{\ensuremath{\exists}}
\newunicodechar{∀}{\ensuremath{\forall}}
\newunicodechar{∛}{\ensuremath{\sqrt[3]{\,}}}
\newunicodechar{∜}{\ensuremath{\sqrt[4]{\,}}}
\newunicodechar{⃗}{\ensuremath{{}^{\to}}}
\newunicodechar{ⁿ}{\ifmmode^{n}\else\textsuperscript{\ensuremath{n}}\fi}
\newunicodechar{ˣ}{\ifmmode^{x}\else\textsuperscript{\ensuremath{x}}\fi}
\newunicodechar{ᵇ}{\ifmmode^{b}\else\textsuperscript{\ensuremath{b}}\fi}
\newunicodechar{ʳ}{\ifmmode^{r}\else\textsuperscript{\ensuremath{r}}\fi}
\newunicodechar{ᵘ}{\ifmmode^{u}\else\textsuperscript{\ensuremath{u}}\fi}
\newunicodechar{ʸ}{\ifmmode^{y}\else\textsuperscript{\ensuremath{y}}\fi}
\newunicodechar{ᵃ}{\ifmmode^{a}\else\textsuperscript{\ensuremath{a}}\fi}
\newunicodechar{ᵉ}{\ifmmode^{e}\else\textsuperscript{\ensuremath{e}}\fi}
\newunicodechar{ᵏ}{\ifmmode^{k}\else\textsuperscript{\ensuremath{k}}\fi}
\newunicodechar{ᵖ}{\ifmmode^{p}\else\textsuperscript{\ensuremath{p}}\fi}
\newunicodechar{ᴺ}{\ifmmode^{N}\else\textsuperscript{\ensuremath{N}}\fi}
\newunicodechar{ᶜ}{\ifmmode^{c}\else\textsuperscript{\ensuremath{c}}\fi}
\newunicodechar{ʰ}{\ifmmode^{h}\else\textsuperscript{\ensuremath{h}}\fi}
\newunicodechar{ᵠ}{\ifmmode^{q}\else\textsuperscript{\ensuremath{q}}\fi}
\newunicodechar{ₙ}{\ifmmode_{n}\else\textsubscript{\ensuremath{n}}\fi}
\newunicodechar{ₐ}{\ifmmode_{a}\else\textsubscript{\ensuremath{a}}\fi}
\newunicodechar{ᵢ}{\ifmmode_{i}\else\textsubscript{\ensuremath{i}}\fi}
\newunicodechar{ᵣ}{\ifmmode_{r}\else\textsubscript{\ensuremath{r}}\fi}
\newunicodechar{ₖ}{\ifmmode_{k}\else\textsubscript{\ensuremath{k}}\fi}
\newunicodechar{ᵦ}{\ifmmode_{\beta}\else\textsubscript{\ensuremath{\beta}}\fi}
\newunicodechar{ₓ}{\ifmmode_{x}\else\textsubscript{\ensuremath{x}}\fi}
\newfontfamily\popdisplay{ArchivoBlack-Regular.ttf}[Path=../../../fonts/]
\newfontfamily\poplogo{SpaceGrotesk-Bold.ttf}[Path=../../../fonts/]
\newfontfamily\popsemi{PlusJakartaSans-SemiBold.ttf}[Path=../../../fonts/]
\newfontfamily\popxbold{PlusJakartaSans-ExtraBold.ttf}[Path=../../../fonts/]
\newfontfamily\popxboldls{PlusJakartaSans-ExtraBold.ttf}[Path=../../../fonts/,LetterSpace=3.0]

\setlist[enumerate]{leftmargin=1.7em,itemsep=5pt,topsep=4pt}
\setlist[itemize]{leftmargin=1.7em,itemsep=4pt,topsep=4pt,label={\color{poporange}\textbullet}}
\setlength{\parindent}{0pt}
\setlength{\parskip}{5pt}
\renewcommand{\arraystretch}{1.25}

% pgfplots : style Pop par défaut
\pgfplotsset{
  every axis/.append style={axis line style={popink,line width=1pt},tick style={popink},
    grid style={popline,line width=0.5pt},label style={font=\small},tick label style={font=\footnotesize}},
  popcurve/.style={poporange,line width=1.8pt,no markers,smooth},
}
% Même style disponible dans un \draw TikZ simple (hors pgfplots)
\tikzset{popcurve/.style={poporange,line width=1.8pt}}
\tikzset{popgrid/.style={popline,very thin}}
\colorlet{popcurve}{poporange} % le nom du style est parfois utilisé comme couleur

% ============================================================
% Pill / squiggle
% ============================================================
\newcommand{\poppill}[3]{%
  \tikz[baseline=-0.55ex]\node[draw=popink,line width=1.2pt,rounded corners=2.6pt,%
    fill=#2,inner xsep=6.5pt,inner ysep=3pt]{%
    \popxboldls\fontsize{7}{7}\selectfont\color{#3}\MakeUppercase{#1}};%
}
\newcommand{\popsquiggle}[1]{%
  \tikz[baseline=-0.4ex]{
    \draw[#1,line width=2.6pt,line cap=round]
      plot [smooth] coordinates {(0,0.09) (0.34,0.01) (0.68,0.21) (1.02,0.01) (1.36,0.10)};
  }%
}
% Mot-clé surligné (marqueur orange sous le texte)
\newcommand{\cle}[1]{{\popxbold\color{popdark}#1}}

% ============================================================
% En-tête / pied
% ============================================================
\pagestyle{fancy}
\fancyhf{}
\renewcommand{\headrulewidth}{0pt}
\renewcommand{\footrulewidth}{0pt}
\lhead{\raisebox{-0.15ex}{\poplogo\fontsize{13}{13}\selectfont\color{popink}OnBuch\color{poporange}+}}
\rhead{\poppill{\DOCMATIERE\ \textbullet\ \DOCNIVEAU\ \textbullet\ Leçon \DOCLECON}{white}{popink}}
\lfoot{\popxbold\fontsize{7.6}{7.6}\selectfont\color{popmuted}\MakeUppercase{Cours signé OnBuch+}\ \textbullet\ {\popsemi\DOCTITRE}}
\rfoot{\tikz[baseline=-0.6ex]\node[rounded corners=8.5pt,fill=popink,minimum height=17pt,minimum width=17pt,inner xsep=5pt,inner sep=0]{\popxbold\fontsize{8}{8}\selectfont\color{popcream}\thepage\,/\,\pageref*{LastPage}};}

% ============================================================
% Titres
% ============================================================
\newcommand{\chaptitre}[2]{%
  \begin{tcolorbox}[enhanced,colback=poporange,colframe=popink,boxrule=2.6pt,arc=11pt,
      drop shadow=popink,left=15pt,right=15pt,top=12pt,bottom=11pt,before skip=3pt,after skip=18pt]
    \poppill{#2}{popink}{popcream}\par\vspace{9pt}
    {\raggedright\hyphenpenalty=10000\exhyphenpenalty=10000\popdisplay\fontsize{22}{24}\selectfont\color{popcream}\MakeUppercase{#1}\par}
  \end{tcolorbox}
}
% Section numérotée : pastille orange avec le numéro + titre Archivo
\newcounter{popsec}
\newcommand{\coursec}[1]{%
  \stepcounter{popsec}\setcounter{popsub}{0}%
  \par\vspace{16pt}\noindent
  \begin{minipage}{\linewidth}
  \tikz[baseline=-0.7ex]\node[draw=popink,line width=1.6pt,fill=poporange,rounded corners=4pt,minimum size=22pt,inner sep=0]{\popdisplay\fontsize{12}{12}\selectfont\color{popcream}\thepopsec};%
  \hspace{8pt}{\popdisplay\fontsize{15}{17}\selectfont\color{popink}\MakeUppercase{#1}}\par
  \vspace{4pt}\noindent{\color{poporange}\rule{36pt}{3.6pt}}
  \end{minipage}\par\nopagebreak\vspace{8pt}
}
\newcounter{popsub}
\newcommand{\courssub}[1]{%
  \stepcounter{popsub}%
  \par\vspace{9pt}\noindent{\popxbold\fontsize{11.5}{13}\selectfont\color{popink}{\color{poporange}\thepopsec.\thepopsub}\hspace{6pt}#1}\par\nopagebreak\vspace{4pt}
}

% ============================================================
% Boîtes pédagogiques — même recette : fond blanc, bordure épaisse colorée,
% ombre dure encre, badge plein. Titre optionnel [#1].
% ============================================================
\tcbset{popbase/.style={breakable,enhanced,colback=white,boxrule=2.3pt,arc=9pt,
  shadow={3pt}{-3pt}{0pt}{popink},left=14pt,right=14pt,top=11pt,bottom=11pt,before skip=11pt,after skip=15pt,
  fontupper=\popsemi\fontsize{10.6}{14.5}\selectfont\color{popink},
  fontlower=\popsemi\fontsize{10.6}{14.5}\selectfont\color{popink}}}
% Coche dessinée (le glyphe ✓ n'existe pas dans les polices du document)
\newcommand{\popcheck}[1]{\tikz[baseline=-0.5ex]\draw[#1,line width=1.6pt,line cap=round,line join=round](-0.13,0.0)--(-0.04,-0.09)--(0.13,0.1);}
\providecommand{\coche}{\popcheck{popgreen}}
\providecommand{\resultat}[1]{\par\smallskip\noindent\fcolorbox{popgreen}{popgreenL}{\parbox{\dimexpr\linewidth-2\fboxsep-2\fboxrule\relax}{\textbf{Résultat.} #1}}\par\smallskip}
\newcommand{\popboxhead}[3]{\poppill{#1}{#2}{white}\ifx\relax#3\relax\else\hspace{6pt}{\popxbold\fontsize{10.6}{12}\selectfont\color{popink}#3}\fi\par\vspace{7pt}}

\newtcolorbox{aretenir}[1][]{popbase,colframe=popgreen,before upper={\popboxhead{À retenir}{popgreen}{#1}}}
\newtcolorbox{exemplebox}[1][]{popbase,colframe=poporange,before upper={\popboxhead{Exemple}{poporange}{#1}}}
\newtcolorbox{definition}[1][]{popbase,colframe=popblue,before upper={\popboxhead{Définition}{popblue}{#1}}}
\newtcolorbox{propriete}[1][]{popbase,colframe=poppurple,before upper={\popboxhead{Propriété}{poppurple}{#1}}}
\newtcolorbox{methode}[1][]{popbase,colframe=popgold,before upper={\popboxhead{Méthode}{popgold}{#1}}}
\newtcolorbox{attention}[1][]{popbase,colframe=poppink,before upper={\popboxhead{Attention}{poppink}{#1}}}
\newtcolorbox{experience}[1][]{popbase,colframe=popblue,colback=popblueL,before upper={\popboxhead{Expérience}{popblue}{#1}}}
% « Le savais-tu ? » : carte violette pleine (contexte, culture, Cameroun)
\newtcolorbox{savaistu}[1][]{popbase,colback=poppurple,colframe=popink,
  fontupper=\popsemi\fontsize{10.4}{14.2}\selectfont\color{white},
  before upper={\poppill{Le savais-tu\,?}{white}{poppurple}\ifx\relax#1\relax\else\hspace{6pt}{\popxbold\color{white}#1}\fi\par\vspace{7pt}}}
% Exercice résolu : énoncé en haut, \tcblower puis la solution sur fond crème
\newcounter{popexo}\newcounter{popexr}
\newtcolorbox{exoresolu}[1][]{popbase,colframe=poporange,colbacklower=popcream,
  segmentation style={popink,line width=1.2pt,dashed},
  before upper={\stepcounter{popexr}\popboxhead{Exercice résolu \thepopexr}{poporange}{#1}},
  before lower={\poppill{Solution}{popink}{popcream}\par\vspace{6pt}}}
% Exercice (non résolu en place) — la correction est dans la partie Corrigés
\newtcolorbox{exercice}[1][]{popbase,colframe=popink,
  before upper={\stepcounter{popexo}\popboxhead{Exercice \thepopexo}{popink}{#1}}}
% Corrigé
\newtcolorbox{corrige}[1][]{popbase,colframe=popgreen,colback=popgreenL,
  before upper={\popboxhead{Corrigé}{popgreen}{#1}}}
% Objectifs / prérequis / bilan
\newtcolorbox{objectifs}{popbase,colframe=popink,colback=poporangeL,
  before upper={\popboxhead{Objectifs de la leçon}{popink}{}}}
\newtcolorbox{prerequis}{popbase,colframe=popink,colback=popblueL,
  before upper={\popboxhead{Prérequis}{popblue}{}}}
\newtcolorbox{bilan}{popbase,colframe=popink,colback=white,boxrule=2.8pt,
  before upper={\popboxhead{Fiche bilan}{poporange}{L'essentiel à connaître pour l'examen}}}
% Figure encadrée avec légende
\newtcolorbox{popfigure}[1][]{enhanced,breakable=false,colback=white,colframe=popline,boxrule=1.4pt,arc=8pt,
  left=10pt,right=10pt,top=10pt,bottom=8pt,before skip=10pt,after skip=12pt,halign=center,
  lower separated=false,halign lower=center,
  fontlower=\popsemi\fontsize{9}{11.5}\selectfont\color{popmuted},
  #1}
\newcommand{\legende}[1]{\par\vspace{6pt}{\popsemi\fontsize{9}{11.5}\selectfont\color{popmuted}#1}}

% ============================================================
% Signature OnBuch+
% ============================================================
\newcommand{\poplogoinline}[1]{{\poplogo\fontsize{#1}{#1}\selectfont\color{popink}OnBuch\color{poporange}+}}
\newcommand{\signatureonbuch}{%
  \begin{tcolorbox}[enhanced,colback=popink,colframe=popink,boxrule=2.2pt,arc=10pt,
      drop shadow=poporange,left=14pt,right=14pt,top=10pt,bottom=10pt,before skip=0pt,after skip=0pt]
    \tikz[baseline=-0.7ex]\node[circle,fill=popgreen,draw=popcream,line width=1.3pt,minimum size=22pt,inner sep=0]{\popcheck{white}};%
    \hspace{9pt}\begin{minipage}[c]{0.62\linewidth}
      {\popxbold\fontsize{12}{13}\selectfont\color{popcream}Signé\ \textbullet\ {\poplogo OnBuch\color{poporange}+}}\par\vspace{2pt}
      {\raggedright\popsemi\fontsize{8}{10.5}\selectfont\color{popline}\MakeUppercase{Équipe pédagogique OnBuch+}\\\MakeUppercase{Conforme au programme officiel MINESEC}\par}
    \end{minipage}\hfill
    {\poplogo\fontsize{22}{22}\selectfont\color{popcream}OnBuch\color{poporange}+}
  \end{tcolorbox}%
}

% ============================================================
% Page de garde
% #1 titre · #2 matière · #3 niveau · #4 module · #5 n° leçon · #6 durée ·
% #7 description / objectifs (texte ou itemize)
% ============================================================
\newcommand{\pagegarde}[7]{%
  \thispagestyle{empty}
  \newgeometry{a4paper,margin=1.7cm,top=1.6cm,bottom=1.4cm}%
  \begin{tikzpicture}[remember picture,overlay]
    \fill[poporange!18] ([xshift=-3.2cm,yshift=-3.6cm]current page.north east) circle (4.6cm);
    \fill[poppurple!14] ([xshift=2.4cm,yshift=7.6cm]current page.south west) circle (3.8cm);
  \end{tikzpicture}%
  {\poplogo\fontsize{30}{30}\selectfont\color{popink}OnBuch\color{poporange}+}\hfill
  \raisebox{0.6ex}{\poppill{Cours signé \textbullet\ Premium}{popink}{popcream}}\par
  \vspace{1pt}\popsquiggle{poppurple}\par
  \vspace{8pt}
  \begin{tcolorbox}[enhanced,colback=poporange,colframe=popink,boxrule=2.8pt,arc=13pt,
      drop shadow=popink,left=18pt,right=18pt,top=12pt,bottom=13pt,after skip=10pt]
    \poppill{#2 \textbullet\ #3}{popink}{popcream}\hspace{5pt}\poppill{#4}{white}{popink}\par\vspace{8pt}
    {\popdisplay\fontsize{14}{14}\selectfont\color{popink}LEÇON #5}\par\vspace{4pt}
    {\raggedright\hyphenpenalty=10000\exhyphenpenalty=10000\popdisplay\fontsize{25}{27}\selectfont\color{popcream}\MakeUppercase{#1}\par}
    \vspace{8pt}
    {\popxbold\fontsize{9.5}{9.5}\selectfont\color{popink}\MakeUppercase{Durée conseillée : #6}}
  \end{tcolorbox}
  \begin{tcolorbox}[enhanced,colback=white,colframe=popink,boxrule=2.2pt,arc=10pt,
      drop shadow=popink,left=16pt,right=16pt,top=10pt,bottom=10pt,after skip=8pt]
    \poppill{Dans cette leçon}{poppurple}{white}\par\vspace{5pt}
    \popsemi\fontsize{9.3}{12.6}\selectfont\color{popink}#7
  \end{tcolorbox}
  \noindent\begin{minipage}[t]{0.487\linewidth}
    \begin{tcolorbox}[enhanced,colback=popgreen,colframe=popink,boxrule=2.2pt,arc=10pt,
        drop shadow=popink,left=11pt,right=11pt,top=10pt,bottom=10pt,height=3.1cm,valign=center]
      \tcbox[colback=white,colframe=popink,boxrule=1.3pt,arc=5pt,boxsep=0pt,nobeforeafter,
        left=3pt,right=3pt,top=3pt,bottom=3pt,valign=center]{\qrcode[height=1.9cm]{https://play.google.com/store/apps/details?id=com.luvvix.onbuch&hl=fr}}%
      \hspace{8pt}\begin{minipage}[c]{0.5\linewidth}
        {\popdisplay\fontsize{11}{12.5}\selectfont\color{white}\MakeUppercase{Télécharge OnBuch}}\par\vspace{4pt}
        {\raggedright\popsemi\fontsize{8.4}{11}\selectfont\color{white}Gratuit \textbullet\ Google Play\\Tuteur IA, annales, concours, résultats\par}
      \end{minipage}
    \end{tcolorbox}
  \end{minipage}\hfill
  \begin{minipage}[t]{0.487\linewidth}
    \begin{tcolorbox}[enhanced,colback=poppurple,colframe=popink,boxrule=2.2pt,arc=10pt,
        drop shadow=popink,left=11pt,right=11pt,top=10pt,bottom=10pt,height=3.1cm,valign=center]
      \tikz[baseline=-0.7ex]\node[circle,fill=white,draw=popink,line width=1.3pt,minimum size=1.6cm,inner sep=0]{\popdisplay\fontsize{24}{24}\selectfont\color{poppurple}+};%
      \hspace{8pt}\begin{minipage}[c]{0.6\linewidth}
        {\popdisplay\fontsize{11}{12.5}\selectfont\color{white}\MakeUppercase{Passe à OnBuch+}}\par\vspace{4pt}
        {\raggedright\popsemi\fontsize{8.4}{11}\selectfont\color{white}Tous les cours signés, Tuteur IA illimité, fiches PDF — onglet OnBuch+ de l'app\par}
      \end{minipage}
    \end{tcolorbox}
  \end{minipage}
  \vspace{8pt}
  \popcontenu
  \vfill
  \signatureonbuch
  \restoregeometry
  \newpage
}

% Rangée « Ce cours contient » (page de garde)
\newcommand{\poptile}[3]{% #1 couleur #2 gros texte #3 légende
  \begin{tcolorbox}[enhanced,colback=#1,colframe=popink,boxrule=2pt,arc=9pt,shadow={2.5pt}{-2.5pt}{0pt}{popink},
      left=6pt,right=6pt,top=9pt,bottom=9pt,halign=center,valign=center,height=1.75cm,width=\linewidth,nobeforeafter]
    {\popdisplay\fontsize{12.5}{13}\selectfont\color{white}\MakeUppercase{#2}}\par\vspace{3pt}
    {\centering\popsemi\fontsize{7.8}{9.5}\selectfont\color{white}#3\par}
  \end{tcolorbox}}
\newcommand{\popcontenu}{%
  \par\noindent{\popxboldls\fontsize{7.5}{7.5}\selectfont\color{popmuted}CE COURS SIGNÉ CONTIENT}\par\vspace{7pt}
  \noindent\begin{minipage}[t]{0.235\linewidth}\poptile{popblue}{Cours}{complet, illustré, pas à pas}\end{minipage}\hfill
  \begin{minipage}[t]{0.235\linewidth}\poptile{poporange}{Méthodes}{et exercices résolus}\end{minipage}\hfill
  \begin{minipage}[t]{0.235\linewidth}\poptile{poppink}{Exercices}{type examen + corrigés}\end{minipage}\hfill
  \begin{minipage}[t]{0.235\linewidth}\poptile{popgreen}{Fiche bilan}{l'essentiel à retenir}\end{minipage}\par}

% Sommaire
\newtcolorbox{sommaire}{enhanced,colback=white,colframe=popink,boxrule=2.2pt,arc=10pt,
  drop shadow=popink,left=18pt,right=18pt,top=13pt,bottom=13pt,before skip=6pt,after skip=20pt,
  before upper={\poppill{Sommaire}{poppurple}{white}\par\vspace{9pt}\popsemi\fontsize{10.6}{14.5}\selectfont\color{popink}}}

% Signature de fin de cours — poussée tout en bas de la dernière page
\newcommand{\findecours}{%
  \par\vfill\nopagebreak
  \begin{center}\popsquiggle{poporange}\end{center}\vspace{4pt}
  \signatureonbuch
}
\AtBeginDocument{\def\llbracket{\mathopen{[\mkern-3.5mu[}}\def\rrbracket{\mathclose{]\mkern-3.5mu]}}} % intervalles d'entiers (glyphe ⟦ absent de la police)
% Glyphes absents de Fira Math : redéfinis à partir de glyphes disponibles
\AtBeginDocument{%
  \def\setminus{\mathbin{\backslash}}\let\smallsetminus\setminus
  \def\vdots{\mathord{\vbox{\baselineskip3.2pt\lineskiplimit0pt\kern4pt\hbox{.}\hbox{.}\hbox{.}}}}%
  \def\ddots{\mathinner{\mkern1mu\raise6.5pt\hbox{.}\mkern2mu\raise3.6pt\hbox{.}\mkern2mu\raise0.7pt\hbox{.}\mkern1mu}}%
  \def\triangle{\mathord{\scalebox{0.9}{$\symup{\Delta}$}}}%
  \def\nabla{\mathord{\raisebox{\depth}{\scalebox{1}[-1]{$\symup{\Delta}$}}}}%
  \def\checkmark{\popcheck{popdark}}%
  \def\star{\mathbin{\ast}}\def\bigstar{\mathord{\ast}}%
  \def\diamond{\mathbin{\scalebox{0.6}{$\Diamond$}}}%
  \def\bigcup{\mathop{\mathchoice{\raisebox{-0.3em}{\scalebox{1.7}{$\cup$}}}{\scalebox{1.25}{$\cup$}}{\cup}{\cup}}\displaylimits}%
  \def\bigcap{\mathop{\mathchoice{\raisebox{-0.3em}{\scalebox{1.7}{$\cap$}}}{\scalebox{1.25}{$\cap$}}{\cap}{\cap}}\displaylimits}%
  \def\lesseqqgtr{\mathrel{\lessgtr}}\def\bigoplus{\mathop{\oplus}\displaylimits}%
}
\providecommand{\ssi}{si et seulement si} % abréviation fréquente
\AtBeginDocument{\providecommand{\wideparen}{\overparen}} % arc de cercle

%% FIGFIX-BEGIN v3 — correctifs globaux des figures (docs/onbuch-generation/tools/figfix)
\usepackage{adjustbox}\usepackage{etoolbox}
% toute figure TikZ/pgfplots plus large que la ligne est réduite à la largeur disponible
\BeforeBeginEnvironment{tikzpicture}{\begin{adjustbox}{max width=\linewidth}}
\AfterEndEnvironment{tikzpicture}{\end{adjustbox}}
% légendes pgfplots lisibles par-dessus les courbes
\pgfplotsset{every axis/.append style={legend style={fill=white,fill opacity=0.92,text opacity=1,draw=black!15,inner sep=3pt}}}
% étiquettes d'arêtes (syntaxe edge["..."]) avec fond blanc
\tikzset{every edge quotes/.append style={fill=white,inner sep=1.2pt}}
% bulles de cartes mentales : texte à la ligne dans la bulle, bulle un peu plus grande
\tikzset{every concept/.append style={text width=2.0cm,align=center,minimum size=2.3cm,inner sep=2pt}}
%% FIGFIX-END
\begin{document}

\pagegarde{La Chine : de la révolution culturelle à l’émergence économique}{Histoire}{Tle Tech}{Histoire – classe de Terminale technique}{N02}{1 séance (fiche de progression harmonisée)}{Cette leçon retrace le parcours de la Chine depuis la Révolution culturelle (1966-1976) jusqu'à son émergence comme deuxième puissance économique mondiale. Elle montre comment les réformes de Deng Xiaoping ont transformé un pays communiste fermé en « atelier du monde », tout en interrogeant les limites de ce modèle et ses implications pour l'Afrique et le Cameroun.
\vspace{4pt}
\begin{multicols}{2}\small
\begin{itemize}
\item À la fin de cette leçon, je sais situer dans le temps et dans l'espace les grandes étapes de l'histoire de la Chine de 1966 à nos jours
\item À la fin de cette leçon, je sais définir la Révolution culturelle et en expliquer les causes et les conséquences
\item À la fin de cette leçon, je sais décrire les réformes engagées par Deng Xiaoping à partir de 1978
\item À la fin de cette leçon, je sais analyser des documents (cartes, photographies, données chiffrées) pour établir les facteurs de l'émergence économique chinoise
\item À la fin de cette leçon, je sais construire un croquis ou un diagramme simple à partir de données statistiques
\item À la fin de cette leçon, je sais identifier les limites et les contradictions de l'émergence chinoise
\end{itemize}\end{multicols}}

\begin{sommaire}
\begin{multicols}{2}
\begin{enumerate}
\item La Chine de Mao et la Révolution culturelle (1966-1976)
\item Le tournant de Deng Xiaoping : les réformes (1978-1989)
\item L'émergence économique : la Chine, « atelier du monde » (années 1990-2010)
\item Les limites et les contradictions de l'émergence chinoise
\item La Chine, l'Afrique et le Cameroun : une présence croissante
\item Activité d'intégration
\item Méthodes et exercices résolus
\item Exercices
\item Corrigés des exercices
\item Fiche bilan
\end{enumerate}
\end{multicols}
\end{sommaire}

\begin{objectifs}
\begin{itemize}
\item À la fin de cette leçon, je sais situer dans le temps et dans l'espace les grandes étapes de l'histoire de la Chine de 1966 à nos jours
\item À la fin de cette leçon, je sais définir la Révolution culturelle et en expliquer les causes et les conséquences
\item À la fin de cette leçon, je sais décrire les réformes engagées par Deng Xiaoping à partir de 1978
\item À la fin de cette leçon, je sais analyser des documents (cartes, photographies, données chiffrées) pour établir les facteurs de l'émergence économique chinoise
\item À la fin de cette leçon, je sais construire un croquis ou un diagramme simple à partir de données statistiques
\item À la fin de cette leçon, je sais identifier les limites et les contradictions de l'émergence chinoise
\item À la fin de cette leçon, je sais établir des liens entre l'émergence de la Chine et des situations concrètes de la vie courante au Cameroun
\item À la fin de cette leçon, je sais rédiger et justifier une conclusion argumentée sur la place de la Chine dans le monde actuel
\end{itemize}
\end{objectifs}

\begin{prerequis}
\begin{itemize}
\item La Seconde Guerre mondiale et la bipolarisation du monde (classe de Première)
\item La décolonisation et l'émergence du tiers-monde (classe de Première)
\item Notions de base : communisme, capitalisme, mondialisation
\item La lecture de cartes et de graphiques statistiques (géographie, classes précédentes)
\item La guerre froide et les relations Est-Ouest
\end{itemize}
\end{prerequis}

\newpage

\coursec{La Chine de Mao et la Révolution culturelle (1966-1976)}

À Yaoundé comme à Douala, les marchés regorgent de produits « made in China » : téléphones, motos, tissus, matériels de construction. Les grands chantiers camerounais (stades, routes, barrages comme Memve'ele) sont souvent financés ou construits par des entreprises chinoises. Pourtant, il y a soixante ans, la Chine était un pays pauvre, déchiré par la Révolution culturelle. Comment ce pays est-il devenu en quelques décennies l'« atelier du monde » et un partenaire incontournable de l'Afrique et du Cameroun ? C'est cette métamorphose que nous allons expliquer, en commençant par comprendre le chaos dont elle est sortie.

\courssub{De la proclamation de la RPC à l'échec du Grand Bond en avant (1949-1961)}

Le \textbf{1er octobre 1949}, sur la place Tian'anmen à Pékin, \textbf{Mao Zedong} proclame la fondation de la \textbf{République populaire de Chine (RPC)}. Le Parti communiste chinois (PCC) l'emporte après une longue guerre civile contre les nationalistes du Kuomintang (KMT) réfugiés à Taïwan. Le nouveau régime entreprend la reconstruction d'un pays exsangue : réforme agraire (répartition des terres aux paysans), nationalisation de l'industrie, alliance avec l'URSS (traité de 1950). Mais dès la fin des années 1950, Mao juge le modèle soviétique trop bureaucratique et trop lent. Il veut une industrialisation rapide, fondée sur la mobilisation de la main-d'œuvre paysanne.

\begin{popfigure}
\begin{tikzpicture}[scale=0.9, every node/.style={font=\small}]
% Axe principal
\draw[thick, popdark] (0,0) -- (16,0);
% Graduations et événements
\foreach \x/\annee/\event/\detail in {
    0/1949/Proclamation RPC/{Mao Zedong proclame la République populaire de Chine},
    3.2/1958/Grand Bond en avant/{Industrialisation forcée, communes populaires, famine},
    7.2/1966/Révolution culturelle/{Déclenchement par Mao, Gardes rouges, chaos},
    11.2/1971/Mort de Lin Biao/{Disparition du dauphin désigné de Mao},
    16/1976/Fin/{Mort de Mao, arrestation de la Bande des Quatre}} {
    \draw[thick, popblue] (\x,0.2) -- (\x,-0.2);
    \node[above=0.1cm, align=center, popblue] at (\x,0.5) {\textbf{\annee}};
    \node[above=0.8cm, align=center, text width=3cm, popdark] at (\x,1.3) {\event};
    \node[below=0.1cm, align=center, text width=3cm, popmuted] at (\x,-0.8) {\detail};
}
% Flèches de liaison
\draw[poporange, thick, ->] (0.5,-1.5) -- (2.7,-1.5) node[midway, above, poporange] {Construction socialiste};
\draw[poporange, thick, ->] (3.7,-1.5) -- (6.7,-1.5) node[midway, above, poporange] {Échec, retrait de Mao};
\draw[poporange, thick, ->] (7.7,-1.5) -- (10.7,-1.5) node[midway, above, poporange] {Chaos, purges};
\draw[poporange, thick, ->] (11.7,-1.5) -- (15.5,-1.5) node[midway, above, poporange] {Fin de l'ère Mao};
\end{tikzpicture}
\legende{Frise chronologique : les grandes étapes de la Chine maoïste de 1949 à 1976.}
\end{popfigure}

En \textbf{1958}, Mao lance le \textbf{Grand Bond en avant} : création de \emph{communes populaires} (fusion de coopératives), \emph{backyard furnaces} (petits fours à acier dans les cours), objectifs de production irréalistes. C'est un \textbf{désastre absolu}. La désorganisation agricole, les fausses statistiques et les calamités naturelles provoquent une \textbf{famine massive (1959-1961)} : entre \textbf{15 et 45 millions de morts} selon les estimations. Mao perd la face et doit s'effacer temporairement de la direction quotidienne de l'État au profit de \textbf{Liu Shaoqi} (président de la République) et \textbf{Deng Xiaoping} (secrétaire général du PCC), qui lancent une politique de redressement économique pragmatique (responsabilités individuelles, marchés ruraux). Mao, marginalisé, craint que la Chine ne prenne la « voie capitaliste » comme l'URSS selon lui. Il prépare sa revanche.

\courssub{Le déclenchement de la Révolution culturelle (1966) : objectifs et acteurs}

En \textbf{mai 1966}, Mao lance la \textbf{Grande Révolution culturelle prolétarienne}. Officiellement, il s'agit de « continuer la révolution sous la dictature du prolétariat ». En réalité, c'est une \textbf{lutte de pouvoir} pour éliminer ses rivaux (Liu, Deng) qualifiés de « \emph{révisionnistes} » et de « \emph{autorités prenant la voie capitaliste} », et pour imposer sa vision d'une révolution permanente, égalitariste et anti-bureaucratique.

\begin{definition}[La Révolution culturelle (1966-1976)]
Mouvement politique et social lancé par Mao Zedong en 1966 pour éliminer ses opposants au sein du Parti communiste chinois (les « révisionnistes ») et « purifier » la révolution en mobilisant la jeunesse contre les « Quatre Vieilleries » (vieilles idées, vieille culture, vieux usages, vieilles habitudes). Elle plonge la Chine dans le chaos pendant dix ans.
\end{definition}

L'instrument principal de Mao est la \textbf{jeunesse étudiante et lycéenne}. Dès août 1966, des millions de jeunes forment des \textbf{Gardes rouges} (\emph{Hongweibing}). Ils brandissent le \textbf{Petit Livre rouge}, recueil de citations de Mao, érigé en vérité absolue et en passeport pour l'action révolutionnaire. Mao les reçoit en masse sur la place Tian'anmen, leur donnant une légitimité directe.

\begin{definition}[Les Gardes rouges]
Groupes de jeunes militants (étudiants, lycéens), fanatisés par le culte de la personnalité de Mao Zedong, chargés d'attaquer les « ennemis de la révolution » (intellectuels, cadres, anciens, religieux) et de détruire les symboles de l'ancien monde. Ils agissent souvent avec une violence extrême, hors de tout contrôle légal.
\end{definition}

\begin{savaistu}[Le Petit Livre rouge]
Le \emph{Recueil de citations du président Mao} (Petit Livre rouge) a été imprimé à \textbf{plus d'un milliard d'exemplaires} entre 1964 et 1976, dans des dizaines de langues. C'est l'un des livres les plus diffusés de l'histoire, surpassé seulement par la Bible. Il servait de « boussole » idéologique : on le brandissait lors des réunions, on en récitait des passages pour justifier n'importe quelle action, y compris les pires violences.
\end{savaistu}

\begin{popfigure}
\begin{tikzpicture}[scale=0.8, every node/.style={font=\small}]
% Représentation schématique d'une photo de Gardes rouges
% Fond : foule stylisée
\foreach \i in {1,...,15} {
    \pgfmathsetmacro{\x}{rand*3}
    \pgfmathsetmacro{\y}{-rand*1.5}
    \pgfmathsetmacro{\s}{0.8+rnd*0.4}
    \begin{scope}[shift={(\x,\y)}, scale=\s]
        \draw[fill=popred!80!popdark] (0,0) ellipse ({0.3} and {0.5}); % Corps (uniforme)
        \draw[fill=poppink] (0,0.55) circle (0.25); % Tête
        % Bras levé avec livre
        \draw[popred!80!popdark, line width=2pt] (0.2,0.2) -- (0.6,1.0);
        \draw[fill=popred!90!black] (0.55,1.0) rectangle (0.85,1.3); % Livre rouge
    \end{scope}
}
% Premier plan : un Garde rouge détaillé
\begin{scope}[shift={(0, -2.5)}]
    \draw[fill=popred!70!black] (-0.4,-1) rectangle (0.4,0.5); % Torse
    \draw[fill=poppink] (0,0.8) circle (0.35); % Tête
    % Casquette
    \draw[fill=popgreen!60!black] (-0.4,1.05) -- (0.4,1.05) -- (0.5,1.3) -- (-0.5,1.3) -- cycle;
    \draw[fill=popyellow] (0,1.25) circle (0.1); % Étoile
    % Bras droit levé
    \draw[popred!70!black, line width=4pt, line cap=round] (0.3,0) -- (1.0,1.5);
    % Livre rouge dans la main
    \draw[fill=popred!90!black, popdark] (0.85,1.4) rectangle (1.35,1.9);
    \node[popwhite, font=\tiny\bfseries, rotate=90] at (1.1,1.65) {MAO};
    % Bras gauche poing levé
    \draw[poppink, line width=3pt, line cap=round] (-0.3,0) -- (-1.0,1.2);
    \draw[fill=poppink] (-1.1,1.2) circle (0.15);
\end{scope}
% Légendes fléchées
\draw[popblue, thick, <->] (-3.5,2) -- (-3.5,0.5) node[midway, left, align=right, text width=2.5cm] {Uniforme kaki / vert\\ type Armée de libération};
\draw[popblue, thick, <->] (3.5,1.8) -- (3.5,0.2) node[midway, right, align=left, text width=2.5cm] {Casquette militaire\\ avec étoile rouge};
\draw[popblue, thick, <->] (2.5,2.5) -- (1.5,1.8) node[midway, right, align=left, text width=2.5cm] {Petit Livre rouge\\ brandi comme une arme};
\draw[popblue, thick, <->] (-2.5,2.5) -- (-1.5,1.5) node[midway, left, align=right, text width=2.5cm] {Poing levé :\\ serment de fidélité};
\node[align=center, popdark, font=\bfseries] at (0,-4) {Gardes rouges brandissant le Petit Livre rouge (1966-1968)};
\end{tikzpicture}
\legende{Schéma de légende d'une photographie type de Gardes rouges : jeunesse militarisée, culte du chef, livre unique comme référence absolue.}
\end{popfigure}

\courssub{Les manifestations : destruction, humiliation, déportation (1966-1969)}

Le mot d'ordre est : « \emph{Il est juste de se révolter} ». Les Gardes rouges s'en prennent aux \textbf{Quatre Vieilleries} :
\begin{itemize}
    \item \textbf{Idées} : confucianisme, pensée bourgeoise, religion.
    \item \textbf{Culture} : livres brûlés, tableaux lacérés, temples rasés, opéras interdits (remplacés par les « opéras modèles » révolutionnaires de Jiang Qing, l'épouse de Mao).
    \item \textbf{Usages} : noms de rues changés, salutations traditionnelles interdites, vêtements « bourgeois » déchirés.
    \item \textbf{Habitudes} : famille nucléaire critiquée, fidélité au Parti avant la famille.
\end{itemize}

Les \textbf{humiliations publiques} (\emph{struggle sessions}) deviennent quotidiennes : professeurs, médecins, écrivains, cadres du Parti, anciens notables sont contraints de porter des pancartes infamantes, de s'agenouiller, de subir des violences physiques devant des foules hystériques. \textbf{Liu Shaoqi} et \textbf{Deng Xiaoping} sont arrêtés, humiliés, envoyés en détention (Liu meurt en détention en 1969).

À partir de \textbf{1968}, pour canaliser le chaos (les Gardes rouges se battent entre factions rivales), Mao lance le mouvement \textbf{« Haut en bas, bas en haut »} : \textbf{17 millions de jeunes urbains} (dont Xi Jinping, futur président) sont envoyés de force à la campagne pour être « rééduqués par les pauvres et moyens paysans ». C'est la génération « perdue » (\emph{Zhiqing}), privée d'études supérieures pendant une décennie.

\begin{exemplebox}[Chiffres clés de la Révolution culturelle]
\begin{itemize}
    \item \textbf{Morts directs} : estimations entre \textbf{1 et 2 millions} (violences, suicides, conditions de détention).
    \item \textbf{Jeunes envoyés à la campagne} : environ \textbf{17 millions} (1968-1978).
    \item \textbf{Patrimoine détruit} : \textbf{4 922} sites historiques classés détruits à Pékin seul ; bibliothèques brûlées, artefacts brisés.
    \item \textbf{Éducation} : universités fermées (1966-1970), examen d'entrée (Gaokao) supprimé jusqu'en 1977. Une génération entière d'intellectuels manquants.
\end{itemize}
\end{exemplebox}

\begin{exoresolu}[Analyse d'un témoignage : la « génération perdue »]
\textbf{Énoncé :} Extrait d'un témoignage de la « littérature des cicatrices » (fin des années 1970) : « Nous avions 15 ans en 1968. On nous a dit : "La campagne est un vaste ciel et une vaste terre où l'on peut s'épanouir." En réalité, c'était l'exil. Nous avons porté des charges de 50 kg sur des sentiers de montagne, mangé de l'écorce d'arbre, perdu nos illusions. Dix ans plus tard, nous revenions dans des villes que nous ne reconnaissions plus, sans diplôme, sans métier, avec une colère rentrée. »

\textbf{Questions :}
\begin{enumerate}
    \item Identifiez la politique d'État décrite dans ce témoignage.
    \item Expliquez les conséquences à long terme pour la société chinoise.
\end{enumerate}
\tcblower
\textbf{Solution détaillée :}
\begin{enumerate}
    \item Il s'agit de la politique des \textbf{jeunes intellectuels envoyés à la campagne} (\emph{Shangshan xiaxiang}), décidée par Mao en décembre 1968 pour désamorcer le chaos des Gardes rouges en ville. C'est une déportation massive de la jeunesse urbaine (lycéens, étudiants) vers les zones rurales reculées.
    \item \textbf{Conséquences :} (1) \emph{Perte de capital humain} : interruption des études pour une génération entière (creux démographique dans les cadres supérieurs actuels). (2) \emph{Traumatisme psychologique} : sentiment de trahison par l'État et le Parti, base du mécontentement qui mènera aux manifestations de 1989. (3) \emph{Rupture de transmission culturelle} : coupure entre traditions familiales et idéologie maoïste. (4) \emph{Ressentiment social} : ces jeunes, de retour en ville fin des années 1970, formeront le noyau dur des demandes de réforme et de vérité historique.
\end{enumerate}
\end{exoresolu}

\courssub{Bilans et fin de la Révolution culturelle (1969-1976)}

Après 1969, l'armée (PLA) sous le maréchal \textbf{Lin Biao} (désigné successeur de Mao) rétablit un ordre militaire. Mais Lin Biao tente un coup d'État avorté et meurt dans un crash d'avion en \textbf{1971} (fuite vers l'URSS). Mao, vieillissant, oscille entre les « radicaux » (la \textbf{Bande des Quatre}, menée par Jiang Qing) et les « modérés » (revenu de Deng Xiaoping en 1973, puis Zhou Enlai). L'économie stagne, la production industrielle chute, l'isolement diplomatique est total (sauf rapprochement avec les USA en 1972, visite de Nixon).

\begin{aretenir}[Bilan global de la Révolution culturelle (1966-1976)]
\begin{itemize}
    \item \textbf{Humain} : 1 à 2 millions de morts, dizaines de millions de vies brisées (persécutions, exil, suicides).
    \item \textbf{Économique} : Gaspillage de 10 ans de développement ; retard technologique accru sur l'Occident et le Japon ; agriculture stagnante.
    \item \textbf{Culturel} : Destruction massive du patrimoine ; rupture de la transmission du savoir ; système éducatif détruit.
    \item \textbf{Politique} : Affaiblissement de la légitimité du PCC ; crise de succession non résolue ; armée politisée.
\end{itemize}
La Révolution culturelle ne s'achève \textbf{officiellement qu'en octobre 1976} : \textbf{Mao meurt le 9 septembre}. Un mois plus tard, le nouveau président \textbf{Hua Guofei} fait arrêter la \textbf{Bande des Quatre} (Jiang Qing, Zhang Chunqiao, Yao Wenyuan, Wang Hongwen). Le pays est exsangue, la population aspire au retour à la normale. C'est dans ce vide que \textbf{Deng Xiaoping}, réhabilité une seconde fois, va imposer sa ligne : la \emph{Réforme et Ouverture}.
\end{aretenir}

\begin{attention}[Erreur fréquente : confusion des dates]
Ne confondez \textbf{JAMAIS} :
\begin{itemize}
    \item La \textbf{Révolution communiste (1949)} : prise de pouvoir par le PCC, fin de la guerre civile, proclamation de la RPC.
    \item La \textbf{Révolution culturelle (1966-1976)} : mouvement interne au pouvoir lancé par Mao 17 ans plus tard pour purger le Parti, qui a détruit la société chinoise.
\end{itemize}
Au Probatoire / Bac, une confusion sur ces dates ou ces natures (prise de pouvoir vs purge interne) entraîne une perte de points majeure.
\end{attention}

\courssub{Exercice d'application : Chronologie commentée}

\begin{exercice}[Mise en ordre et analyse]
Reclassez les événements suivants dans l'ordre chronologique et, pour chacun, précisez s'il marque un \textbf{renforcement} ou un \textbf{affaiblissement} du pouvoir personnel de Mao :
\begin{enumerate}
    \item Arrestation de la Bande des Quatre.
    \item Proclamation de la RPC.
    \item Lancement du Grand Bond en avant.
    \item Mort de Lin Biao.
    \item Déclenchement de la Révolution culturelle.
    \item Mort de Mao.
\end{enumerate}
\end{exercice}

\begin{corrige}[Exercice n°1]
\begin{enumerate}
    \item \textbf{1949 : Proclamation de la RPC} $\rightarrow$ \textbf{Renforcement} absolu (fondateur).
    \item \textbf{1958 : Grand Bond en avant} $\rightarrow$ \textbf{Affaiblissement} (échec économique, recul au 2e rang).
    \item \textbf{1966 : Révolution culturelle} $\rightarrow$ \textbf{Renforcement tactique} (élimination rivaux, mobilisation jeunesse), mais chaos ingérable.
    \item \textbf{1971 : Mort de Lin Biao} $\rightarrow$ \textbf{Affaiblissement} (échec du successeur désigné, mystère, perte de confiance).
    \item \textbf{1976 : Mort de Mao} $\rightarrow$ Fin de l'ère, pouvoir personnel éteint.
    \item \textbf{Octobre 1976 : Arrestation Bande des Quatre} $\rightarrow$ Liquidation de l'héritage radical de Mao par ses successeurs.
\end{enumerate}
\end{corrige}

\coursec{Le tournant de Deng Xiaoping : les réformes (1978-1989)}

À la mort de Mao Zedong en 1976, la Chine sort épuisée de dix années de Révolution culturelle : l'économie est désorganisée, les campagnes pauvres, les universités fermées ou affaiblies, et le pays est isolé sur la scène internationale. Comment relever un pays d'environ un milliard d'habitants sans renier le régime communiste ? C'est le défi que va relever Deng Xiaoping à partir de 1978, en lançant des réformes économiques spectaculaires tout en maintenant le monopole politique du Parti communiste chinois. Cette section étudie ce double mouvement : l'ouverture économique d'une part, la continuité autoritaire d'autre part.

\courssub{L'arrivée au pouvoir de Deng Xiaoping et les Quatre Modernisations}

Après la mort de Mao (9 septembre 1976) et l'arrestation de la « Bande des Quatre » (octobre 1976), une brève lutte de pouvoir oppose Hua Guofeng, successeur désigné par Mao, aux réformateurs. C'est \cle{Deng Xiaoping} qui s'impose : au troisième plénum du XI\textsuperscript{e} comité central du PCC, en \cle{décembre 1978}, il devient l'homme fort de la Chine, sans jamais cumuler officiellement tous les titres suprêmes. Sa légitimité repose sur une idée simple : le socialisme ne doit pas signifier la pauvreté, il doit apporter la prospérité au peuple.

\begin{definition}[Les Quatre Modernisations]
Les \cle{Quatre Modernisations} désignent le programme de réformes lancé par Deng Xiaoping en 1978 pour transformer la Chine en grande puissance d'ici à l'an 2000. Elles portent sur quatre secteurs : l'\textbf{agriculture}, l'\textbf{industrie}, la \textbf{défense} et les \textbf{sciences et technologies}. Ce programme, d'abord énoncé par Zhou Enlai en 1975, devient sous Deng la feuille de route officielle du régime.
\end{definition}

L'intuition de Deng est pragmatique : plutôt que de débattre d'idéologie, il faut juger les politiques à leurs résultats concrets. La priorité est donnée à l'agriculture, car c'est là que vivent environ 80 \% des Chinois et que la misère est la plus profonde.

\begin{savaistu}[« Le chat noir ou le chat blanc »]
Pour justifier l'ouverture au marché, Deng Xiaoping reprenait un proverbe du Sichuan : « \emph{Peu importe que le chat soit noir ou blanc, pourvu qu'il attrape les souris.} » Autrement dit : peu importe qu'une méthode soit « capitaliste » ou « socialiste », pourvu qu'elle fasse produire davantage. Il déclarait aussi : « \emph{Enrichissez-vous !} » et « \emph{Laisser certains s'enrichir d'abord} », rompant avec l'égalitarisme maoïste.
\end{savaistu}

\courssub{La réforme agricole : la décollectivisation progressive}

\textbf{Du système de responsabilité familiale à la hausse des rendements}\par

Sous Mao, les terres étaient organisées en \cle{communes populaires} : la terre, les outils et les récoltes appartenaient au collectif, et le paysan recevait des points de travail sans lien direct avec ses efforts. Résultat : faible motivation, rendements médiocres, famines récurrentes.

La réforme naît d'une expérience spontanée : en 1978, dans le village de \cle{Xiaogang} (province de l'Anhui), dix-huit paysans signent secrètement un pacte pour diviser les terres collectives entre les familles, chacune s'engageant à livrer son quota à l'État et à garder le surplus. Leurs récoltes explosent dès la première année. Au lieu de les punir, Deng Xiaoping généralise l'expérience : c'est le \cle{système de responsabilité familiale}, étendu à toute la Chine entre 1980 et 1983.

\begin{methode}[Comprendre le mécanisme de la réforme agricole]
\begin{enumerate}
\item \textbf{Situation initiale} : la terre appartient au collectif ; le paysan n'a aucun intérêt personnel à produire plus.
\item \textbf{Changement} : la terre reste propriété collective, mais elle est \emph{louée} à chaque famille pour 15 ans (puis 30 ans) ; la famille doit livrer un quota fixe à l'État à prix administré.
\item \textbf{Incitation} : tout ce qui est produit au-delà du quota peut être vendu librement sur le marché au profit de la famille.
\item \textbf{Résultat} : l'effort individuel est récompensé, donc la production augmente fortement.
\end{enumerate}
\end{methode}

Les résultats sont spectaculaires : la production céréalière passe d'environ 300 millions de tonnes en 1978 à plus de 400 millions de tonnes en 1984, soit une hausse d'environ un tiers en six ans. Les revenus ruraux augmentent, la famine recule, et des \cle{entreprises rurales} (les « TVE », entreprises de bourg et de village) se créent pour employer les paysans excédentaires et produire des biens de consommation. L'excédent de main-d'œuvre rurale alimentera ensuite les usines des zones côtières.

\courssub{L'ouverture économique : les Zones économiques spéciales}

\textbf{Les quatre ZES de 1980}\par

Pour attirer les capitaux, la technologie et les méthodes de gestion étrangères sans ouvrir tout le pays d'un coup, Deng Xiaoping choisit une stratégie prudente : expérimenter l'économie de marché dans des espaces limités, contrôlables, situés en périphérie. En \cle{1980}, quatre \cle{Zones économiques spéciales} (ZES) sont créées : \textbf{Shenzhen}, \textbf{Zhuhai}, \textbf{Shantou} et \textbf{Xiamen}.

\begin{definition}[Zone économique spéciale (ZES)]
Une \cle{Zone économique spéciale} est un territoire délimité qui bénéficie d'avantages fiscaux (impôts réduits) et douaniers (importations de machines et de matières premières en franchise) ainsi que d'une réglementation du travail souple, afin d'attirer les capitaux étrangers, les technologies modernes et de développer les exportations.
\end{definition}

Le choix des localisations n'est pas un hasard : Shenzhen et Zhuhai jouxtent respectivement \textbf{Hong Kong} et \textbf{Macao}, Xiamen fait face à \textbf{Taiwan}, et Shantou est liée aux réseaux de la diaspora chinoise d'Asie du Sud-Est. Il s'agit de capter les investissements des Chinois d'outre-mer et des firmes de Hong Kong, proches culturellement et géographiquement.

En \cle{1984}, l'expérience jugée concluante, l'ouverture est étendue à \textbf{14 villes côtières} (de Dalian au nord à Beihai au sud, en passant par Tianjin, Shanghai et Canton), puis à des régions entières : deltas du Yangzi et de la Rivière des Perles, péninsule du Shandong. En 1988, l'île de Hainan devient à son tour une ZES géante.

\begin{methode}[Analyser une carte économique]
\begin{enumerate}
\item \textbf{Lire la légende} : identifier les symboles (ici : ZES, villes ouvertes, ports, frontières).
\item \textbf{Localiser les zones dynamiques} : observer la répartition spatiale — ici, toutes les zones ouvertes se concentrent sur le \emph{littoral}, du nord-est au sud-est de la Chine.
\item \textbf{Formuler une hypothèse explicative} : le littoral offre des ports pour exporter, la proximité de Hong Kong, Macao et Taiwan, et des liaisons maritimes avec les marchés mondiaux ; l'intérieur du pays reste à l'écart de la croissance.
\end{enumerate}
\end{methode}

\begin{popfigure}
\begin{center}
\begin{tikzpicture}[scale=0.85, font=\small]
% Contour très simplifié de la Chine
\draw[thick, popdark, fill=popcream] (0,4.5) -- (1.5,5.2) -- (3,5.0) -- (4.5,5.5) -- (6,5.2) -- (7.2,4.6) -- (7.6,3.6) -- (7.0,2.6) -- (6.2,1.8) -- (5.2,1.2) -- (4.2,0.6) -- (3.0,0.9) -- (2.0,1.6) -- (1.0,2.2) -- (0.2,3.2) -- cycle;
% Mer de Chine
\node[popblue, font=\itshape] at (8.6,1.5) {Mer de Chine};
\draw[->, popblue] (8.2,1.9) -- (7.7,2.5);
% Nord
\draw[->, thick, popdark] (0.4,5.6) -- (0.4,6.3); \node[above, popdark] at (0.4,6.3) {N};
% Villes côtières ouvertes (14) - points bleus
\foreach \x/\y in {6.9/4.7, 7.3/4.3, 7.5/3.7, 7.6/3.1, 7.5/2.6, 7.3/2.2, 7.1/1.8, 6.6/1.3, 5.6/0.9}
  \fill[popblue] (\x,\y) circle (1.6pt);
% Etiquettes de quelques villes ouvertes
\node[above, popblue, font=\footnotesize] at (6.9,4.75) {Dalian};
\node[right, popblue, font=\footnotesize] at (7.35,4.3) {Tianjin};
\node[right, popblue, font=\footnotesize] at (7.65,3.1) {Shanghai};
\node[right, popblue, font=\footnotesize] at (6.65,1.35) {Canton};
% ZES - carrés orange
\foreach \x/\y in {6.9/1.6, 6.4/1.15, 6.15/1.0, 5.85/0.95}
  \fill[poporange] (\x,\y) rectangle ++(0.14,0.14);
% Etiquettes ZES
\node[right, poporange, font=\footnotesize] at (7.05,1.65) {Xiamen};
\node[right, poporange, font=\footnotesize] at (6.55,1.2) {Shantou};
\node[below left, poporange, font=\footnotesize] at (6.1,0.95) {Shenzhen};
\node[below, poporange, font=\footnotesize] at (5.85,0.85) {Zhuhai};
% Hong Kong et Macao
\fill[poppurple] (6.35,0.75) circle (1.4pt); \node[below, poppurple, font=\footnotesize] at (6.45,0.6) {Hong Kong};
\fill[poppurple] (5.7,0.7) circle (1.4pt); \node[below left, poppurple, font=\footnotesize] at (5.6,0.62) {Macao};
% Taiwan
\draw[thick, popmuted, fill=popmuted!20] (7.9,1.3) ellipse (0.18 and 0.45);
\node[popmuted, font=\footnotesize] at (8.35,1.75) {Taiwan};
% Pékin
\node[popgreen, font=\footnotesize] at (6.55,4.35) {$\bigstar$};
\node[above, popgreen, font=\footnotesize] at (6.55,4.45) {Pékin};
% Légende
\node[anchor=west, font=\footnotesize] at (0.3,1.9) {\textcolor{poporange}{\rule{0.18cm}{0.18cm}} ZES (1980)};
\node[anchor=west, font=\footnotesize] at (0.3,1.5) {\textcolor{popblue}{\textbullet} Villes côtières ouvertes (1984)};
\node[anchor=west, font=\footnotesize] at (0.3,1.1) {\textcolor{poppurple}{\textbullet} Hong Kong, Macao};
\end{tikzpicture}
\end{center}
\legende{Croquis schématique : localisation des quatre ZES de 1980 (Shenzhen, Zhuhai, Shantou, Xiamen) et de quelques villes côtières ouvertes en 1984. Toutes se concentrent sur le littoral, face à Hong Kong, Macao et Taiwan. Échelle indicative, contour simplifié.}
\end{popfigure}

\begin{exemplebox}[Shenzhen, la ville-champignon]
\cle{Shenzhen} n'était en 1980 qu'un ensemble de bourgs et de villages de pêcheurs d'environ \textbf{30 000 habitants}, choisi parce qu'il touchait la frontière de Hong Kong. Grâce aux investissements étrangers et à l'afflux de migrants ruraux, c'est devenu en quarante ans une métropole de plus de \textbf{12 millions d'habitants}, surnommée la « Silicon Valley chinoise », siège de géants technologiques comme Huawei et Tencent. C'est l'exemple le plus frappant de la rapidité de la transformation chinoise.
\end{exemplebox}

\courssub{Investissements étrangers et exportations : les premiers résultats}

Le mécanisme est le suivant : les firmes étrangères (d'abord de Hong Kong et Taiwan, puis du Japon, des États-Unis et d'Europe) installent des usines dans les ZES pour profiter d'une main-d'œuvre abondante, disciplinée et très bon marché. La Chine fournit le terrain et les ouvriers ; l'étranger fournit le capital, les machines et les débouchés. Les produits — textiles, jouets, électronique — sont exportés massivement, ce qui rapporte des devises. Le commerce extérieur chinois est multiplié par plus de quatre entre 1978 et 1989, et le PIB croît en moyenne d'environ 9 à 10 \% par an, un rythme presque sans précédent dans l'histoire.

\begin{popfigure}
\begin{center}
\begin{tikzpicture}
\begin{axis}[
  ybar, bar width=14pt,
  width=0.85\linewidth, height=6.2cm,
  ymin=0, ymax=1300,
  ylabel={PIB (milliards de dollars)},
  xlabel={Année},
  symbolic x coords={1978,1980,1985,1990,1995,2000},
  xtick=data,
  nodes near coords, every node near coord/.append style={font=\footnotesize},
  ymajorgrids, grid style={popline!40},
  axis line style={popdark},
]
\addplot[fill=popblue, draw=popdark] coordinates {(1978,150) (1980,190) (1985,310) (1990,360) (1995,735) (2000,1210)};
\end{axis}
\end{tikzpicture}
\end{center}
\legende{Évolution du PIB chinois de 1978 à 2000 (en milliards de dollars courants, ordres de grandeur Banque mondiale). Le décollage s'accélère nettement après 1990 : le PIB est multiplié par environ 8 entre 1978 et 2000.}
\end{popfigure}

\courssub{« Socialisme de marché » : ouverture économique, fermeture politique}

\begin{attention}[La Chine n'est pas devenue un pays capitaliste au sens politique]
Erreur fréquente dans les copies : affirmer que les réformes ont transformé la Chine en démocratie capitaliste. En réalité, le \cle{Parti communiste chinois} conserve le \textbf{monopole absolu du pouvoir} : aucun parti d'opposition n'est autorisé, la presse reste contrôlée, l'armée est au service du Parti. La formule officielle est celle d'« \cle{économie socialiste de marché} » : le marché est un \emph{outil} au service du socialisme, pas un régime politique. Deng Xiaoping pose dès 1979 les « Quatre Principes cardinaux » (voie socialiste, dictature démocratique populaire, direction du PCC, marxisme-léninisme-pensée de Mao) qui interdisent toute contestation du système.
\end{attention}

Cette contradiction entre libéralisation économique et verrouillage politique va provoquer des tensions. Les réformes créent aussi des effets pervers : inflation (les prix dépassent 18 \% en 1988), corruption des cadres du Parti qui profitent des doubles circuits de prix, inégalités croissantes entre littoral riche et intérieur rural. Les étudiants et les intellectuels, qui ont goûté à l'ouverture, réclament davantage de libertés.

\courssub{Les limites démocratiques : du printemps de Pékin à Tian'anmen}

Dès l'hiver 1978-1979, le « \cle{printemps de Pékin} » voit fleurir des dazibaos (affiches murales) sur le « Mur de la démocratie », où l'on réclame une « cinquième modernisation » : la démocratie. Deng fait rapidement fermer le mur et emprisonner le dissident Wei Jingsheng : le message est clair, la modernisation économique ne s'accompagnera pas d'un pluralisme politique.

Dix ans plus tard, la contestation reprend à plus grande échelle. À la mort du réformateur Hu Yaobang, en avril 1989, les étudiants occupent la place \cle{Tian'anmen} au cœur de Pékin. Pendant sept semaines, étudiants, ouvriers et habitants réclament la démocratie, la liberté de la presse et la lutte contre la corruption ; une statue de la « Déesse de la démocratie » est érigée face au portrait de Mao. Dans la nuit du \textbf{3 au 4 juin 1989}, Deng Xiaoping et les durs du Parti ordonnent l'intervention de l'armée : les chars et les soldats tirent sur la foule. Le bilan, jamais officiellement reconnu, est estimé entre plusieurs centaines et plus d'un millier de morts. La répression de Tian'anmen montre au monde entier que le PCC ne tolérera aucune remise en cause de son pouvoir, tout en poursuivant — après une brève pause — les réformes économiques.

\courssub{La politique de l'enfant unique (1979)}

Parallèlement aux réformes, le régime s'attaque à un problème jugé vital : la démographie. Avec près d'un milliard d'habitants en 1980, les dirigeants craignent que la croissance de la population n'absorbe tous les gains de la croissance économique. En \cle{1979}, la \cle{politique de l'enfant unique} est instaurée : chaque couple urbain ne doit avoir qu'un seul enfant (deux dans certaines campagnes si le premier est une fille). Le respect de la règle est encouragé par des avantages (logement, emploi, scolarité) et sanctionné par des amendes, voire des stérilisations et des avortements forcés.

Cette politique permet d'éviter, selon les estimations officielles, environ 300 à 400 millions de naissances, et contribue à l'augmentation du revenu par habitant. Mais elle a des conséquences lourdes : déséquilibre entre garçons et filles (environ 120 garçons pour 100 filles à la naissance dans les années 1990, du fait des préférences pour les garçons), vieillissement accéléré de la population, génération d'« enfants uniques » sans frères ni sœurs. Assouplie à partir de 2013, elle est abandonnée en 2015 (deux enfants autorisés) puis en 2021 (trois enfants).

\begin{exoresolu}[Analyser la stratégie de Deng Xiaoping]
\textbf{Énoncé.} Document : « Peu importe que le chat soit noir ou blanc, pourvu qu'il attrape les souris » (Deng Xiaoping, 1962, repris après 1978) ; « Laisser certains s'enrichir d'abord » (Deng Xiaoping). Questions : 1) Que signifie la première citation ? 2) En quoi la création des ZES illustre-t-elle cette philosophie ? 3) Quelle limite politique Deng impose-t-il aux réformes ?
\tcblower
\textbf{Solution détaillée.}
1) La première citation exprime le \textbf{pragmatisme} de Deng Xiaoping : une méthode économique ne doit pas être jugée selon son étiquette idéologique (« capitaliste » ou « socialiste ») mais selon son \textbf{efficacité concrète} à produire des richesses et à améliorer la vie du peuple. C'est une rupture avec le dogmatisme maoïste.
2) Les ZES illustrent ce pragmatisme : la Chine socialiste y autorise des entreprises étrangères, des salaires négociés, des profits privés — des pratiques « capitalistes » — parce qu'elles apportent capitaux, technologies et exportations. L'expérience est d'abord limitée à quatre zones (prudence), puis généralisée après succès (14 villes ouvertes en 1984).
3) La limite politique est le maintien du \textbf{monopole du PCC} : les « Quatre Principes cardinaux » interdisent toute contestation du socialisme et de la direction du Parti. La répression du printemps de Pékin (1979) puis de Tian'anmen (4 juin 1989) démontre que l'ouverture économique ne s'accompagne d'aucune libéralisation politique.
\end{exoresolu}

\begin{aretenir}[L'essentiel de la section]
\begin{itemize}
\item \textbf{1978} : Deng Xiaoping prend le pouvoir et lance les \cle{Quatre Modernisations} (agriculture, industrie, défense, sciences et technologies).
\item \textbf{Campagnes} : décollectivisation et système de responsabilité familiale $\rightarrow$ forte hausse des rendements et des revenus ruraux.
\item \textbf{1980} : création de quatre \cle{ZES} (Shenzhen, Zhuhai, Shantou, Xiamen) ; \textbf{1984} : ouverture de 14 villes côtières $\rightarrow$ afflux d'investissements étrangers et boom des exportations.
\item \textbf{1979} : politique de l'\cle{enfant unique} pour freiner la démographie.
\item \textbf{Continuité politique} : « économie socialiste de marché » sous monopole du PCC ; répression du printemps de Pékin (1979) et de \cle{Tian'anmen} (4 juin 1989).
\end{itemize}
\end{aretenir}

\begin{exercice}[Pour s'entraîner]
1) Définis les Quatre Modernisations et cite leurs quatre secteurs. 2) Explique en cinq lignes le fonctionnement du système de responsabilité familiale et pourquoi il a fait augmenter la production. 3) Sur une carte muette de Chine, place les quatre ZES de 1980 et justifie leur localisation littorale. 4) Pourquoi peut-on dire que la Chine de Deng est « ouverte économiquement mais fermée politiquement » ? Appuie ta réponse sur deux faits datés.
\end{exercice}

\begin{corrige}[Exercice]
1) Programme lancé par Deng Xiaoping en 1978 pour moderniser l'agriculture, l'industrie, la défense et les sciences et technologies. 2) La terre collective est louée aux familles qui doivent livrer un quota fixe à l'État et peuvent vendre le surplus au marché : l'effort est récompensé, donc la motivation et les rendements augmentent. 3) Shenzhen et Zhuhai (près de Hong Kong et Macao), Shantou et Xiamen (face à Taiwan) : le littoral offre des ports, la proximité des capitaux de la diaspora et l'accès aux marchés mondiaux. 4) Ouverte économiquement : ZES de 1980, investissements étrangers, exportations. Fermée politiquement : répression du Mur de la démocratie (1979) et massacre de Tian'anmen (4 juin 1989), monopole du PCC maintenu.
\end{corrige}

\coursec{L'émergence économique : la Chine, « atelier du monde » (années 1990-2010)}

Après la répression de Tian'anmen en 1989, la Chine traverse une période d'incertitude politique et de sanctions internationales. Pourtant, loin de marquer l'arrêt des réformes, cet événement précipite une décision stratégique majeure : Deng Xiaoping, alors âgé de 87 ans, entreprend en 1992 son célèbre « voyage dans le Sud » (Shenzhen, Zhuhai, Guangzhou, Shanghai). Ce périple relance l'ouverture et l'expérimentation capitaliste, scellant le destin de la Chine comme puissance économique mondiale. C'est le point de départ de l'émergence fulgurante que nous allons analyser.

\courssub{L'accélération des réformes après 1992 : le « voyage dans le Sud »}

Au début des années 1990, le camp conservateur au sein du Parti communiste chinois (PCC) freine les réformes, craignant une perte de contrôle idéologique. Deng Xiaoping, bien que retiré des fonctions officielles, conserve une autorité morale suprême. En janvier-février 1992, il effectue une tournée d'inspection dans les zones économiques spéciales (ZES) du Sud. Il y déclare notamment : « Celui qui ne réforme pas doit quitter la scène » et « Ouvrir encore plus largement ».

\begin{aretenir}[Le voyage de Deng Xiaoping dans le Sud (1992)]
Ce voyage n'est pas une simple visite ; c'est un acte politique décisif. Il légitime l'économie de marché socialiste, encourage l'investissement privé et étranger, et impose l'idée que « le développement est la seule vérité dure ». Il débloque la situation politique et lance la décennie de croissance la plus rapide de l'histoire chinoise.
\end{aretenir}

Les conséquences sont immédiates : le 14e Congrès du PCC (octobre 1992) officialise l'objectif d'une « économie de marché socialiste ». Les investissements directs étrangers (IDE) affluent : ils passent de \num{11} milliards USD en 1991 à plus de \num{50} milliards en 1993. La Chine devient la « usine du monde » pour l'assemblage et la sous-traitance.

\courssub{L'adhésion à l'OMC (2001) : intégration complète à la mondialisation}

L'aboutissement institutionnel de cette ouverture est l'adhésion de la Chine à l'Organisation mondiale du commerce (OMC) le 11 décembre 2001, après 15 années de négociations acharnées.

\begin{definition}[Adhésion à l'OMC]
L'adhésion à l'OMC engage la Chine à respecter les règles du commerce international : baisse drastique des droits de douane, suppression des barrières non tarifaires, ouverture des marchés des services (banque, télécoms, distribution) aux entreprises étrangères, protection de la propriété intellectuelle. En échange, la Chine bénéficie de la clause de la nation la plus favorisée (NPF) permanente, garantissant un accès stable aux marchés occidentaux.
\end{definition}

Cet engagement force une restructuration massive de l'appareil productif chinois : fermetures d'entreprises d'État non rentables, modernisation du cadre juridique, montée en compétence de la main-d'œuvre. C'est le « choc de modernisation » qui propulse la Chine au cœur des chaînes de valeur mondiales.

\courssub{Les facteurs de l'émergence : une conjonction favorable}

Quatre piliers structurels expliquent le « miracle chinois ». Ils ne sont pas le fruit du hasard mais le résultat de politiques délibérées.

\begin{popfigure}
\begin{tikzpicture}[
    mindmap,
    concept color=popblue,
    level 1/.append style={level distance=4cm, sibling angle=90},
    level 2/.append style={level distance=3cm, sibling angle=45},
    every node/.style={font=\small},
    grow cyclic
]
\node[concept, text width=3.5cm, align=center] {Facteurs de l'émergence chinoise}
    child[concept color=popgreen] { node[concept, text width=2.8cm, align=center] {Main-d'œuvre\\ abondante et bon marché}
        child { node[concept, text width=2.2cm, align=center] {Migrants ruraux (mingong)} }
        child { node[concept, text width=2.2cm, align=center] {Coût salarial < \num{5}\% USA (2000)} }
    }
    child[concept color=poporange] { node[concept, text width=2.8cm, align=center] {Investissements\\ étrangers (IDE) \& Transferts technologiques}
        child { node[concept, text width=2.2cm, align=center] {ZES : Shenzhen, Shanghai} }
        child { node[concept, text width=2.2cm, align=center] {Joint-ventures obligatoires} }
    }
    child[concept color=poppurple] { node[concept, text width=2.8cm, align=center] {Infrastructures\\ modernes \& logistique}
        child { node[concept, text width=2.2cm, align=center] {Ports : Shanghai, Ningbo} }
        child { node[concept, text width=2.2cm, align=center] {Autoroutes, TGV, aéroports} }
    }
    child[concept color=poppink] { node[concept, text width=2.8cm, align=center] {État stratège\\ \& Politique industrielle}
        child { node[concept, text width=2.2cm, align=center] {Planification quinquennale} }
        child { node[concept, text width=2.2cm, align=center] {Soutien aux champions nationaux} }
    };
\end{tikzpicture}
\legende{Schéma mental : les quatre piliers de l'émergence économique chinoise (années 1990-2010).}
\end{popfigure}

\begin{enumerate}
    \item \textbf{Main-d'œuvre abondante et bon marché :} Le « dividende démographique » (pic de la population en âge de travailler) et l'exode rural massif (\emph{mingong}, plus de \num{150} millions de migrants vers les villes côtières vers 2010) offrent une réserve de main-d'œuvre docile, peu coûteuse et non syndiquée.
    \item \textbf{Investissements étrangers et transferts de technologie :} La politique des « joint-ventures » (entreprises mixtes) impose le transfert de technologie en échange de l'accès au marché chinois. Les IDE apportent capital, gestion moderne et accès aux réseaux de distribution mondiaux.
    \item \textbf{Infrastructures modernes :} L'État investit massivement (routes, ports, barrages, puis TGV). Le port de Shanghai devient le premier port à conteneurs du monde en 2010. La logistique efficace réduit les coûts de transaction.
    \item \textbf{État stratège :} Contrairement au dogme néolibéral, l'État chinois garde le contrôle des secteurs stratégiques (banque, énergie, télécoms, transports) via les entreprises d'État (SOE) et pilote l'allocation du crédit vers l'exportation et l'infrastructure.
\end{enumerate}

\courssub{Le « Made in China » : du textile à l'électronique}

La spécialisation internationale de la Chine suit le « vol d'oies sauvages » : elle récupère les industries que les « tigres asiatiques » (Corée du Sud, Taïwan, Hong Kong, Singapour) délocalisent en montant en gamme.

\begin{exemplebox}[Spécialisation progressive : l'exemple de Shenzhen]
Dans les années 1980-1990, Shenzhen assemble des jouets, des textiles et des chaussures (faible valeur ajoutée). Dans les années 2000, elle devient la capitale mondiale de l'électronique grand public : assemblages pour Apple (Foxconn), Huawei, ZTE, DJI (drones). La ville passe du « \emph{made in China} » (fabriqué en Chine) au « \emph{designed in China} » (conçu en Chine) pour une partie croissante de sa production.
\end{exemplebox}

Les secteurs phares de l'exportation chinoise vers 2010 :
\begin{itemize}
    \item \textbf{Textile \& habillement :} \num{~30}\% des exportations mondiales.
    \item \textbf{Électronique \& machines :} Ordinateurs, téléphones, composants. La Chine devient le premier exportateur mondial de produits de haute technologie (selon la classification OCDE) dès 2004.
    \item \textbf{Jouets, meubles, chaussures :} Dominations quasi monopolistiques sur certains segments.
\end{itemize}

\courssub{Des résultats spectaculaires : croissance, PIB, commerce}

Les chiffres parlent d'eux-mêmes. La croissance moyenne annuelle du PIB réel oscille entre \num{9}\% et \num{10}\% sur trois décennies (1980-2010), un record historique pour une grande économie.

\begin{popfigure}
\begin{tikzpicture}
\begin{axis}[
    width=\linewidth,
    height=8cm,
    xlabel={Année},
    ylabel={Part dans le PIB mondial (\%)},
    xmin=1975, xmax=2025,
    ymin=0, ymax=20,
    xtick={1980,1990,2000,2010,2020},
    ytick={0,2,4,6,8,10,12,14,16,18,20},
    grid=major,
    grid style={popline},
    axis line style={popdark},
    tick style={popdark},
    legend pos=north west,
    legend style={fill=popcream, draw=popline},
]
\addplot[popcurve, very thick, mark=*, mark size=2.5pt, mark options={fill=popblue}] coordinates {
    (1980, 1.8)
    (1990, 1.7)
    (2000, 3.6)
    (2010, 9.2)
    (2020, 17.4)
};
\addlegendentry{Part de la Chine dans le PIB mondial (parité de pouvoir d'achat, FMI)}
% Ligne de tendance pour l'œil
\addplot[poporange, dashed, thick, domain=1978:2020] {0.0005*(x-1978)^2 + 1.8};
\end{axis}
\end{tikzpicture}
\legende{Évolution de la part de la Chine dans le PIB mondial (PPA). Source : FMI, World Economic Outlook. La courbe montre l'accélération post-2001 (OMC).}
\end{popfigure}

\begin{exoresolu}[Analyse statistique : La montée en puissance chinoise]
\textbf{Énoncé :} À partir des données suivantes, calculez le taux de variation moyen annuel de la part de la Chine dans le PIB mondial entre 2000 et 2010. Interprétez le résultat.
\begin{center}
\begin{tabular}{|c|c|}\hline
Année & Part dans le PIB mondial (PPA, \%) \\ \hline
2000 & 3,6 \\ \hline
2010 & 9,2 \\ \hline
\end{tabular}
\end{center}

\tcblower
\textbf{Solution détaillée :}
\begin{enumerate}
    \item \textbf{Formule du taux de variation moyen annuel (TVMA) :} $TVMA = \left( \frac{V_{final}}{V_{initial}} \right)^{\frac{1}{n}} - 1$, où $n$ est le nombre d'années.
    \item \textbf{Application :} $V_{initial} = \num{3,6}$, $V_{final} = \num{9,2}$, $n = 10$.
    \item \textbf{Calcul :} $TVMA = \left( \frac{\num{9,2}}{\num{3,6}} \right)^{\num{0,1}} - 1 = (\num{2,555})^{\num{0,1}} - 1 \approx \num{1,098} - 1 = \num{0,098}$.
    \item \textbf{Résultat :} Le taux de variation moyen annuel est d'environ \textbf{\num{9,8} \% par an}.
\end{enumerate}
\textbf{Interprétation :} La part de la Chine dans la richesse mondiale a presque triplé en dix ans, croissant près de deux fois plus vite que l'économie mondiale elle-même (dont la croissance moyenne est \num{~3-4}\%/an). Cela confirme une convergence rapide vers le centre de l'économie mondiale.
\end{exoresolu}

\begin{aretenir}[Résultats majeurs vers 2010]
\begin{itemize}
    \item \textbf{2e PIB mondial (nominal) :} La Chine dépasse le Japon en 2010 (environ \num{6 000} milliards USD contre \num{5 500} pour le Japon). Elle ne devra plus qu'aux États-Unis.
    \item \textbf{1er exportateur mondial :} Dépasse l'Allemagne en 2009 ; exportations \num{~1 200} milliards USD en 2009, \num{~1 580} milliards USD en 2010.
    \item \textbf{Réserves de change :} Les plus importantes du monde ($> \num{2 800}$ milliards USD en 2010), fruit des excédents commerciaux massifs.
    \item \textbf{Réduction de la pauvreté :} Plus de \textbf{800 millions de personnes} sorties de l'extrême pauvreté (seuil \num{1,90} \$/jour) depuis 1978, dont la majeure partie après 1990. C'est la plus grande réduction de pauvreté de l'histoire.
\end{itemize}
\end{aretenir}

\courssub{Les symboles de la puissance retrouvée : Pékin 2008 et Shanghai 2010}

La puissance économique s'accompagne d'une diplomatie du « soft power » pour légitimer le régime à l'intérieur et améliorer l'image à l'extérieur.

\begin{itemize}
    \item \textbf{Jeux olympiques de Pékin (août 2008) :} Préparés pendant 7 ans, coût estimé $> \num{40}$ milliards USD (infrastructures comprises). Cérémonie d'ouverture grandiose (Zhang Yimou), \num{51} médailles d'or (1er au tableau). Message : la Chine est une grande puissance moderne, pacifique et organisée.
    \item \textbf{Exposition universelle de Shanghai (mai-octobre 2010) :} Thème « Meilleure ville, meilleure vie ». \num{73} millions de visiteurs (record). Pavillon chinois (le « Crown of the East ») deux fois plus grand que le suivant. Vitrine de l'urbanisation chinoise et de ses ambitions technologiques.
\end{itemize}

\begin{savaistu}[Le « hard power » matériel : acier et ciment]
Pour construire ces infrastructures et ces villes, la Chine consomme des quantités phénoménales de matières premières. En 2010, elle produit \textbf{près de la moitié de l'acier mondial} (\num{~600} Mt sur \num{1 400} Mt, soit \num{~45}\%) et \textbf{plus de la moitié du ciment mondial} (\num{~1,8} Gt sur \num{3,3} Gt). Entre 2011 et 2013 seulement, la Chine a utilisé plus de ciment que les États-Unis pendant tout le 20e siècle. C'est la matérialité brute de l'émergence.
\end{savaistu}

\courssub{Les nouvelles ambitions : Nouvelles Routes de la Soie et montée en gamme}

Dès la fin de la période étudiée (2010), la Chine prépare l'après-« atelier du monde ». Le modèle basé sur l'exportation bas coût, l'investissement massif et l'endettement montre des signes d'essoufflement (surcapacités, bulle immobilière, vieillissement démographique).

\begin{definition}[Initiative « Belt and Road » (BRI) / Nouvelles Routes de la Soie]
Lancée par Xi Jinping en 2013 (Kazakhstan puis Indonésie), la BRI est une stratégie géoéconomique massive visant à connecter la Chine à l'Eurasie, l'Afrique et l'Europe par des corridors routiers, ferroviaires, portuaires et énergétiques. Objectifs : écouler les surcapacités industrielles chinoises (acier, ciment, BTP), sécuriser les approvisionnements énergétiques, internationaliser le yuan, et créer un espace d'influence politique alternatif à l'ordre occidental.
\end{definition}

\begin{exemplebox}[Montée en gamme technologique : « Made in China 2025 »]
Pour échapper au « piège du revenu intermédiaire », la Chine vise l'autonomie dans 10 secteurs de pointe : robotique, aéronautique, naval, ferroviaire haute vitesse, énergie nucléaire, biotech, matériaux nouveaux, véhicules électriques, équipements agricoles, TIC (Huawei, ZTE, 5G, IA). L'État subventionne massivement la R\&D (dépenses R\&D $> \num{2}\%$ du PIB dès 2014) et protège son marché intérieur pour laisser émerger des champions nationaux capables de concurrencer les GAFAM occidentaux.
\end{exemplebox}

\courssub{Les limites et contradictions de l'émergence (aperçu)}

Cette section sera développée dans la partie 4 du cours, mais il est crucial de ne pas dresser un tableau idyllique. L'émergence repose sur des déséquilibres structurels majeurs :

\begin{attention}[Piège fréquent : Confondre PIB total et niveau de vie]
\textbf{Erreur :} « La Chine est la 2e puissance mondiale, donc les Chinois sont riches. »
\textbf{Réalité :} En 2010, le PIB par habitant chinois est d'environ \num{4 500} USD (nominal), soit le \textbf{90e-100e rang mondial}, loin derrière les États-Unis (\num{~48 000} USD) ou la France (\num{~40 000} USD). La Chine est un pays \emph{à revenu intermédiaire} (selon la Banque mondiale). La richesse est concentrée sur la côte Est ; l'intérieur reste pauvre. L'inégalité (coefficient de Gini $> \num{0,47}$) est parmi les plus élevées d'Asie.
\end{attention}

\begin{itemize}
    \item \textbf{Modèle de croissance déséquilibré :} Sur-investissement (\num{~45-48}\% du PIB), sous-consommation des ménages (\num{~35}\% du PIB), dépendance aux exportations.
    \item \textbf{Coûts écologiques et sanitaires :} Pollution de l'air (PM2,5), de l'eau, des sols ; « villages du cancer » ; émissions de CO2 (1er émetteur mondial depuis 2006).
    \item \textbf{Vieillissement rapide :} Conséquence de la politique de l'enfant unique (1979-2015). La population en âge de travailler diminue depuis 2012. Fin du dividende démographique.
    \item \textbf{Risque financier :} Endettement des gouvernements locaux, bulle immobilière (Evergrande), système bancaire opaque.
    \item \textbf{Contrôle politique renforcé :} Depuis 2012 (Xi Jinping), recentralisation du pouvoir, répression au Xinjiang, loi de sécurité nationale à Hong Kong (2020), contrôle accru du secteur privé (Jack Ma / Ant Group, 2020).
\end{itemize}

\courssub{La Chine, l'Afrique et le Cameroun : une présence croissante (introduction)}

L'émergence chinoise a besoin de ressources (pétrole, minerais, bois, produits agricoles) et de marchés pour ses excédents de capitaux et de capacités industrielles. L'Afrique devient un partenaire privilégié.

\begin{exemplebox}[La Chine au Cameroun : chiffres clés (vers 2010-2020)]
\begin{itemize}
    \item \textbf{Commerce :} La Chine est le 1er partenaire commercial du Cameroun (excédent chinois massif). Exportations camerounaises : bois, coton, pétrole, minerais. Importations : machines, textiles, électronique, véhicules, BTP.
    \item \textbf{Infrastructures :} Barrage de Memve'ele (Sinohydro), Stade d'Olembé (CMEC), Route Kribi-Lolabé, projets routiers structurants (dont sections Yaoundé-Douala), Centrale à gaz de Kribi.
    \item \textbf{Financement :} Prêts concessionnels et commerciaux (Exim Bank of China, China Development Bank). Dette envers la Chine = part significative de la dette extérieure camerounaise.
    \item \textbf{Investissements privés :} Usines de transformation de bois, coton, riz (ex. Sodecoton partenariats), télécoms (Huawei, ZTE, China Unicom via Camtel).
\end{itemize}
Cette relation, qualifiée de « gagnant-gagnant » par Pékin, soulève des questions de souveraineté, de transfert de technologie réel, de normes sociales/environnementales et de soutenabilité de la dette, qui seront au cœur de la section 5.
\end{exemplebox}

\begin{methode}[Construire un diagramme en bâtons pour comparer des parts de marché]
Pour visualiser la montée en puissance de la Chine face à ses concurrents (USA, UE, Japon, Allemagne) :
\begin{enumerate}
    \item \textbf{Rassembler les données :} Tableau des parts de marché mondial à l'exportation pour 2000 et 2010 (source OMC).
    \item \textbf{Choisir l'échelle :} Axe des ordonnées (vertical) de 0\% à 20\% (max observé \num{~15-18}\%). Pas de 2\%.
    \item \textbf{Graduer les axes :} Axe des abscisses : pays/groupes. Axe des ordonnées : pourcentage.
    \item \textbf{Tracer les bâtons :} Deux bâtons groupés par pays (couleur 2000, couleur 2010). Largeur identique, espacement régulier.
    \item \textbf{Légender :} Titre clair (« Parts de marché mondial à l'exportation : 2000 vs 2010 »), source, unité (\%), légende des couleurs.
    \item \textbf{Analyser :} Croissance spectaculaire de la barre Chine (\num{3,9}\% $\to$ \num{10,4}\%), stagnation ou baisse des parts USA/Japon/UE.
\end{enumerate}
\end{methode}

\begin{exercice}[Entraînement : La Chine dans la mondialisation]
\begin{enumerate}
    \item En vous appuyant sur la courbe de la part de la Chine dans le PIB mondial (figure ci-dessus), identifiez les deux points de rupture majeurs de la courbe et associez-les à des événements politiques précis.
    \item Expliquez pourquoi l'adhésion à l'OMC en 2001 a été un accélérateur plus puissant que la création des ZES en 1980.
    \item À partir de l'exemple de Shenzhen, décrivez le processus de « montée en gamme » d'un territoire chinois.
    \item « La Chine est la deuxième puissance économique mondiale. » Cette affirmation est-elle suffisante pour qualifier la Chine de pays développé ? Justifiez votre réponse en utilisant la notion de PIB par habitant.
\end{enumerate}
\end{exercice}

\begin{corrige}[Exercice : La Chine dans la mondialisation]
\begin{enumerate}
    \item \textbf{1992 (Voyage de Deng / 14e Congrès) :} Rupture de pente nette, passage d'une croissance de la part mondiale quasi plate (\num{~1,7-1,8}\%) à une croissance soutenue. \textbf{2001 (Adhésion OMC) :} Accélération brutale de la courbe (convexité marquée), la part double en 9 ans (\num{4}\% $\to$ \num{9}\%).
    \item Les ZES étaient des zones expérimentales limitées géographiquement. L'OMC impose l'ouverture de \emph{tout} le territoire chinois aux règles du commerce mondial, verrouille les réformes par des engagements internationaux juridiquement contraignants (mécanisme de règlement des différends), et garantit l'accès au marché occidental (fin des menaces de sanctions annuelles type « Most Favored Nation » aux USA).
    \item Shenzhen : 1. Assemblage low-cost (textile, jouets, plastique) $\to$ 2. Électronique grand public (montage PC, téléphones pour marques étrangères) $\to$ 3. Conception/R\&D/propres marques (Huawei, DJI, BYD, Tencent) $\to$ 4. Exportation de capital et de normes (5G, paiement mobile, smart cities).
    \item Non. Le PIB total mesure la puissance économique globale (capacité d'influence, marché, puissance militaire potentielle). Le PIB par habitant mesure le niveau de vie moyen. En 2010, PIB/hab Chine $\approx$ \num{4 500} \$ (nominal) vs USA $\approx$ \num{48 000} \$. La Chine reste un pays en développement / à revenu intermédiaire. La puissance et le développement ne sont pas synonymes.
\end{enumerate}
\end{corrige}

\coursec{Les limites et les contradictions de l'émergence chinoise}

La croissance spectaculaire qui a fait de la Chine « l'atelier du monde » depuis les années 1990 ne s'est pas faite sans coûts. Si le \emph{Produit Intérieur Brut} (PIB) a été multiplié par dix en trois décennies, ce « miracle » repose sur des déséquilibres profonds qui menacent sa soutenabilité. Comprendre ces limites, c'est dépasser la seule lecture économique pour analyser les tensions sociales, écologiques, démographiques et politiques qui structurent la Chine d'aujourd'hui. Cette section prolonge logiquement la précédente : l'émergence économique a engendré ses propres contradictions.

\courssub{Des inégalités sociales et spatiales majeures : les deux Chine}

L'industrialisation côtière a créé une fracture territoriale sans précédent. Le développement suit le tracé des \emph{Zones Économiques Spéciales} (ZES) et des grandes métropoles portuaires (Shanghai, Shenzhen, Guangzhou), tandis que l'intérieur du pays — provinces du Centre, de l'Ouest et du Nord-Est — reste largement agricole et pauvre.

\begin{popfigure}
\begin{tikzpicture}[scale=0.9, every node/.style={font=\small}]
% Contour très simplifié de la Chine (polygone approximatif)
\draw[popdark, thick] (0,0) -- (8,0) -- (9,2) -- (9,5) -- (7,7) -- (5,8) -- (2,7) -- (0,5) -- (0,2) -- cycle;
% Ligne de partage (ligne Hu Huanyong / Heihe-Tengchong approximative)
\draw[poporange, dashed, thick] (1,1) -- (8,6);
% Littoral développé (hachures denses)
\fill[pattern=north east lines, pattern color=popblue!50] (0,0) -- (8,0) -- (9,2) -- (9,5) -- (7,7) -- (5,8) -- (3,6) -- (1,2) -- cycle;
% Intérieur rural (hachures légères)
\fill[pattern=north west lines, pattern color=poporange!30] (1,1) -- (8,6) -- (7,7) -- (5,8) -- (3,6) -- (1,2) -- cycle;
% Villes côtières
\foreach \x/\y/\name in {1.5/1.5/Shanghai, 6.5/1.5/Shenzhen, 7.5/3.5/Guangzhou, 4/5.5/Pékin} {
  \fill[popblue] (\x,\y) circle (3pt);
  \node[anchor=west, popdark] at (\x+0.2,\y) {\name};
}
% Villes intérieures
\foreach \x/\y/\name in {4/2.5/Chongqing, 5.5/4.5/Xi'an, 2.5/4/Chengdu} {
  \fill[poporange] (\x,\y) circle (3pt);
  \node[anchor=west, popdark] at (\x+0.2,\y) {\name};
}
% Flèches migrations (mingong)
\draw[->, popgreen, very thick] (2.5,2) -- (4,1.5) node[midway, above, sloped, popdark] {Mingong};
\draw[->, popgreen, very thick] (3.5,3.5) -- (5.5,2.5);
\draw[->, popgreen, very thick] (1.5,3.5) -- (3,2.5);
% Flèche Nord
\draw[->, thick, popdark] (9.5, 6) -- (9.5, 7) node[above] {N};
% Légende manuelle
\node[anchor=south west, draw, fill=popcream, rounded corners, inner sep=4pt, align=center] at (0.2,0.2){
  \begin{tabular}{ll}
    \textcolor{popblue}{\rule{5pt}{5pt}} Littoral industrialisé (ZES) & \textcolor{poporange}{\rule{5pt}{5pt}} Intérieur rural / ancien industriel \\
    \textcolor{popgreen}{\rule{10pt}{2pt}} Migrations \emph{mingong} & \textcolor{poporange}{\rule{5pt}{5pt}} Ligne de fracture (Hu Huanyong)
  \end{tabular}
};
\end{tikzpicture}
\legende{Croquis de synthèse : la Chine du littoral développé et la Chine de l'intérieur. La ligne pointillée (ligne Hu Huanyong) sépare la Chine dense et riche (Est/Sud-Est) de la Chine vaste et pauvre (Ouest/Nord). Les flèches vertes symbolisent le flux massif des \emph{mingong} vers les usines côtières.}
\end{popfigure}

\begin{methode}[Réaliser un croquis de synthèse en histoire-géographie]
\emph{Objectif :} Représenter l'essentiel d'un phénomène spatial sans noyer la carte sous les détails.
\begin{enumerate}
\item \textbf{Analyser le sujet} : identifier l'idée force (ici : fracture côtière / intérieur + migrations).
\item \textbf{Choisir 3 à 4 figurés maximum} : un fond (hachures/couleurs pour les deux Chine), une ligne de fracture, des points (villes), des flèches (flux).
\item \textbf{Hiérarchiser l'information} : Titre clair, légende structurée (fond, lignes, points, flèches), orientation (Nord) et échelle indicative.
\item \textbf{Soigner l'exécution} : traits nets, coloriage propre, écriture lisible. Éviter le superflu (fleuves, relief détaillé) sauf s'ils expliquent le phénomène.
\end{enumerate}
\end{methode}

\begin{exemplebox}[Chiffres clés des inégalités]
Le coefficient de Gini (mesure des inégalités de revenus) chinois dépasse \num{0,46} (seuil d'alerte international), contre environ \num{0,35} en France. Le revenu par habitant de Shanghai ou Pékin approche celui du Portugal, tandis que celui du Guizhou ou du Yunnan reste proche de celui du Ghana. \textbf{L'écart de revenu entre le littoral et l'intérieur rural peut aller de 1 à 3, voire 1 à 4.}
\end{exemplebox}

Cette fracture spatiale alimente une migration interne massive : les \cle{mingong} (travailleurs paysans-ouvriers). On en compte environ \textbf{300 millions} (chiffre incluant les migrants de courte et longue durée, source Bureau national des statistiques 2023). Ils quittent l'intérieur pour travailler dans les usines, les chantiers et les services des métropoles côtières, souvent dans des conditions précaires (heures supplémentaires imposées, absence de protection sociale, logement insalubre).

\begin{definition}[Le hukou (户口) : le permis de résidence]
Le \cle{hukou} est le système chinois d'enregistrement des résidences qui lie les droits sociaux (école publique gratuite, santé subventionnée, logement social, retraites) au lieu d'habitation officiel (rural ou urbain). Héréditaire à l'origine, il crée une \textbf{citoyenneté à deux vitesses} : un migrant rural en ville n'a pas accès aux services urbains pour lui et ses enfants, le maintenant dans une précarité structurelle malgré sa contribution à l'économie urbaine.
\end{definition}

\begin{attention}[Piège à éviter : confusion statut / présence]
Ne pas confondre \emph{population résidente} (ceux qui vivent sur place) et \emph{population enregistrée} (ceux qui ont le hukou local). Les statistiques officielles de population urbaine sous-estiment souvent la réalité urbaine car des millions de \emph{mingong} y vivent sans y être enregistrés. La réforme du hukou (assouplissement dans les villes moyennes, maintien strict dans les mégapoles comme Pékin/Shanghai) progresse mais reste très inégale.
\end{attention}

\courssub{Une urgence environnementale planétaire}

Le modèle « usine du monde » repose sur une consommation massive d'énergie (charbon environ 55\% du mix électrique) et de ressources, avec une réglementation environnementale longtemps laxiste ou mal appliquée localement.

\begin{itemize}
\item \textbf{Pollution de l'air} : Les pics de \emph{PM2,5} (particules fines) à Pékin et dans la plaine du Nord de la Chine ont rendu l'air irrespirable des dizaines de jours par an (« airpocalypse » de 2013). Cela a forcé l'État à lancer un « Plan d'action pour la prévention et le contrôle de la pollution de l'air » (2013) avec des résultats réels (baisse des PM2,5), mais la situation reste critique l'hiver.
\item \textbf{Pollution de l'eau et des sols} : Environ 60\% des nappes phréatiques sont polluées ; des « villages du cancer » apparaissent le long des rivières industrielles. La réhabilitation des sols coûte des milliards.
\item \textbf{Émissions de CO$_2$} : La Chine est devenue le \textbf{premier émetteur mondial de CO$_2$} (annuel) vers 2006, dépassant les États-Unis (mais pas en émissions cumulées historiques ni par habitant). Elle s'est engagée à atteindre le pic d'émissions avant 2030 et la neutralité carbone en 2060 (« double objectif »), ce qui implique une transformation radicale de son modèle énergétique.
\end{itemize}

\begin{exoresolu}[Analyse d'un document : la courbe des émissions de CO$_2$]
\textbf{Document :} Graphique montrant l'évolution des émissions annuelles de CO$_2$ (en milliards de tonnes) : Chine (croissance exponentielle 2000-2013, plateau 2013-2019, reprise post-Covid), États-Unis (lente baisse), UE (baisse), Inde (hausse).
\textbf{Question :} Expliquer la responsabilité différenciée de la Chine dans le réchauffement climatique.
\textbf{Solution :}
\begin{enumerate}
\item \textbf{Constat chiffré} : La Chine émet $\approx$ 11-12 Gt CO$_2$/an (environ 30\% du total mondial), soit le double des USA.
\item \textbf{Causes structurelles} : Charbon (électricité, sidérurgie, ciment), exportations (émissions « délocalisées » par l'Occident), urbanisation rapide (béton, acier).
\item \textbf{Nuance indispensable (responsabilité partagée)} :
\begin{itemize}
\item Émissions par habitant : Chine ($\approx$ 8 t/hab) < USA ($\approx$ 14 t/hab) < Moyenne mondiale ($\approx$ 4,7 t/hab).
\item Émissions cumulées (1850-2020) : USA $\approx$ 20\% ; Chine $\approx$ 11\%.
\item Commerce : $\approx$ 15-20\% des émissions chinoises produisent des biens exportés (consommation occidentale).
\end{itemize}
\item \textbf{Conclusion} : La Chine est l'acteur \emph{actuel} incontournable pour limiter le réchauffement (+1,5°C impossible sans elle), mais la justice climatique exige une différenciation des efforts (financement, transfert technologique).
\end{enumerate}
\end{exoresolu}

\courssub{Le défi démographique : le vieillissement avant l'enrichissement}

La \cle{politique de l'enfant unique} (1979-2015), imposée pour freiner la croissance démographique et favoriser l'accumulation de capital, a brutalement réduit le taux de fécondité (indice synthétique $\approx$ 1,2-1,3 aujourd'hui, bien en dessous du seuil de renouvellement 2,1).

\begin{savaistu}[L'assouplissement progressif face au mur démographique]
Face au vieillissement accéléré et à la baisse de la population active, le régime a assoupli la règle :
\begin{itemize}
\item 2013 : Deux enfants si un parent est enfant unique.
\item \textbf{2015 : Fin officielle de l'enfant unique} $\rightarrow$ politique des \textbf{deux enfants}.
\item \textbf{2021 : Politique des \textbf{trois enfants}} + mesures d'accompagnement (congés parentaux, subventions, interdiction des cours privés payants pour réduire le coût de l'éducation).
\end{itemize}
\textbf{Résultat :} La natalité continue de chuter (9,56 millions de naissances en 2022, record bas depuis 1949). Le coût de la vie, l'éducation, le logement et la place des femmes sur le marché du travail freinent plus que la loi.
\end{savaistu}

\begin{propriete}[Conséquences structurelles du vieillissement]
\begin{enumerate}
\item \textbf{Érosion de la main-d'œuvre} : La population en âge de travailler (15-64 ans) diminue depuis 2012. Fin du « dividende démographique » (main-d'œuvre abondante et bon marché) $\rightarrow$ hausse des salaires, perte de compétitivité « low-cost », relocalisation partielle vers Vietnam, Inde, Mexique.
\item \textbf{Explosion de la charge de dépendance} : Ratio retraités/actifs en forte hausse. Système de retraites par répartition sous tension (déficit des caisses provinciales).
\item \textbf{Phénomène « 4-2-1 »} : Un enfant unique doit potentiellement subvenir à 2 parents + 4 grands-parents. Charge familiale immense, frein à la consommation et à la natalité.
\item \textbf{Déséquilibre des sexes} : Préférence culturelle pour les garçons + politique de l'enfant unique $\rightarrow$ avortements sélectifs $\rightarrow$ $\approx$ 30 millions d'hommes « en trop » (difficultés de mariage, tensions sociales).
\end{enumerate}
\end{propriete}

\courssub{L'absence de démocratie : un régime autoritaire high-tech}

La croissance économique n'a pas entraîné de libéralisation politique, contredisant la thèse de la « modernisation » (Lipset). Le \cle{Parti communiste chinois (PCC)} conserve le monopole du pouvoir.

\begin{itemize}
\item \textbf{Institutions} : Pas de séparation des pouvoirs. Le PCC contrôle l'État (Assemblée populaire), l'armée (Armée populaire de libération) et la justice. Le Secrétaire général du PCC (Xi Jinping depuis 2012) cumule les titres (Président, Chef de l'armée). Suppression de la limite de deux mandats (2018) $\rightarrow$ pouvoir personnel renforcé.
\item \textbf{Contrôle de l'information} : « Grand Firewall » (censure Internet), contrôle des médias, disparition de la presse d'investigation. Surveillance numérique généralisée : reconnaissance faciale, crédit social (notation des citoyens/entreprises), application de traçage (sanitaire puis sécuritaire).
\item \textbf{Répression des dissidences} : Mouvement de la place Tian'anmen (1989) réprimé ; mouvement pro-démocratie à Hong Kong (2019) étouffé par la \emph{Loi de sécurité nationale} (2020) ; internement massif de Ouïghours et autres minorités turcophones au \textbf{Xinjiang} (camps de « rééducation », travail forcé, stérilisations) qualifié de « génocide » par plusieurs parlements occidentaux ; politique d'assimilation forcée au \textbf{Tibet} et en \textbf{Mongolie intérieure}.
\end{itemize}

\begin{attention}[Erreur fréquente : Puissance économique $\neq$ Démocratie]
Ne pas écrire : « La Chine est devenue démocratique car elle est riche. » \textbf{Faux.} La Chine combine \textbf{libéralisme économique} (marché, propriété privée, ouverture au commerce) et \textbf{régime autoritaire} (parti unique, censure, répression). C'est un modèle de « capitalisme d'État autoritaire » ou « capitalisme bureaucratique ». Le « contrat social » implicite : le PCC assure la prospérité et la stabilité $\leftrightarrow$ la population accepte l'absence de libertés politiques. Ce contrat est testé par le ralentissement économique.
\end{attention}

\courssub{Tensions internationales : la rivalité sino-américaine}

L'émergence chinoise bouleverse l'ordre international post-1945 (hégémonie US). La relation est devenue structurellement conflictuelle (« piège de Thucydide » : puissance montante vs puissance établie).

\begin{center}
\begin{tabularx}{\linewidth}{|l|X|X|}
\hline
\rowcolor{popblueL} \textbf{Domaine} & \textbf{Points de friction} & \textbf{Enjeux} \\
\hline
\textbf{Commerce / Tech} & Guerre commerciale (droits de douane Trump/Biden), sanctions sur puces (semi-conducteurs), Huawei, TikTok, investissements US en Chine. & Souveraineté technologique, découplage (\emph{decoupling}) ou « dérisquage » (\emph{de-risking}). \\
\hline
\textbf{Taïwan} & Revendication chinoise (province rebelle), loi anti-sécession (2005), exercices militaires, statut ambigu US (One China policy + Taiwan Relations Act : vente d'armes). & Risque de conflit armé majeur (blocus / invasion). Ligne rouge absolue pour Pékin. \\
\hline
\textbf{Mer de Chine} & Revendications « ligne en neuf traits » (îles Spratleys, Paracels), militarisation récifs, arbitrage La Haye 2016 (rejeté par Chine), tensions Philippines/Vietnam. & Contrôle routes maritimes (commerce, énergie), projection stratégique. \\
\hline
\textbf{Droits de l'Homme} & Xinjiang, Hong Kong, Tibet $\rightarrow$ sanctions occidentales (UE, US, UK, Canada) sur officiels/entités. & Valeurs universelles vs non-ingérence / souveraineté. \\
\hline
\end{tabularx}
\end{center}

\begin{popfigure}
\begin{tikzpicture}
\begin{axis}[
    width=\linewidth,
    height=8cm,
    ybar=5pt,
    bar width=18pt,
    ymin=0, ymax=90000,
    ylabel={PIB par habitant (USD courants, approx. 2023)},
    ylabel style={font=\small},
    symbolic x coords={Cameroun, Chine, États-Unis},
    xtick=data,
    xticklabel style={rotate=0, font=\small, anchor=north},
    nodes near coords,
    nodes near coords align={vertical},
    nodes near coords style={font=\tiny, rotate=90, anchor=west, yshift=2pt},
    ymajorgrids=true,
    grid style={popline},
    legend style={at={(0.5,-0.25)}, anchor=north, legend columns=-1, font=\small},
]
\addplot[popblue, fill=popblue!60] coordinates {(Cameroun,1600) (Chine,12600) (États-Unis,81600)};
\legend{PIB/hab (USD 2023 approx.)}
\end{axis}
\end{tikzpicture}
\legende{Comparaison des PIB par habitant (USD courants, estimations FMI/Banque Mondiale 2023). La Chine reste un pays à revenu intermédiaire (tranche supérieure), loin du niveau US, mais très au-dessus du Cameroun. L'écart Chine/US $\approx$ 1 à 6,5 ; Chine/Cameroun $\approx$ 1 à 8. Source : FMI World Economic Outlook 2024.}
\end{popfigure}

\begin{aretenir}[Synthèse : Les 5 contradictions majeures de l'émergence chinoise]
\begin{enumerate}
\item \textbf{Spatiale/Sociale} : Littoral riche / Intérieur pauvre ; 300 M \emph{mingong} précarisés par le \emph{hukou}.
\item \textbf{Écologique} : 1er émetteur CO$_2$, pollution air/eau/sols critique ; transition verte indispensable mais coûteuse.
\item \textbf{Démographique} : Vieillissement rapide (politique enfant unique), baisse population active, natalité effondrée malgré 3 enfants autorisés.
\item \textbf{Politique} : Régime autoritaire (PCC unique), surveillance numérique, répression minorités (Xinjiang, Tibet), pas de transition démocratique.
\item \textbf{Géopolitique} : Rivalité systémique USA (Tech, Taïwan, Mer de Chine, Valeurs) ; dépendance mutuelle économique mais découplage stratégique.
\end{enumerate}
\end{aretenir}

\coursec{La Chine, l'Afrique et le Cameroun : une présence croissante}

L'Afrique est devenue un partenaire stratégique majeur pour la Chine, à la fois comme source de matières premières (pétrole, minerais, bois, produits agricoles), comme marché d'écoulement pour ses produits manufacturés et ses excédents de capacités industrielles (BTP, acier, ciment), et comme relais diplomatique (ONU, Forum sur la coopération sino-africaine - FCSA).

\courssub{Le cadre institutionnel : le FCSA et les Nouvelles Routes de la Soie}

Depuis 2000, le \cle{Forum sur la coopération sino-africaine (FCSA)} (FOCAC en anglais) structure la relation tous les trois ans (sommets Pékin 2006, 2018, Dakar 2021, Pékin 2024). La Chine y annonce des enveloppes massives (prêts concessionnels, dons, lignes de crédit, investissements directs). L'initiative \cle{Ceinture et Route} (BRI / Nouvelles Routes de la Soie), lancée en 2013, intègre l'Afrique comme maillon essentiel (voies maritimes via ports, corridors ferroviaires/routières).

\begin{exemplebox}[Chiffres clés de la relation Chine-Afrique (2023)]
\begin{itemize}
\item \textbf{Commerce bilatéral} : $\approx$ 282 milliards USD (Chine = 1er partenaire commercial de l'Afrique depuis 2009).
\item \textbf{Structure} : Exportations Afrique $\rightarrow$ Chine : matières premières (pétrole, cuivre, cobalt, fer, bois, coton). Importations Chine $\rightarrow$ Afrique : machines, électronique, textiles, véhicules, acier.
\item \textbf{Investissements (IDE)} : Stock d'IDE chinois en Afrique $\approx$ 50-60 milliards USD (secteurs : mines, infrastructures, manufacturier, services).
\item \textbf{Aide/Dette} : La Chine est le premier créancier bilatéral de nombreux pays africains (dont le Cameroun). Problématique de la « dette opaque » et des clauses de confidentialité.
\end{itemize}
\end{exemplebox}

\courssub{Le cas camerounais : un partenariat « tous azimuts »}

Le Cameroun, « Afrique en miniature », illustre la densité de la présence chinoise.

\begin{center}
\begin{tabularx}{\linewidth}{|l|X|}
\hline
\rowcolor{popblueL} \textbf{Secteur} & \textbf{Projets / Réalisations emblématiques} \\
\hline
\textbf{Infrastructures portuaires} & Port en eau profonde de \textbf{Kribi} (phase 1 achevée 2014, phase 2 en cours) : opérateur CHEC (China Harbour Engineering Company). Plateforme logistique pour l'hinterland (Tchad, RCA). \\
\hline
\textbf{Énergie} & Barrage de \textbf{Lom Pangar} (exploitation 2017, Sinohydro/CWE), barrage de \textbf{Mekin}, centrale à gaz de \textbf{Kribi} (216 MW, CMEC). Objectif : résorber le déficit énergétique. \\
\hline
\textbf{Transport routier/fer} & Routes : Yaoundé-Nsimalen, Kribi-Lolodorf, routes du Nord. Études/réhabilitation \textbf{Camrail} (chemin de fer). Projet corridor Douala-Ndjamena. \\
\hline
\textbf{Équipements sociaux/sportifs} & Stades d'\textbf{Olembe} (CAN 2021, CMEC) et de \textbf{Japoma} (Douala). Hôpital de référence de \textbf{Yaooundé} (don), Hôpital de \textbf{Douala} (prêt). \\
\hline
\textbf{Industrie / ZES} & Zone économique spéciale de \textbf{Kribi} (bois, transformation). Usines de transformation du coton, du bois, cimenteries (Cimencam - actionnariat chinois partiel via Sinoma). \\
\hline
\textbf{Coopération humaine} & Instituts Confucius (Yaoundé, Douala), bourses d'études (des centaines d'étudiants camerounais en Chine/an), missions médicales, formation militaire/police. \\
\hline
\end{tabularx}
\end{center}

\begin{attention}[Enjeux et critiques : une relation asymétrique ?]
\begin{itemize}
\item \textbf{Déséquilibre commercial} : Excédent structurel chinois. Le Cameroun exporte du brut (bois en grumes, coton brut, pétrole, aluminium) et importe du manufacturé. Faible transfert de technologie et de valeur ajoutée locale.
\item \textbf{Endettement} : La dette envers la Chine (créancier bilatéral n°1) pèse sur la soutenabilité budgétaire (FMI). Négociations de rééchelonnement (Cadre commun G20, initiative DSSI).
\item \textbf{Gouvernance / Environnement / Social} : Critiques sur l'opacité des contrats (clauses de confidentialité), l'impact écologique (déforestation pour le bois, mines), les conditions de travail (salaires bas, sécurité), l'emploi de main-d'œuvre chinoise au détriment du local.
\item \textbf{Diplomatie} : Soutien chinois au Cameroun aux Nations Unies (non-ingérence), alignement fréquent sur les votes d'intérêt chinois (Taïwan, Xinjiang, Hong Kong).
\end{itemize}
\end{attention}

\begin{exercice}[Entraînement Bac : Commentaire de document]
\textbf{Document :} Extrait du discours de Xi Jinping au 20e Congrès du PCC (octobre 2022) : « Nous devons promouvoir une développement de haute qualité [...] réaliser la modernisation chinoise [...] assurer la sécurité nationale [...] ne jamais hégémonisme. »
\textbf{Consigne :} En vous appuyant sur le document et vos connaissances, montrez que le modèle chinois actuel tente de résoudre ses contradictions internes par un renforcement du contrôle politique et une réorientation économique, tout en affirmant sa puissance extérieure.
\textbf{Pistes :} 1. Contexte (20e Congrès, 3e mandat Xi). 2. « Développement haute qualité » = réponse aux limites (environnement, inégalités, low-cost). 3. « Sécurité nationale » = primauté du politique/contrôle (hukou, tech, Xinjiang, Covid-zéro). 4. « Modernisation chinoise » = modèle alternatif à l'occidental (autoritaire + marché). 5. « Pas d'hégémonisme » vs réalités (Taïwan, Mer de Chine, Nouvelles Routes de la Soie).
\end{exercice}

\begin{exercice}[Entraînement Bac : Sujet de dissertation / composition]
\textbf{Sujet :} « La Chine, puissance émergente : entre affirmation internationale et fragilités internes. » Discutez cette affirmation en vous appuyant sur l'évolution de la Chine depuis 1978.
\textbf{Plan type :}
\begin{enumerate}
\item \textbf{Une émergence économique spectaculaire, fondement de la puissance.} (Réformes Deng, ouverture, ZES, usine du monde, 2e PIB mondial, BRI/FCSA, influence Afrique/Cameroun).
\item \textbf{Des fragilités internes structurelles qui menacent le modèle.} (Inégalités spatiales/hukou, urgence écologique, piège démographique, régime autoritaire/surveillance, endettement interne/bulle immobilière).
\item \textbf{Une affirmation géopolitique assumée mais contestée.} (Rivalité US, Taïwan/Mer de Chine, leadership Sud global/BRICS, soft power, mais isolement relatif sur valeurs/droits de l'homme, dépendance technologique/énergétique).
\end{enumerate}
\textbf{Conclusion :} La Chine est une puissance \emph{incontournable} mais \emph{incomplète} (puissance économique/étatique forte, puissance normative/sociale faible). Son avenir dépend de la résolution de ses contradictions internes (transition haute qualité, gestion vieillissement, légitimité du PCC).
\end{exercice}

\coursec{La Chine, l'Afrique et le Cameroun : une présence croissante}

Après avoir analysé les contradictions internes du modèle chinois -- inégalités sociales, coût environnemental, tensions politiques --, nous devons désormais observer comment cette puissance émergente projette son influence à l'extérieur. L'Afrique, et le Cameroun en particulier, constituent un terrain d'observation privilégié pour comprendre la stratégie globale de la Chine : sécuriser ses approvisionnements, écouler sa production industrielle et asseoir son leadership diplomatique. Cette section propose une enquête documentée sur cette relation asymétrique mais structurante.

\courssub{Le FOCAC : cadre institutionnel de la coopération sino-africaine}

\begin{definition}[Forum sur la coopération sino-africaine (FOCAC)]
Le \cle{FOCAC} (Forum on China-Africa Cooperation) est une instance de dialogue politique et de coopération économique collective créée en \textbf{2000} à Pékin. Il réunit la Chine et les pays africains (hors Eswatini, qui reconnaît Taïwan) tous les trois ans, alternativement en Chine et en Afrique. Il structure les relations par des plans d'action triennaux (infrastructures, santé, agriculture, formation, paix et sécurité).
\end{definition}

Contrairement aux relations bilatérales classiques, le FOCAC offre à la Chine un cadre \emph{multilatéral} pour négocier avec le continent comme un bloc, tout en conservant des accords bilatéraux très forts. Pour les pays africains, c'est une tribune pour revendiquer un « partenariat gagnant-gagnant » et diversifier leurs partenaires traditionnels (France, UE, institutions de Bretton Woods).

\begin{savaistu}[Une diplomatie du sommet]
Depuis 2000, les sommets FOCAC (Pékin 2000, Addis-Abeba 2003, Pékin 2006, Charm el-Cheikh 2009, Pékin 2012, Johannesburg 2015, Pékin 2018, Dakar 2021) ont vu les engagements financiers chinois passer de quelques milliards à \textbf{60 milliards USD} (2018) puis \textbf{40 milliards USD} (2021, dont 10 milliards de facilités de crédit), marquant une maturation de la relation : moins de « gros chèques » spectaculaires, plus de projets ciblés (santé, numérique, vert).
\end{savaistu}

\courssub{Les moteurs de la présence chinoise en Afrique : ressources, marchés, diplomatie}

L'engagement chinois en Afrique répond à trois impératifs stratégiques, souvent imbriqués :

\begin{enumerate}
    \item \textbf{Sécuriser les approvisionnements en matières premières.} La Chine est le premier consommateur mondial de nombreux minerais (cobalt, cuivre, fer) et d'hydrocarbures. L'Afrique lui fournit environ \textbf{20\% de ses importations de pétrole} et une part massive de minerais critiques (cobalt de RDC, fer d'Afrique du Sud, bauxite de Guinée). Les infrastructures (chemins de fer, ports) financées par la Chine relient souvent les zones d'extraction aux ports d'exportation (ex. : corridor Lobito en Angola, chemin de fer Addis-Abeba--Djibouti).
    \item \textbf{Conquérir des débouchés pour son industrie.} Face à la saturation de son marché intérieur et aux barrières protectionnistes occidentales, l'Afrique représente un marché de \textbf{1,4 milliard de consommateurs} pour les produits manufacturés chinois (textile, électronique, machines, véhicules, BTP). La Zone de libre-échange continentale africaine (ZLECAf, 2021) intéresse au premier chef Pékin.
    \item \textbf{Asseoir un leadership diplomatique alternatif.} La Chine promeut le principe de \emph{non-ingérence dans les affaires intérieures}, séduisant des régimes critiqués par l'Occident sur la gouvernance. Elle obtient en retour des soutiens aux Nations Unies (Taïwan, Xinjiang, Hong Kong, Mer de Chine méridionale) et promeut son modèle de développement (Initiative « la Ceinture et la Route », lancée en 2013).
\end{enumerate}

\begin{exemplebox}[Chiffres clés : la Chine, premier partenaire commercial de l'Afrique]
Depuis \textbf{2009}, la Chine a détrôné les États-Unis et l'Union européenne comme premier partenaire commercial de l'Afrique.
\begin{itemize}
    \item Volume des échanges : \textbf{282 milliards USD} (2023), contre 11 milliards en 2000 ($\times$25).
    \item Structure : L'Afrique exporte principalement des matières premières (pétrole, minerais, bois, coton) et importe des produits manufacturés (machines, électronique, textile, acier).
    \item Investissements directs (IDE) : Stock d'IDE chinois en Afrique $\approx$ 50 milliards USD (2022), concentrés dans l'extraction, les infrastructures, puis de plus en plus les services et la manufacture.
\end{itemize}
\end{exemplebox}

\courssub{La présence chinoise au Cameroun : infrastructures visibles et produits du quotidien}

Le Cameroun, « Afrique en miniature », illustre la densité de la coopération sino-camerounaise, formalisée dès 1971 mais accélérée depuis les années 2000.

\begin{popfigure}
\begin{tikzpicture}[scale=0.9, every node/.style={font=\small}]
    % Contour très simplifié du Cameroun (polygone approximatif)
    \draw[popdark, thick, fill=popcream!30] 
        (0,0) -- (7,0) -- (8,2) -- (7.5,5) -- (5,6.5) -- (2,6) -- (0.5,4) -- (0,2) -- cycle;
    
    % Villes principales
    \node[circle, fill=popblue, inner sep=2pt, label=above:\textbf{Yaoundé}] (Y) at (3.5,4.5) {};
    \node[circle, fill=popblue, inner sep=2pt, label=right:\textbf{Douala}] (D) at (6.5,1.5) {};
    \node[circle, fill=popblue, inner sep=2pt, label=below:\textbf{Kribi}] (K) at (5.5,0.5) {};
    \node[circle, fill=popblue, inner sep=2pt, label=left:\textbf{Memve'ele}] (M) at (1.5,1.5) {};
    \node[circle, fill=popblue, inner sep=2pt, label=above:\textbf{Bertoua}] (B) at (5.5,4) {};
    \node[circle, fill=popblue, inner sep=2pt, label=above:\textbf{Garoua}] (G) at (4,6) {};
    \node[circle, fill=popblue, inner sep=2pt, label=above:\textbf{Maroua}] (Ma) at (2.5,6.2) {};

    % Projets chinois (étoiles rouges)
    \node[star, star points=5, star point ratio=2.25, fill=poporange, draw=popdark, inner sep=3pt, label={[popdark]90:Barrage de Memve'ele}] at (M) {};
    \node[star, star points=5, star point ratio=2.25, fill=poporange, draw=popdark, inner sep=3pt, label={[popdark]0:Port en eau profonde de Kribi}] at (K) {};
    \node[star, star points=5, star point ratio=2.25, fill=poporange, draw=popdark, inner sep=3pt, label={[popdark]90:Stade d'Olembe (Yaoundé)}] at ($(Y)+(0.5,0.5)$) {};
    \node[star, star points=5, star point ratio=2.25, fill=poporange, draw=popdark, inner sep=3pt, label={[popdark]-90:Stade de Japoma (Douala)}] at ($(D)+(0.5,-0.5)$) {};
    \node[star, star points=5, star point ratio=2.25, fill=poporange, draw=popdark, inner sep=3pt, label={[popdark]180:Route Yaoundé--Bertoua (tronçon)}] at (4.5,4.2) {};
    \node[star, star points=5, star point ratio=2.25, fill=poporange, draw=popdark, inner sep=3pt, label={[popdark]0:Route Kribi--Lolabé}] at (6,1) {};
    \node[star, star points=5, star point ratio=2.25, fill=poporange, draw=popdark, inner sep=3pt, label={[popdark]90:Assemblée nationale (Yaoundé)}] at ($(Y)+(-0.5,0.5)$) {};
    
    % Flèche Nord
    \node[draw=popdark, fill=popblueL, inner sep=3pt] at (7.5, 5.5) {\textbf{N}};
    
    % Légende manuelle
    \begin{scope}[shift={(0,-1.5)}]
        \node[star, star points=5, star point ratio=2.25, fill=poporange, draw=popdark, inner sep=3pt] (leg1) at (0,0) {};
        \node[right] at (leg1.east) {Grands projets d'infrastructures (entreprises chinoises)};
        \node[circle, fill=popblue, inner sep=2pt] (leg2) at (0,-0.6) {};
        \node[right] at (leg2.east) {Villes principales};
    \end{scope}
\end{tikzpicture}
\legende{Carte schématique du Cameroun : localisation des grands projets réalisés avec la participation d'entreprises chinoises (TP : compléter avec les axes routiers, le chemin de fer, les zones minières).}
\end{popfigure}

\begin{itemize}
    \item \textbf{Infrastructures de transport et énergie :} Le \textbf{barrage de Memve'ele} (211 MW, mis en service 2019, Sinohydro) sur le fleuve Ntem (Sud) ; le \textbf{port en eau profonde de Kribi} (phase 1 achevée 2014, CHEC/China Harbour) et sa zone industrielle, futur hub logistique ; les \textbf{stades de la CAN 2021} (Olembe à Yaoundé par CCECC, Japoma à Douala par PICC) ; le \textbf{Palais des Congrès} et le nouveau \textbf{Palais de l'Assemblée nationale} à Yaoundé ; de nombreux tronçons routiers (Yaoundé--Bertoua, Kribi--Lolabé, routes urbaines).
    \item \textbf{Produits manufacturés :} Les marchés camerounais (Marché Central, Mokolo, Nkolndongo, Douala) sont inondés de produits chinois : textiles (pagnes, vêtements), chaussures, électronique (téléphones Transsion/Tecno/Infinix dominent le marché), motos, quincaillerie, ustensiles ménagers. La Chine est le \textbf{premier fournisseur du Cameroun} (environ 20--25\% des importations).
    \item \textbf{Entreprises et main-d'œuvre :} Grandes entreprises d'État (SOE) : CCECC (chemins de fer, routes, stades), Sinohydro (barrages), CHEC (ports), CAMC, ZTE/Huawei (télécoms, ville intelligente). Présence croissante d'entreprises privées et de commerçants chinois dans les villes.
\end{itemize}

\begin{savaistu}[Le stade d'Olembe, symbole de la coopération]
Le stade Paul Biya d'Olembe (60 000 places), inauguré en 2021 pour la CAN, a été construit par \textbf{CCECC} (China Civil Engineering Construction Corporation) pour un coût d'environ 160 milliards FCFA. Il illustre la « diplomatie du stade » chinoise : infrastructure prestigieuse, livraison rapide, financement par prêt concessionnel, mais aussi polémiques sur la qualité des finitions et la gestion post-événementielle.
\end{savaistu}

\courssub{Financements, dette et le piège de la « diplomatie du piège à dette » ?}

Le modèle de financement chinois diffère de l'aide occidentale (dons, prêts très concessionnels du FMI/Banque mondiale) : ce sont majoritairement des \textbf{prêts concessionnels} (taux bas $\approx$ 2\%, longue maturité 15--20 ans, différé 5--7 ans) et des \textbf{prêts commerciaux} (taux marché, garantis par des ressources naturelles -- \emph{resource-backed loans}).

\begin{aretenir}[Structure de la dette camerounaise vis-à-vis de la Chine]
\begin{itemize}
    \item La Chine est le \textbf{premier créancier bilatéral} du Cameroun (devant la France), représentant $\approx$ \textbf{30 à 35\% de la dette extérieure publique} (environ 2 500--3 000 milliards FCFA sur $\approx$ 8 000--9 000 milliards FCFA de dette extérieure totale, chiffres 2023-2024).
    \item Projets phares financés par dette : Kribi (port, centrale à gaz), Memve'ele, routes, stades, fibre optique (Huawei), centrale de Lom Pangar (partiellement).
    \item Risque : Le Cameroun est classé en \textbf{risque élevé de surendettement} par le FMI (2023-2024). Le service de la dette chinoise pèse sur le budget et limite les marges de manœuvre (ajustement structurel FMI 2021-2024).
\end{itemize}
\end{aretenir}

\begin{attention}[Piège rhétorique : « La Chine achète tout » / « Piège à dette »]
\begin{itemize}
    \item \textbf{Éviter :} « Les Chinois nous colonisent / achètent le pays. » \rightarrow \textbf{Raisonner en économiste politique :} La Chine prête pour financer \emph{ses} entreprises (SOE) qui construisent \emph{ses} infrastructures pour extraire \emph{ses} ressources ou accéder \emph{à ses} marchés. C'est une relation contractuelle, asymétrique, mais signée par des États souverains.
    \item \textbf{Nuancer le « debt-trap diplomacy » :} Aucune saisie d'actif stratégique (port, aéroport) n'a eu lieu en Afrique pour défaut de paiement (contrairement au port de Hambantota au Sri Lanka, souvent cité mais contexte différent). En revanche, la \emph{conditionnalité ressources-contre-infrastructures} (ex. : pétrole coté pour Kribi) hypothèque les revenus futurs.
\end{itemize}
\end{attention}

\courssub{Opportunité ou nouvelle dépendance ? Le grand débat}

\begin{center}
\begin{tabularx}{\linewidth}{|l|X|X|}
\hline
\rowcolor{popblueL} \textbf{Dimension} & \textbf{Arguments « Opportunité / Partenariat gagnant-gagnant »} & \textbf{Arguments « Dépendance / Risques »} \\ \hline
\textbf{Infrastructures} & Comblent le déficit massif (énergie, ports, routes) ; rapidité d'exécution ; coûts souvent inférieurs aux appels d'offres occidentaux. & Qualité parfois inégale ; maintenance négligée ; clauses de confidentialité des contrats ; main-d'œuvre chinoise importée (peu de transfert de compétences). \\ \hline
\textbf{Économie / Emploi} & Prix bas pour les consommateurs (téléphones, vêtements, motos) ; création d'emplois sur chantiers ; installation d'usines (ex. : cimenterie Cimencam modernisée, usines de transformation bois/agro). & \textbf{Concurrence déloyale} : effondrement du textile local (pagnes, coton), fermeture d'unités de transformation (riz, bois) ; importation de main-d'œuvre qualifiée chinoise ; faible création de valeur locale. \\ \hline
\textbf{Ressources / Environnement} & Investissements dans l'exploitation minière (fer Mbalam/Nabeba -- projet bloqué, cobalt, or) et forestière. & Extraction « prédatrice » : bois brut exporté (interdit théoriquement mais contourné), minerais sans transformation locale, dégradation environnementale (minage artisanal encadré par des opérateurs chinois). \\ \hline
\textbf{Diplomatie / Souveraineté} & Alternative au « pré carré » français / conditionnalités FMI (gouvernance, droits de l'homme) ; soutien à l'ONU ; formation de cadres (bourse d'études Chine). & Alignement diplomatique attendu (vote ONU, Taïwan) ; dépendance technologique (Huawei/ZTE : « ville intelligente », surveillance, données) ; opacité des contrats (clauses de non-divulgation). \\ \hline
\end{tabularx}
\end{center}

\courssub{Comparaison avec les partenaires traditionnels : France, Banque mondiale, Union européenne}

\begin{exoresolu}[Analyse comparative : Chine vs Partenaires traditionnels au Cameroun]
\textbf{Énoncé :} À partir des données ci-dessous, comparez la nature de l'engagement de la Chine et de la France (premier partenaire historique) au Cameroun.

\textbf{Données :}
\begin{itemize}
    \item \textbf{France :} 1er investisseur historique (stock IDE), présence dans l'énergie (Perenco, TotalEnergies), grande distribution (CFAO, Carrefour), banque (BICEC, Afriland -- capitaux français), télécoms (Orange). Aide au développement (AFD) : dons/prêts très concessionnels \emph{conditionnés} à la gouvernance, climat, genre. Accords de défense. Monnaie : Franc CFA (garantie par Trésor français).
    \item \textbf{Chine :} 1er partenaire commercial (échanges), 1er créancier bilatéral, 1er constructeur d'infrastructures publiques. Prêts concessionnels/commerciaux \emph{non conditionnés} à la gouvernance politique. Pas de présence militaire formelle (sauf formation). Entreprises d'État dominantes.
    \item \textbf{Banque mondiale / FMI :} Prêts concessionnels (IDA) / non concessionnels ; \emph{conditionnalité forte} (réformes structurelles : transparence budgétaire, secteur énergie, filet social, climat). Appui technique et normatif.
\end{itemize}

\textbf{Solution détaillée :}
\begin{enumerate}
    \item \textbf{Nature du flux :} France = IDE (capital productif, long terme, reprise de bénéfices) + Aide publique au développement (APD). Chine = Prêts souverains (créance publique) + IDE minier/infrastructure (SOE). BM/FMI = Prêts programmes/projets + Expertise normative.
    \item \textbf{Conditionnalité :} France/BM/FMI = \textbf{Conditionnalité politique et économique} (gouvernance, droits humains, transparence, climat). Chine = \textbf{Conditionnalité commerciale et diplomatique} (accès ressources, marchés, soutien ONU, non-reconnaissance Taïwan). Principe de non-ingérence affiché.
    \item \textbf{Secteurs cibles :} France = Secteurs matures, services, rente (pétrole, distribution, finance). Chine = Infrastructures lourdes (BTP, énergie, transport), extraction, numérique. BM/FMI = Budget de l'État, secteurs sociaux, réforme administration, énergie.
    \item \textbf{Risque pour le Cameroun :} France = Dépendance monétaire (FCFA), capture de la rente par grands groupes, désindustrialisation historique. Chine = Surendettement public, enclaves extractives, désindustrialisation par concurrence importations, opacité contractuelle. BM/FMI = Austérité budgétaire, retrait État services sociaux, inégalité.
    \item \textbf{Complémentarité / Concurrence :} Sur les infrastructures (barrages, ports, routes), la Chine a \textbf{remplacé} les entreprises françaises (Bouygues, Vinci, Razel) souvent jugées trop chères/lentes. Sur le numérique, Huawei/ZTE concurrencent Orange/Thales. Sur la dette, la Chine refuse le Club de Paris (réstructuration coordonnée) préférant le bilatéral, compliquant les plans du FMI (ex. : restructuration dette Zambie, Ghana, Tchad -- Cadre commun G20 dont Chine est partie prenante réticente).
\end{enumerate}
\end{exoresolu}

\courssub{Importations camerounaises : la place dominante de la Chine}

\begin{popfigure}
\begin{tikzpicture}
    % Diagramme circulaire - Répartition indicative des origines des importations du Cameroun (moyenne 2020-2023)
    % Données : Chine 22%, France 12%, Nigeria 10%, Belgique 6%, Inde 5%, USA 4%, Allemagne 3%, Autres 38%
    \def\angleChine{79.2}
    \def\angleFrance{43.2}
    \def\angleNigeria{36}
    \def\angleBelgique{21.6}
    \def\angleInde{18}
    \def\angleUSA{14.4}
    \def\angleAllemagne{10.8}
    \def\angleAutres{136.8}
    
    \begin{scope}[shift={(0,0)}]
        % Parts du camembert
        \draw[fill=poporange] (0,0) -- (3,0) arc (0:\angleChine:3) -- cycle;
        \draw[fill=popblue] (0,0) -- (\angleChine:3) arc (\angleChine:\angleChine+\angleFrance:3) -- cycle;
        \draw[fill=popgreen] (0,0) -- (\angleChine+\angleFrance:3) arc (\angleChine+\angleFrance:\angleChine+\angleFrance+\angleNigeria:3) -- cycle;
        \draw[fill=poppurple] (0,0) -- (\angleChine+\angleFrance+\angleNigeria:3) arc (\angleChine+\angleFrance+\angleNigeria:\angleChine+\angleFrance+\angleNigeria+\angleBelgique:3) -- cycle;
        \draw[fill=poppink] (0,0) -- (\angleChine+\angleFrance+\angleNigeria+\angleBelgique:3) arc (\angleChine+\angleFrance+\angleNigeria+\angleBelgique:\angleChine+\angleFrance+\angleNigeria+\angleBelgique+\angleInde:3) -- cycle;
        \draw[fill=popgold] (0,0) -- (\angleChine+\angleFrance+\angleNigeria+\angleBelgique+\angleInde:3) arc (\angleChine+\angleFrance+\angleNigeria+\angleBelgique+\angleInde:\angleChine+\angleFrance+\angleNigeria+\angleBelgique+\angleInde+\angleUSA:3) -- cycle;
        \draw[fill=popmuted] (0,0) -- (\angleChine+\angleFrance+\angleNigeria+\angleBelgique+\angleInde+\angleUSA:3) arc (\angleChine+\angleFrance+\angleNigeria+\angleBelgique+\angleInde+\angleUSA:\angleChine+\angleFrance+\angleNigeria+\angleBelgique+\angleInde+\angleUSA+\angleAllemagne:3) -- cycle;
        \draw[fill=popline] (0,0) -- (\angleChine+\angleFrance+\angleNigeria+\angleBelgique+\angleInde+\angleUSA+\angleAllemagne:3) arc (\angleChine+\angleFrance+\angleNigeria+\angleBelgique+\angleInde+\angleUSA+\angleAllemagne:360:3) -- cycle;
        
        % Légende à droite (positionnement manuel pour stabilité)
        \begin{scope}[shift={(4.5, 2.8)}]
            \node[anchor=west, font=\small] at (0,0) {\color{poporange}\rule{0.5cm}{0.3cm} Chine : 22\%};
            \node[anchor=west, font=\small] at (0,-0.7) {\color{popblue}\rule{0.5cm}{0.3cm} France : 12\%};
            \node[anchor=west, font=\small] at (0,-1.4) {\color{popgreen}\rule{0.5cm}{0.3cm} Nigeria : 10\%};
            \node[anchor=west, font=\small] at (0,-2.1) {\color{poppurple}\rule{0.5cm}{0.3cm} Belgique : 6\%};
            \node[anchor=west, font=\small] at (0,-2.8) {\color{poppink}\rule{0.5cm}{0.3cm} Inde : 5\%};
            \node[anchor=west, font=\small] at (0,-3.5) {\color{popgold}\rule{0.5cm}{0.3cm} USA : 4\%};
            \node[anchor=west, font=\small] at (0,-4.2) {\color{popmuted}\rule{0.5cm}{0.3cm} Allemagne : 3\%};
            \node[anchor=west, font=\small] at (0,-4.9) {\color{popline}\rule{0.5cm}{0.3cm} Autres : 38\%};
        \end{scope}
    \end{scope}
\end{tikzpicture}
\legende{Répartition indicative des origines des importations du Cameroun (moyenne 2020-2023, sources : INS, Douanes, Comtrade). La Chine s'impose comme le premier fournisseur ($\approx$22\%), devant la France ($\approx$12\%). Le Nigeria figure haut grâce aux flux pétroliers.}
\end{popfigure}

\begin{methode}[Rédiger une conclusion argumentée sur la coopération sino-africaine]
\textbf{Problématique type :} « La coopération sino-camerounaise est-elle une opportunité de développement ou un facteur de nouvelle dépendance ? »

\begin{enumerate}
    \item \textbf{Rappeler la problématique} et définir les termes (opportunité = infrastructures, prix bas, diversification ; dépendance = dette, désindustrialisation, alignement diplomatique).
    \item \textbf{Synthétiser les acquis par grands axes équilibrés} (au moins 3) :
        \begin{itemize}
            \item \textit{Axe 1 : Infrastructures et connectivité} (réel apport, mais opacité, qualité, main-d'œuvre).
            \item \textit{Axe 2 : Commerce et industrie} (pouvoir d'achat préservé, mais concurrence fatale au tissu local, primarisation des exportations).
            \item \textit{Axe 3 : Finances et souveraineté} (alternative au FMI/Club de Paris, mais surendettement, clauses ressources, dépendance technologique).
        \end{itemize}
    \item \textbf{Nuancer :} La responsabilité incombe aussi à l'État camerounais (négociation, planification, contrôle, politique industrielle locale, diversification des partenaires).
    \item \textbf{Ouvrir sur une perspective :} « L'enjeu pour le Cameroun n'est pas de choisir \emph{entre} la Chine et l'Occident, mais de \emph{négocier} avec tous sur la base d'une stratégie nationale claire (SND30, ZLECAf) : exiger transfert de technologie, transformation locale des ressources, transparence contractuelle, et utiliser la concurrence entre partenaires pour obtenir de meilleures conditions. »
\end{enumerate}
\end{methode}

\begin{exercice}[Sujet d'entraînement : Commentaire de document]
\textbf{Document :} Extrait du discours de Xi Jinping au Sommet FOCAC de Dakar (novembre 2021) : « La Chine continuera de soutenir l'Afrique dans la lutte contre le Covid-19 [...] mettra en œuvre neuf programmes [...] dont le programme de promotion du commerce, le programme d'investissement, le programme de réduction de la pauvreté... La Chine accordera 10 milliards de dollars de facilités de crédit aux institutions financières africaines... »

\textbf{Questions :}
\begin{enumerate}
    \item Situez le document (auteur, date, contexte, nature).
    \item Identifiez les trois grands axes de la coopération annoncée.
    \item Montrez en quoi cette annonce illustre la stratégie chinoise en Afrique (ressources, marchés, diplomatie).
    \item En vous appuyant sur vos connaissances, discutez les limites de ce type d'engagement pour un pays comme le Cameroun.
\end{enumerate}
\end{exercice}

\begin{corrige}[Exercice n°1]
\begin{enumerate}
    \item \textbf{Analyse du document :} Discours de Xi Jinping (Président chinois), Sommet FOCAC 8 (Dakar, Sénégal), 29 nov. 2021. Contexte : Post-Covid, relance africaine, rivalité Sino-US. Nature : Discours officiel d'engagement politique et financier.
    \item \textbf{Trois axes :} (1) Santé / Covid (vaccins, hôpital CDC Afrique) ; (2) Commerce / Investissement / Pauvreté (9 programmes) ; (3) Financement (10 mds USD facilités de crédit + annulation dettes sans intérêt arrivant à échéance pour PMA).
    \item \textbf{Illustration de la stratégie :} Santé = soft power / image responsable. Commerce/Investissement = débouchés pour excédents chinois (BTP, manufacture) + sécurisation ressources (programme « promotion du commerce » = importations africaines). Financement = maintien de la dépendance par la dette / liquidités pour projets chinois (SOE), alternative au FMI.
    \item \textbf{Limites pour le Cameroun :} (a) Facilités de crédit $\neq$ dons : alourdit dette (déjà risque élevé FMI). (b) « Programme promotion commerce » : favorise exportations brutes (bois, coton, pétrole) vs transformation locale (ZLECAf). (c) Non-ingérence = pas de pression sur gouvernance/environnement (risques projets miniers/forestiers). (d) Concurrence produits chinois (textile, acier) tue industrie naissante camerounaise. (e) Conditionnalité implicite : alignement votes ONU, 5G Huawei, données numériques.
\end{enumerate}
\end{corrige}

\newpage

\coursec{Activité d'intégration}

\courssub{Objectifs de l'activité}
Cette activité d'intégration vise à mobiliser l'ensemble des savoirs et savoir-faire de la leçon sur l'émergence chinoise. Elle permet à l'élève de :
\begin{itemize}
    \item \textbf{Analyser et croiser des documents hétérogènes} (statistiques, cartographiques, iconographiques) pour construire une argumentation historique et géographique.
    \item \textbf{Maîtriser le langage cartographique} : choisir des figurés pertinents (points, surfaces, flèches), hiérarchiser l'information et réaliser une légende organisée.
    \item \textbf{Manipuler des données quantitatives} : passer d'un tableau statistique à une représentation graphique (diagramme en bâtons) en respectant les règles de construction (échelle, unités, titres).
    \item \textbf{Ancrer les connaissances dans le contexte camerounais} : identifier et localiser les grands projets de coopération sino-camerounais pour en mesurer l'impact concret sur le territoire national.
    \item \textbf{Rédiger une synthèse argumentée} (10 à 15 lignes) répondant à une problématique complexe, en articulant le processus historique interne à la Chine (réformes, ouverture) et ses effets externes (investissements, infrastructure, dette).
\end{itemize}

\courssub{Situation}
\textbf{Contexte :} Dans le cadre de la préparation du \emph{Baccalauréat technique d'Histoire}, vous intégrez une cellule d'analyse prospective du Ministère de l'Économie, de la Planification et de l'Aménagement du Territoire (MINEPAT). Votre mission est de produire une \emph{note de synthèse cartographiée et illustrée} à l'intention des décideurs, intitulée : \og \emph{La Chine : trajectoire d'une puissance émergente et implications pour le Cameroun} \fg{}.

Le dossier documentaire mis à votre disposition comprend quatre pièces :
\begin{enumerate}
    \item \textbf{Document 1 (Carte de base) :} Un fond de carte de la Chine (échelle 1/50 000 000) représentant le littoral Pacifique, les frontières, les grands fleuves (Fleuve Jaune, Fleuve Bleu/Yangtsé, Fleuve des Perles) et la localisation des principales métropoles (Pékin, Shanghai, Guangzhou, Shenzhen, Tianjin, Wuhan, Chongqing, Hong Kong).
    \item \textbf{Document 2 (Données statistiques) :} Tableau de l'évolution du PIB nominal de la Chine (en milliards de dollars US courants) de 1978 à 2020.
    \begin{center}
    \begin{tabularx}{\linewidth}{|l|X|}\hline
    \textbf{Année} & \textbf{PIB (Mds \$ US)} \\ \hline
    1978 & 150 \\ \hline
    1980 & 191 \\ \hline
    1990 & 361 \\ \hline
    2000 & 1 211 \\ \hline
    2001 (entrée OMC) & 1 339 \\ \hline
    2005 & 2 286 \\ \hline
    2010 & 6 101 \\ \hline
    2015 & 11 062 \\ \hline
    2020 & 14 723 \\ \hline
    \end{tabularx}
    \end{center}
    \item \textbf{Document 3 (Texte de référence) :} Extrait du discours de Deng Xiaoping lors de sa tournée du Sud (janvier 1992) : \og \emph{Le développement est la seule vérité dure. [...] Il faut s'ouvrir davantage, pas se refermer. [...] Les zones économiques spéciales sont des fenêtres de l'ouverture, des fenêtres pour acquérir la technologie, la gestion moderne, les connaissances.} \fg{} Liste des Zones Économiques Spéciales (ZES) de première génération (1980) : Shenzhen, Zhuhai, Shantou, Xiamen. Extension majeure : Hainan (1988), Shanghai-Pudong (1990), et ouverture des villes côtières.
    \item \textbf{Document 4 (Dossier photographique Cameroun) :} Trois clichés légendés :
    \begin{itemize}
        \item \textit{Photo A :} Port en eau profonde de Kribi (phase 1 inaugurée 2014), réalisé par \emph{China Harbour Engineering Company (CHEC)}, financé à 85\% par la \emph{Exim Bank of China}. Capacité : 1,5 million de conteneurs EVP/an.
        \item \textit{Photo B :} Barrage hydroélectrique de Memve'ele (mis en service 2019), construit par \emph{Sinohydro}, financé par prêt concessionnel chinois. Puissance : 211 MW. Ligne de transport 225 kV vers Yaoundé.
        \item \textit{Photo C :} Échangeur de Nsimalen (Yaoundé) / Palais des Congrès de Yaoundé (rénové/construit par entreprises chinoises), symbolisant l'urbanisme et la diplomatie.
    \end{itemize}
    \item \textbf{Document 5 (Carte de base) :} Fond de carte du Cameroun (échelle 1/10 000 000) avec les 10 régions, le réseau hydrographique principal, les villes de Yaoundé, Douala, Kribi, Memve'ele (Sud), et le tracé du pipeline Tchad-Cameroun.
\end{enumerate}

\courssub{Consignes}
En groupe de 3 ou 4 élèves, vous réaliserez le travail suivant sur papier Canson (format A3 pour les cartes, A4 pour le graphique et le texte) ou sur support numérique (logiciel de cartographie / tableur). La restitution orale désignera un rapporteur par groupe.

\begin{enumerate}
    \item \textbf{Cartographie de l'émergence chinoise (Document 1 + 3) :} Sur le fond de carte de la Chine fourni, construisez une carte intitulée \og \emph{La façade maritime chinoise : moteur de l'émergence (1980-2020)} \fg{}. Votre légende, organisée en 3 items maximum, doit faire apparaître :
    \begin{itemize}
        \item Les 4 ZES historiques de 1980 (figuré surfacique ou ponctuel spécifique).
        \item Les 3 principales métropoles portuaires et industrielles du littoral (Shanghai, Guangzhou, Shenzhen) hiérarchisées par leur poids démographique/économique (figuré ponctuel proportionnel ou ordonné).
        \item Le sens de l'ouverture sur le monde (flèches de flux maritimes vers l'Asie, l'Europe, l'Amérique, l'Afrique).
    \end{itemize}
    \item \textbf{Analyse quantitative : le décollage économique (Document 2) :} À partir du tableau fourni, tracez un \textbf{diagramme en bâtons verticaux} représentant l'évolution du PIB chinois (1978-2020). Respectez les règles : axes gradués (échelle adaptée aux valeurs), titre précis, unité indiquée, source citée. Identifiez graphiquement (flèche ou couleur distincte) la rupture de 1978 et l'accélération post-2001 (OMC).
    \item \textbf{La présence chinoise au Cameroun (Documents 4 + 5) :} Sur le fond de carte du Cameroun, localisez avec précision les trois projets illustrés par les photos (Kribi, Memve'ele, Yaoundé). Utilisez des figurés ponctuels distincts par nature d'infrastructure (port, énergie, urbain/administratif). Réalisez une légende de 3 items. Tracez une flèche de flux commercial majeure (Douala/Kribi $\rightarrow$ Chine : bois, coton, pétrole ; Chine $\rightarrow$ Douala/Kribi : machines, textiles, électronique).
    \item \textbf{Rédaction de la synthèse (10 à 15 lignes) :} En vous appuyant \textbf{exclusivement} sur les cartes, le graphique et les documents du dossier, répondez à la problématique : \og \textbf{Comment la Chine est-elle devenue une puissance mondiale et que change-t-elle pour le Cameroun ?} \fg{} Votre réponse doit articuler : le tournant de 1978 (politique), le levier des ZES/OMC (économique), la mutation structurelle (usine du monde), et la traduction concrète au Cameroun (infrastructures, financement, dépendance asymétrique).
\end{enumerate}

\courssub{Ressources à mobiliser}
\begin{methode}[Rappel méthodologique : Construire un diagramme en bâtons]
\begin{enumerate}
    \item \textbf{Choisir l'échelle :} Valeur max (14 723) $\rightarrow$ arrondir à 15 000. Hauteur dispo $\approx$ 15 cm $\rightarrow$ 1 cm pour 1 000 Mds \$. (Ou 0,5 cm pour 500 Mds \$).
    \item \textbf{Tracer les axes :} Abscisses = années (échelle temporelle non proportionnelle car dates irrégulières, mais ordre chronologique respecté ; espacement égal pour lisibilité). Ordonnées = PIB (Mds \$).
    \item \textbf{Tracer les bâtons :} Largeur constante, espacement régulier. Colorier différemment les périodes : 1978-1990 (pré-réforme/amorce), 1990-2001 (réformes), 2001-2020 (OMC/boom).
    \item \textbf{Légender :} Titre centré, source (Banque Mondiale / FMI), unité.
\end{enumerate}
\end{methode}

\begin{methode}[Rappel : Langage cartographique (Figurés)]
\begin{itemize}
    \item \textbf{Figuré ponctuel (point, carré, triangle) :} Localisation précise (ville, site industriel, barrage). Taille $\propto$ importance (proportionnel) ou nature (qualitatif).
    \item \textbf{Figuré surfacique (hachures, couleur, tramage) :} Étendue d'un phénomène (ZES, zone agricole, forêt).
    \item \textbf{Figuré linéaire (flèche, trait) :} Flux, axes de communication, frontières. Épaisseur $\propto$ intensité du flux.
    \item \textbf{Légende :} Titre, items classés (surfaciques $\rightarrow$ linéaires $\rightarrow$ ponctuels), figurés reproduits \textbf{exactement} comme sur la carte.
\end{itemize}
\end{methode}

\begin{aretenir}[Savoirs clés de la leçon à réinvestir]
\begin{itemize}
    \item \textbf{1966-1976 :} \textbf{Révolution culturelle} (Mao Zedong) : chaos politique, purge des cadres, économie paralysée, rupture avec l'URSS (1960), isolement international. Nécessité d'un changement radical après la mort de Mao (1976).
    \item \textbf{1978 :} III\textsuperscript{e} Plenum du XI\textsuperscript{e} Comité Central $\rightarrow$ Deng Xiaoping lance \emph{Réformes et Ouverture} (\textit{Gaige Kaifang}). Fin de l'économie planifiée rigide, responsabilité familiale en agriculture, création ZES.
    \item \textbf{ZES (Shenzhen, etc.) :} \cle{Incitations fiscales}, autonomie de gestion, main-d'œuvre bon marché, attraction des \cle{IDE} (Investissements Directs Étrangers), transfert de technologie. \og \emph{Une fenêtre ouverte sur le monde} \fg.
    \item \textbf{2001 :} Adhésion à l'\cle{OMC} $\rightarrow$ accès aux marchés mondiaux, explosion exportations, \og \emph{Atelier du monde} \fg.
    \item \textbf{Stratégie \og Aller vers l'extérieur \fg (\textit{Zou Chu Qu}) :} Investissements massifs en Afrique (ressources, infrastructures, marchés) via \emph{Exim Bank}, \emph{China Development Bank}, entreprises d'État (SOE).
    \item \textbf{Cameroun :} Partenaire stratégique (zone CEMAC). Projets structurants (énergie, transport, urbain). Enjeux : désenclavement, industrialisation vs endettement, balance commerciale déficitaire, transfert de technologie limité.
\end{itemize}
\end{aretenir}

\courssub{Démarche et corrigé commenté}
\begin{exoresolu}[Corrigé type de l'activité d'intégration]
\textbf{Étape 1 : Carte de la façade maritime chinoise (Consigne 1)}
\begin{popfigure}
\begin{tikzpicture}[scale=0.65, every node/.style={font=\footnotesize}]
    % Coastline simplified (polygon)
    \draw[popblue!80!popdark, line width=1.2pt, fill=popblue!5!popcream] 
        (0,0) -- (1,0.2) -- (2,0.5) -- (3,1) -- (4,1.5) -- (5,2) -- (6,2.5) -- (7,3) -- (8,3.5) -- (9,3.8) -- (10,4) -- (11,3.8) -- (12,3.5) -- (13,3) -- (14,2.5) -- (15,2) -- (16,1.5) -- (17,1) -- (18,0.5) -- (19,0) -- cycle;
    % North Arrow
    \draw[->, popdark, line width=1pt] (17.5, 3.5) -- ++(0, 0.6) node[above, popdark] {N};
    % Rivers
    \draw[popblue, line width=0.8pt] (4,1.5) -- (3,0.5) -- (2,0) node[left, popdark] {Fleuve Jaune};
    \draw[popblue, line width=1pt] (8,3.5) -- (7,2) -- (6,1) -- (5,0.5) node[below, popdark] {Yangtsé};
    \draw[popblue, line width=0.8pt] (12,3) -- (13,1.5) -- (14,0.5) node[right, popdark] {Fleuve des Perles};
    % SEZ 1980 (Surface hachurée)
    \tikzset{sezstyle/.style={pattern=north east lines, pattern color=poporange!80!popdark, fill=poporange!15, draw=poporange!80!popdark, thick}}
    \filldraw[sezstyle] (1.5,0.5) circle (0.4) node[above=2pt] {Shenzhen};
    \filldraw[sezstyle] (2.5,0.3) circle (0.35) node[below=2pt] {Zhuhai};
    \filldraw[sezstyle] (3.5,0.8) circle (0.35) node[above=2pt] {Shantou};
    \filldraw[sezstyle] (5.5,2.2) circle (0.35) node[above=2pt] {Xiamen};
    \filldraw[sezstyle] (1.5,-0.5) circle (0.5) node[below=4pt] {Hainan (88)};
    \filldraw[sezstyle] (8.5,3.8) circle (0.4) node[above=2pt] {Pudong (90)};
    % Metropoles (Points proportionnels)
    % Beijing (North) - not on coast but capital
    \node[draw=poppurple, thick, circle, minimum size=6pt, fill=poppurple!30] at (10, 5.5) {}; \node[above, poppurple] at (10, 5.8) {Pékin};
    % Shanghai (Yangtze delta)
    \node[draw=poppurple, thick, circle, minimum size=10pt, fill=poppurple!50] at (8.5, 3.5) {}; \node[left, poppurple] at (8, 3.5) {Shanghai};
    % Tianjin
    \node[draw=poppurple, thick, circle, minimum size=7pt, fill=poppurple!40] at (11, 4.2) {}; \node[above, poppurple] at (11, 4.5) {Tianjin};
    % Wuhan (inland)
    \node[draw=poppurple, thick, circle, minimum size=6pt, fill=poppurple!30] at (8, 1.5) {}; \node[left, poppurple] at (7.5, 1.5) {Wuhan};
    % Chongqing (inland)
    \node[draw=poppurple, thick, circle, minimum size=6pt, fill=poppurple!30] at (5, 0.5) {}; \node[below, poppurple] at (5, 0.1) {Chongqing};
    % Guangzhou / Pearl River Delta
    \node[draw=poppurple, thick, circle, minimum size=9pt, fill=poppurple!50] at (13, 1.8) {}; \node[right, poppurple] at (13.5, 1.8) {Guangzhou};
    % Shenzhen (SEZ + Metro)
    \node[draw=poporange!80!popdark, thick, circle, minimum size=9pt, fill=poporange!50] at (1.8, 0.6) {}; \node[below, poporange!80!popdark] at (1.8, 0.1) {Shenzhen};
    % Hong Kong
    \node[draw=popgreen!70!popdark, thick, circle, minimum size=7pt, fill=popgreen!30] at (2.5, 0.0) {}; \node[below, popgreen!70!popdark] at (2.5, -0.4) {Hong Kong};
    % Arrows (Flux maritimes)
    \draw[->, poporange, line width=2pt, >=Stealth] (15, 2.5) -- ++(2, 0.5) node[right, popdark] {Europe / Afrique};
    \draw[->, poporange, line width=2pt, >=Stealth] (15, 2.0) -- ++(2, -0.5) node[right, popdark] {Moyen-Orient};
    \draw[->, poporange, line width=1.5pt, >=Stealth] (0, 0.5) -- ++(-2, 0) node[left, popdark] {Amérique (Pacifique)};
    \draw[->, poporange, line width=1.5pt, >=Stealth] (0, 3.5) -- ++(-2, 0) node[left, popdark] {Asie du Nord-Est};
    % Legend Box (simulated)
    \node[draw=popline, fill=popcream, rounded corners, inner sep=5pt, anchor=south west, align=center] at (0.5, -2.5){
        \begin{tabular}{ll}
            \textbf{LÉGENDE} & \\
            \tikz\filldraw[pattern=north east lines, pattern color=poporange!80!popdark, fill=poporange!15, draw=poporange!80!popdark, thick] (0,0) rectangle (0.3,0.2); & ZES (1980, 1988, 1990) \\
            \tikz\draw[poppurple, thick, fill=poppurple!50] (0,0) circle (0.15); & Métropoles industrielles / portuaires (taille $\propto$ poids) \\
            \tikz\draw[->, poporange, line width=2pt, >=Stealth] (0,0) -- (0.5,0); & Flux maritimes d'exportation / importation \\
        \end{tabular}
    };
\end{tikzpicture}
\legende{Carte corrigée : La façade maritime chinoise, moteur de l'émergence. Les ZES (hachurées orange) concentrent l'investissement initial. Les métropoles (cercles violets) structurent l'urbanisation côtière. Les flèches orange illustrent l'intégration aux chaînes de valeur mondiales.}
\end{popfigure}

\textbf{Justification des choix cartographiques :}
\begin{itemize}
    \item \textbf{ZES :} Figuré surfacique hachuré (orange) pour marquer la \textbf{zone} de droit dérogatoire. Regroupement des 4 ZES 1980 (Delta Fleuve des Perles + Xiamen) + extensions majeures (Hainan, Pudong).
    \item \textbf{Métropoles :} Figuré ponctuel proportionnel (violet). Shanghai (1er port conteneurs mondial) et Guangzhou (cœur Delta Fleuve des Perles) > Tianjin (port Pékin) > Pékin (pouvoir politique/finance) > Wuhan/Chongqing (pivots intérieurs).
    \item \textbf{Flux :} Flèches linéaires orange (couleur \og émergence/énergie \fg) vers les 4 bassins mondiaux. Épaisseur plus forte vers l'Ouest (Europe/USA/Afrique) et le Sud (Asie/ASEAN).
\end{itemize}

\textbf{Étape 2 : Diagramme en bâtons de l'évolution du PIB (Consigne 2)}
\begin{popfigure}
\begin{tikzpicture}
\begin{axis}[
    width=\linewidth,
    height=10cm,
    ybar,
    bar width=12pt,
    ymajorgrids=true,
    ylabel={PIB (Milliards \$ US courants)},
    xlabel={Années},
    xtick=data,
    xticklabels={1978, 1980, 1990, 2000, 2001, 2005, 2010, 2015, 2020},
    x tick label style={rotate=45, anchor=east, font=\footnotesize},
    yticklabel style={font=\footnotesize},
    ymin=0, ymax=16000,
    title={Évolution du PIB nominal de la Chine (1978--2020)},
    title style={font=\normalsize\bfseries, yshift=0.3cm},
    legend style={at={(0.02,0.98)}, anchor=north west, font=\footnotesize, fill=popcream, draw=popline},
    nodes near coords={\pgfmathprintnumber[fixed,precision=0]{\pgfplotspointmeta}},
    nodes near coords style={font=\tiny, rotate=90, anchor=west, yshift=2pt},
    enlarge x limits=0.15,
]
% Data
\addplot[fill=popblue!30, draw=popblue!80!popdark] coordinates {
    (1,150) (2,191) (3,361) (4,1211) (5,1339) (6,2286) (7,6101) (8,11062) (9,14723)
};
% Highlight phases
\addplot[fill=poporange!50, draw=poporange!80!popdark, forget plot] coordinates {(1,150) (2,191) (3,361)};
\addplot[fill=popgreen!50, draw=popgreen!80!popdark, forget plot] coordinates {(4,1211) (5,1339) (6,2286)};
\addplot[fill=poppurple!50, draw=poppurple!80!popdark, forget plot] coordinates {(7,6101) (8,11062) (9,14723)};
\legend{Période 1978-1990 (Amorce), Période 1990-2001 (Réformes), Période 2001-2020 (OMC / Boom)}
\end{axis}
\end{tikzpicture}
\legende{Diagramme en bâtons de l'évolution du PIB chinois. Rupture nette post-1978, accélération exponentielle après l'entrée à l'OMC (2001). Source : Banque Mondiale / FMI.}
\end{popfigure}

\textbf{Commentaire statistique :}
\begin{itemize}
    \item \textbf{1978--1990 :} Croissance lente mais structurelle (150 $\rightarrow$ 361 Mds \$). Mise en place des institutions de marché, agriculture, ZES naissantes.
    \item \textbf{1990--2001 :} Décollage industriel (361 $\rightarrow$ 1 339 Mds \$). \textbf{Taux de variation moyen annuel $\approx$ 12,5\%}. Tournée du Sud (1992), investissements massifs, urbanisation.
    \item \textbf{2001--2020 :} Explosion \og Atelier du monde \fg (1 339 $\rightarrow$ 14 723 Mds \$). \textbf{Multiplication par 11 en 19 ans}. Entrée OMC = levier exportateur maximal. La Chine devient 2\textsuperscript{e} économie mondiale (2010) et 1\textsuperscript{re} en PPA (2014).
\end{itemize}

\textbf{Étape 3 : Carte de la présence chinoise au Cameroun (Consigne 3)}
\begin{popfigure}
\begin{tikzpicture}[scale=0.7, every node/.style={font=\footnotesize}]
    % Cameroon Outline (Very simplified polygon)
    \draw[popgreen!70!popdark, line width=1.5pt, fill=popgreen!5!popcream] 
        (0,0) -- (1,0.5) -- (1.5,1.5) -- (1,2.5) -- (0.5,3.5) -- (0,4) -- (-0.5,3.5) -- (-1,2.5) -- (-1.5,1.5) -- (-1,0.5) -- cycle;
    % North Arrow
    \draw[->, popdark, line width=1pt] (1.5, 3.5) -- ++(0, 0.5) node[above, popdark] {N};
    % Rivers
    \draw[popblue, line width=0.5pt] (0,0) -- (0,1.5) -- (-0.3, 2.5) node[left, popdark] {Sanaga};
    \draw[popblue, line width=0.5pt] (-0.5,0) -- (-0.8, 1.5) node[left, popdark] {Nyong};
    \draw[popblue, line width=0.5pt] (0.5,0) -- (0.8, 1.5) node[right, popdark] {Wouri};
    % Cities
    \node[draw=popdark, fill=popcream, circle, inner sep=1.5pt] (ya) at (0, 2.2) {}; \node[right] at (0.4, 2.2) {Yaoundé};
    \node[draw=popdark, fill=popcream, circle, inner sep=1.5pt] (dl) at (0.8, 0.8) {}; \node[right] at (1.2, 0.8) {Douala};
    \node[draw=popdark, fill=popcream, circle, inner sep=1.5pt] (kr) at (-0.2, 0.3) {}; \node[left] at (-0.6, 0.3) {Kribi};
    \node[draw=popdark, fill=popcream, circle, inner sep=1.5pt] (mv) at (-0.8, 2.0) {}; \node[left] at (-1.3, 2.0) {Memve'ele};
    % Projects (Figurés ponctuels distincts)
    % Port Kribi : Circle with P
    \node[draw=popblue!80!popdark, thick, fill=popblue!20, circle, minimum size=8pt] at (-0.2, -0.3) {P}; \node[below, popblue!80!popdark] at (-0.2, -0.6) {Port Kribi};
    % Barrage Memve'ele : Diamond
    \node[draw=poporange!80!popdark, thick, fill=poporange!20, diamond, minimum size=8pt] at (-0.8, 2.6) {}; \node[above, poporange!80!popdark] at (-0.8, 2.9) {Barrage Memve'ele};
    % Yaoundé Urban/Diplo : Rectangle
    \node[draw=poppurple!80!popdark, thick, fill=poppurple!20, rectangle, minimum size=8pt] at (0.5, 2.8) {}; \node[above, poppurple!80!popdark] at (0.5, 3.1) {Palais Congrès / Échangeurs};
    % Flux Commercial Douala/Kribi <-> Chine
    \draw[<->, popgold!70!popdark, line width=2pt, >=Stealth, bend left=20] (dl) to node[midway, above, sloped, popdark, font=\tiny] {Bois, Coton, Pétrole / Machines, Textiles} ++(2.5, -1.5) node[right, popgold!70!popdark] {CHINE};
    % Legend
    \node[draw=popline, fill=popcream, rounded corners, inner sep=5pt, anchor=north east, align=center] at (2.5, 4){
        \begin{tabular}{ll}
            \textbf{LÉGENDE} & \\
            \tikz\draw[popblue!80!popdark, thick, fill=popblue!20] (0,0) circle (0.1); & Port en eau profonde (Kribi) \\
            \tikz\node[draw=poporange!80!popdark, thick, fill=poporange!20, diamond, minimum size=0.2cm] at (0,0) {}; & Barrage hydroélectrique (Memve'ele) \\
            \tikz\draw[poppurple!80!popdark, thick, fill=poppurple!20] (0,0) rectangle (0.15,0.15); & Infrastructures urbaines / diplomatiques (Yaoundé) \\
            \tikz\draw[<->, popgold!70!popdark, line width=1.5pt, >=Stealth] (0,0) -- (0.5,0); & Flux commerciaux sino-camerounais \\
        \end{tabular}
    };
\end{tikzpicture}
\legende{Carte corrigée : Projets chinois structurants au Cameroun. Kribi (port export), Memve'ele (énergie), Yaoundé (administration/urbanisme). Flux bidirectionnel déséquilibré (matières premières vs manufacturés).}
\end{popfigure}

\textbf{Étape 4 : Rédaction de la synthèse argumentée (Consigne 4)}
\textbf{Proposition de réponse (environ 12 lignes) :}

\og L'émergence chinoise résulte d'un double tournant stratégique initié par Deng Xiaoping en 1978 : la \textbf{réforme interne} (décollectivisation, économie de marché socialiste) et l'\textbf{ouverture externe} via les Zones Économiques Spéciales (Shenzhen, etc.) qui attirent les IDE et les technologies. L'adhésion à l'OMC en 2001 parachève cette insertion : le PIB passe de 150 Mds \$ (1978) à 14 723 Mds \$ (2020), faisant de la Chine l'\og \emph{atelier du monde} \fg{} et la 2\textsuperscript{e} puissance économique. Cette puissance se projette en Afrique par la stratégie \og \emph{Zou Chu Qu} \fg{} (Aller vers l'extérieur). Au Cameroun, elle se traduit par le financement et la réalisation d'infrastructures structurantes par des entreprises d'État (CHEC, Sinohydro) via des prêts concessionnels (Exim Bank) : le port de Kribi (désenclavement/export), le barrage de Memve'ele (souveraineté énergétique), les équipements urbains de Yaoundé. Toutefois, cette présence crée une \textbf{asymétrie} : balance commerciale déficitaire pour le Cameroun (export de matières premières brutes vs import de manufacturés), endettement croissant envers la Chine, et transfert de technologie encore limité. La Chine est ainsi devenue un partenaire \textbf{incontournable mais déséquilibré} pour le développement camerounais. \fg

\textbf{Grille d'évaluation de la rédaction (pour l'enseignant) :}
\begin{center}
\begin{tabularx}{\linewidth}{|l|X|c|}\hline
\textbf{Critère} & \textbf{Indicateurs de réussite} & \textbf{Pts} \\ \hline
Problématique traitée & Réponse explicite aux deux volets (Devenue puissance / Change pour Cameroun) & 2 \\ \hline
Mobilisation docs & Références précises : 1978, ZES, OMC, PIB, Kribi, Memve'ele, Yaoundé & 3 \\ \hline
Articulation Histoire/Géo & Lien causal : Réformes $\rightarrow$ Croissance $\rightarrow$ Projection africaine (Investissements) & 3 \\ \hline
Analyse critique & Mention des limites (asymétrie, dette, matières premières) & 2 \\ \hline
Qualité langue/structure & 10-15 lignes, connecteurs logiques, vocabulaire précis (IDE, OMC, ZES, Exim Bank) & 2 \\ \hline
\textbf{Total} & & \textbf{12} \\ \hline
\end{tabularx}
\end{center}
\end{exoresolu}

\courssub{Bilan de l'activité}
Cette activité d'intégration a permis de valider les acquis suivants, conformément aux attendus du programme de Terminale Technique et aux exigences du Baccalauréat technique :
\begin{enumerate}
    \item \textbf{Maîtrise chronologique et causale :} Les élèves ont remis en ordre les grandes étapes (1978, 1992, 2001) et expliqué le mécanisme \textbf{Réformes $\rightarrow$ ZES/IDE $\rightarrow$ OMC $\rightarrow$ Croissance exponentielle $\rightarrow$ Projection mondiale}.
    \item \textbf{Compétence cartographique :} Ils ont différencié les figurés (surfacique pour les zones/ZES, ponctuel proportionnel pour les métropoles, linéaire pour les flux) et construit des légendes hiérarchisées et normées sur deux espaces distincts (Chine, Cameroun).
    \item \textbf{Compétence statistique :} Le passage du tableau au diagramme en bâtons a nécessité le choix d'une échelle adaptée (gestion de l'hétéroscédasticité forte des données : 150 à 14 723) et la lecture graphique des ruptures (1978, 2001).
    \item \textbf{Ancrage territorial camerounais :} La localisation de Kribi (Sud, façade atlantique), Memve'ele (Sud, fleuve Ntem) et Yaoundé (Centre, cœur politique) a spatialisé la coopération. L'analyse du flux commercial a révélé la structure \og matières premières / manufacturés \fg{} typique de la relation Nord-Sud / Chine-Afrique.
    \item \textbf{Argumentation écrite :} La contrainte de longueur (10-15 lignes) a imposé la concision, la sélection des faits saillants et l'usage du vocabulaire conceptuel (\cle{ZES}, \cle{IDE}, \cle{OMC}, \cle{Exim Bank}, \cle{asymétrie}, \cle{balance commerciale}).
\end{enumerate}
\begin{attention}[Pièges fréquents relevés lors de la correction collective]
\begin{itemize}
    \item \textbf{Carte Chine :} Confondre ZES (zones délimitées) et villes-port (points). Oublier la flèche vers l'Afrique (essentielle pour le lien avec la consigne 3).
    \item \textbf{Graphique :} Utiliser une échelle linéaire inadaptée (bâtons 1978 invisibles) ou une échelle logarithmique (non au programme). Oublier l'unité (Mds \$ US) ou la source.
    \item \textbf{Carte Cameroun :} Placer Kribi au Nord ou Memve'ele au Nord. Utiliser le même figuré pour le port et le barrage.
    \item \textbf{Rédaction :} Faire l'histoire de la Chine \textbf{sans} conclure sur le Cameroun (hors sujet partiel). Inverser le sens du flux (dire que la Chine importe du bois du Cameroun via le port de Kribi \textbf{vers} la Chine, ce qui est juste, mais ne pas mentionner l'import de machines \textbf{depuis} la Chine). Ne pas citer de chiffres (PIB, puissance barrage, capacité port).
\end{itemize}
\end{attention}

\begin{savaistu}[Prolongement possible : La « Nouvelle Route de la Soie » (BRI) et le Cameroun]
Depuis 2013, l'Initiative \og \emph{Belt and Road Initiative} \fg{} (BRI) structure l'action chinoise. Le Cameroun a signé un Mémorandum d'entente BRI en 2017. Le port de Kribi est présenté comme un \og \emph{hub} \fg{} logistique pour l'Afrique centrale (projet de chemin de fer Kribi-Ngaoundéré, liaison avec le Tchad et la RCA). Cela transforme la relation bilatérale en insertion sous-régionale, mais renforce la dépendance aux standards techniques et financiers chinois (contrats EPC \emph{Engineering, Procurement, Construction}).
\end{savaistu}

\newpage

\coursec{Méthodes et exercices résolus}

\courssub{Méthodes et exercices résolus}

\begin{methode}[Analyser un document historique (texte, image, statistique) sur la Chine contemporaine]
\textbf{Objectif :} Appliquer la notion de « source historique » pour comprendre les ruptures (Révolution culturelle, Réformes) et l'émergence économique.
\begin{enumerate}
    \item \textbf{Identifier la nature et l'origine :} Type (discours, affiche, rapport FMI, photo), Auteur (Mao, Deng, journaliste, statisticien), Date, Contexte précis (ex: 9e Congrès PCC, 3e Plénum 1978, adhésion OMC 2001).
    \item \textbf{Dégager le sens général :} Quel est le message principal ? (Mobilisation révolutionnaire, « Quatre modernisations », « Enrichissez-vous », ouverture au marché).
    \item \textbf{Analyser le contenu par thèmes :} Politique (pouvoir, idéologie), Économie (planification, marché, IDE, exportations), Société (Gardes Rouges, migration rurale-urbaine, inégalités), International (rupture URSS, relations US, Afrique).
    \item \textbf{Évaluer la portée et les limites :} Intention de l'auteur (propagande, réforme, alerte), subjectivité, ce que le document \emph{ne dit pas} (répression Tian'anmen 1989, pollution, dette africaine).
    \item \textbf{Conclure en reliant au cours :} En quoi ce document illustre-t-il la transition « Révolution culturelle $\rightarrow$ Émergence économique » ?
\end{enumerate}
\cle{Réflexe Bac :} Toujours citer le document (« L'extrait montre que... », « L'affiche représente... ») et utiliser le vocabulaire du programme (Gang des Quatre, Responsabilité familiale, ZES, OMC, Nouvelles Routes de la Soie).
\end{methode}

\begin{exoresolu}[Commentaire d'un extrait du discours de Deng Xiaoping, « Émanciper l'esprit, chercher la vérité dans les faits » (décembre 1978)]
\textbf{Énoncé :} Vous commenterez le document suivant en mobilisant vos connaissances.\\
\textit{« Nous devons intégrer la vérité universelle du marxisme aux réalités concrètes de la Chine... La pratique est le critère unique de la vérité. Si nous ne faisons pas croître les forces productives, nous ne pouvons pas réaliser le socialisme. La pauvreté n'est pas du socialisme. »} — Deng Xiaoping, séance de clôture de la session de travail du Comité central du PCC, 13 décembre 1978.

\tcblower
\textbf{Solution détaillée}

\textbf{1. Présentation / Identification}
\begin{itemize}
    \item \textbf{Nature :} Extrait de discours politique (transcription officielle).
    \item \textbf{Auteur :} \cle{Deng Xiaoping} (1904-1997), dirigeant \emph{de facto} de la Chine après l'arrestation du Gang des Quatre (octobre 1976) et la chute de Hua Guofeng. Architecte des réformes.
    \item \textbf{Date :} 13 décembre 1978 (veille du \cle{3e Plénum du 11e Comité Central}, 18-22 déc. 1978, acte de naissance officiel de la politique de « Réforme et Ouverture »).
    \item \textbf{Contexte :} Sortie de la \cle{Révolution culturelle} (1966-1976) : chaos politique, économie exsangue, légitimité du PCC ébranlée. Nécessité de rompre avec le dogmatisme maoïste (lutte des classes comme lien principal) sans abandonner le monopole du pouvoir du Parti.
\end{itemize}

\textbf{2. Analyse du contenu (Thèmes et arguments)}
\begin{itemize}
    \item \textbf{Rupture épistémologique :} « La pratique est le critère unique de la vérité ». Rejet du dogmatisme (suivre aveuglément Mao ou Marx-Lénine). Légitime le pragmatisme : une politique est bonne si elle marche (croissance), peu importe son étiquette « capitaliste » ou « socialiste ».
    \item \textbf{Redéfinition du socialisme :} « La pauvreté n'est pas du socialisme ». Le but n'est plus l'égalitarisme révolutionnaire mais le développement des \cle{forces productives} (économie). Le socialisme = richesse collective via croissance.
    \item \textbf{Adaptation du marxisme :} « Intégrer la vérité universelle... aux réalités concrètes de la Chine ». Justification théorique du « socialisme à la chinoise » : économie de marché socialiste, ouverture au capital étranger, agriculture familiale (système de responsabilité).
\end{itemize}

\textbf{3. Portée et Limites (Esprit critique)}
\begin{itemize}
    \item \textbf{Portée :} Texte fondateur. Lance les \cle{Quatre Modernisations} (Agriculture, Industrie, Défense, Science/Technique). Début de la décollectivisation (contrats familiaux), création des \cle{ZES} (Shenzhen 1980), attraction des IDE.
    \item \textbf{Limites / Non-dits :} Silence total sur la \cle{démocratie} et les droits de l'homme. « Émanciper l'esprit » s'arrête aux portes du pouvoir politique (répression de 1989 préfigurée). Ne mentionne pas les coûts sociaux futurs (inégalités, corruption, environnement).
\end{itemize}

\textbf{4. Conclusion}
Ce document marque le \cle{tournant de 1978} : la Chine choisit le \cle{pragmatisme économique} sous contrôle politique autoritaire. Il explique la trajectoire unique d'une « émergence économique » sans transition démocratique, modèle qui structure encore aujourd'hui les relations Chine-Afrique (infrastructures contre ressources, non-ingérence politique).

\textbf{Phrase de conclusion type :} \emph{« En définitive, ce discours de Deng Xiaoping acte le passage du "socialisme pauvre" maoïste au "socialisme de marché" productiviste, fondant la légitimité du PCC sur la performance économique plutôt que sur la pureté idéologique. »}
\end{exoresolu}

\begin{methode}[Construire et exploiter une frise chronologique commentée / Repérer ruptures et continuités]
\textbf{Objectif :} Maîtriser la chronologie 1966-2010 pour structurer une dissertation ou un commentaire.
\begin{enumerate}
    \item \textbf{Fixer les bornes :} 1966 (Début Révolution culturelle) — 1976 (Mort Mao / Arrestation Gang des Quatre) — 1978 (3e Plénum / Deng) — 1989 (Tian'anmen) — 2001 (Entrée OMC) — 2010 (2e économie mondiale).
    \item \textbf{Identifier 3 grandes périodes :} 
    \begin{itemize}
        \item \textbf{1966-1976 :} Chaos révolutionnaire (Rupture politique/sociale, régression économique).
        \item \textbf{1978-1989 :} Réformes graduelles, ouverture contrôlée, début croissance (Continuité politique : parti unique ; Rupture économique : marché).
        \item \textbf{1992-2010 :} Accélération (Voyage Sud Deng 1992), « Atelier du monde », entrée OMC, puissance mondiale (Rupture d'échelle économique).
    \end{itemize}
    \item \textbf{Repérer les continuités :} Monopole PCC, primauté de la stabilité, nationalisme, planification indicative (Plans quinquennaux maintenus).
    \item \textbf{Utiliser des codes visuels :} Flèches montantes (PIB, Exportations), flèches descendantes (Taux de pauvreté), symboles (Étoile rouge $\rightarrow$ Engrenage/Usine $\rightarrow$ Globe).
\end{enumerate}
\cle{Formule à retenir :} \textbf{Taux de variation annuel moyen} = $\left( \frac{V_{\text{final}}}{V_{\text{initial}}} \right)^{\frac{1}{n}} - 1$ (pour croissances composées).
\end{methode}

\begin{exoresolu}[Frise chronologique et calcul de croissance : La Chine 1978-2010]
\textbf{Énoncé :} 
\begin{enumerate}
    \item Classer les événements suivants par ordre chronologique : Mort de Mao ; Voyage du Sud de Deng ; Création de la ZES de Shenzhen ; Entrée à l'OMC ; Répression de Tian'anmen ; 3e Plénum du 11e CC ; Début Révolution culturelle.
    \item Le PIB chinois passe d'environ 150 milliards \$ (1978) à 6 000 milliards \$ (2010). Calculer le taux de croissance annuel moyen sur la période (arrondi à 0,1\%).
    \item Expliquer en deux lignes comment la chronologie éclaire ce chiffre.
\end{enumerate}

\tcblower
\textbf{Solution détaillée}

\textbf{1. Ordre chronologique}
1. 1966 : Début Révolution culturelle \\
2. 1976 : Mort de Mao (septembre) / Arrestation Gang des Quatre (octobre) \\
3. 1978 (déc.) : \cle{3e Plénum du 11e Comité Central} (Lancement réformes) \\
4. 1980 : \cle{Création de la ZES de Shenzhen} (Expérimentation marché/IDE) \\
5. 1989 (juin) : \cle{Répression de Tian'anmen} (Gel politique, poursuite économique) \\
6. 1992 (janv.-fév.) : \cle{Voyage du Sud de Deng} (Relance réformes, « Enrichissez-vous ») \\
7. 2001 (déc.) : \cle{Entrée à l'OMC} (Intégration mondiale définitive)

\textbf{2. Calcul du taux de croissance annuel moyen (TCAM)}
\begin{itemize}
    \item $V_{\text{initial}} = 150$ (milliards \$)
    \item $V_{\text{final}} = 6000$ (milliards \$)
    \item $n = 2010 - 1978 = 32$ années
    \item Formule : $TCAM = \left( \frac{6000}{150} \right)^{\frac{1}{32}} - 1$
    \item $\frac{6000}{150} = 40$ (Le PIB a été multiplié par 40)
    \item $40^{\frac{1}{32}} = e^{\frac{\ln(40)}{32}} \approx e^{\frac{\num{3,6889}}{32}} \approx e^{\num{0,11528}} \approx \num{1,1222}$
    \item $TCAM \approx \num{1,1222} - 1 = \num{0,1222} = \textbf{\num{12,2}\% par an}$.
\end{itemize}
\textbf{Résultat :} \textbf{\num{12,2}\%} de croissance moyenne annuelle sur 32 ans (performance historique inégalée pour une grande économie).

\textbf{3. Lien chronologie / chiffre}
Ce chiffre spectaculaire est le produit direct de la \cle{chronologie des réformes} : la rupture de 1978 (fin planification rigide) + l'accélération de 1992 (libéralisation prix, entreprises privées) + l'intégration 2001 (accès marchés mondiaux) ont créé un cercle vertueux investissement-exportation-croissance ininterrompu, malgré le choc politique de 1989.

\textbf{Phrase de conclusion :} \emph{« La frise chronologique révèle que la croissance à deux chiffres n'est pas un hasard mais le résultat d'une séquence réformiste continue (1978 $\rightarrow$ 1992 $\rightarrow$ 2001) où chaque étape a levé un frein structurel majeur. »}
\end{exoresolu}

\begin{methode}[Rédiger une dissertation historique structurée et argumentée (Sujet de type Bac)]
\textbf{Objectif :} Répondre à une problématique en 3 parties équilibrées (Thèse / Antithèse / Synthèse ou Chronologique / Thématique).
\begin{enumerate}
    \item \textbf{Analyser le sujet :} Mots-clés, bornes chronologiques, concepts (émergence, modèle, contradictions). Formuler la \cle{problématique} (question à laquelle le plan répond).
    \item \textbf{Rédiger l'introduction (ENTONNOIR) :} 
    \begin{itemize}
        \item Accroche (chiffre choc, citation, fait marquant).
        \item Définition des termes du sujet.
        \item Cadrage spatio-temporel.
        \item \textbf{Problématique} (en gras).
        \item \textbf{Annonce du plan} (3 parties, 2-3 sous-parties chacune).
    \end{itemize}
    \item \textbf{Développer en 3 parties (A, B, C) :} Chaque partie = 1 idée force. 
    \begin{itemize}
        \item Titre de partie (thématique ou chronologique).
        \item Sous-parties (arguments) : \textbf{Argument} + \textbf{Exemple précis daté/chiffré} + \textbf{Explication/Lien avec problématique}.
        \item Mini-transition vers partie suivante.
    \end{itemize}
    \item \textbf{Rédiger la conclusion :} 
    \begin{itemize}
        \item Réponse synthétique à la problématique.
        \item Ouverture (actualité : Xi Jinping, Nouvelles Routes de la Soie, rivalité US/Chine, dette Cameroun).
    \end{itemize}
    \item \textbf{Relire :} Vérifier l'orthographe des noms propres (Deng Xiaoping, Zhu Rongji, Xi Jinping), les dates, la cohérence du fil conducteur.
\end{enumerate}
\cle{Reflexe Bac :} Pas de « Je pense que ». Utiliser « L'analyse montre que », « L'exemple de... illustre », « Toutefois », « En revanche ».
\end{methode}

\begin{exoresolu}[Dissertation : « La Chine depuis 1978 : une émergence économique sans modèle politique ? »]
\textbf{Énoncé :} Sujet de dissertation. Traiter le sujet en suivant la méthode.

\tcblower
\textbf{Solution détaillée (Plan détaillé et éléments de rédaction)}

\textbf{INTRODUCTION}
\textit{Accroche :} En 2010, la Chine devient la 2e économie mondiale, multipliant son PIB par 40 en 32 ans. 
\textit{Définitions :} Émergence = transition vers économie de marché intégrée, croissance forte, urbanisation. Modèle politique = Parti-État unique, contrôle idéologique, absence d'élections libres.
\textit{Cadrage :} Chine, 1978 (3e Plénum) à nos jours.
\textbf{Problématique :} \textbf{Comment la Chine a-t-elle concilié une libéralisation économique radicale avec le maintien d'un autoritarisme politique inchangé, et ce modèle « hybride » est-il durable ?}
\textbf{Plan :} 
I. 1978-1989 : La primauté de l'économique sur le politique (Réformes sans démocratie). 
II. 1992-2012 : L'institutionnalisation du « capitalisme d'État » (Croissance comme légitimité). 
III. Depuis 2012 (Xi Jinping) : Les tensions du modèle (Contrôle renforcé, nouvelles contradictions).

\textbf{DÉVELOPPEMENT}

\textbf{I. 1978-1989 : Réformes économiques graduelles, verrouillage politique}
\begin{itemize}
    \item \textbf{A. Rupture économique :} Désengagement de l'État agricole (Système responsabilité familiale $\rightarrow$ +50\% prod. 1978-1984), ZES (Shenzhen), IDE. « La pratique critère de vérité ».
    \item \textbf{B. Continuité politique :} Refus du multipartisme (Deng : « Quatre principes cardinaux »). Répression mouvements étudiants 1986-87 (Hu Yaobang limogé).
    \item \textbf{C. Rupture de 1989 :} Tian'anmen. Démonstration que l'enrichissement ne mènera pas à la libéralisation politique. Sanctions occidentales temporaires.
\end{itemize}

\textbf{II. 1992-2012 : L'âge d'or du « capitalisme d'État » (Modèle Pékin)}
\begin{itemize}
    \item \textbf{A. Relance par le marché (Voyage Sud 1992) :} « Enrichissez-vous ». Entreprises privées légalisées (1993), réforme entreprises d'État (Zhu Rongji, fin 90s : licenciements massifs, fermetures usines déficitaires).
    \item \textbf{B. Intégration mondiale (OMC 2001) :} « Atelier du monde ». Excédents commerciaux massifs, réserves de change (3000 Mrd \$ 2011). Urbanisation forcée (migrants ruraux).
    \item \textbf{C. Contrat social implicite :} « On vous laisse faire de l'argent, vous ne touchez pas au pouvoir ». Stabilité = Légitimité. Contrôle internet (Grand Firewall).
\end{itemize}

\textbf{III. Depuis 2012 (Xi Jinping) : Les limites et contradictions de l'émergence}
\begin{itemize}
    \item \textbf{A. Renforcement autoritaire :} Lutte anti-corruption (élimination rivaux), suppression limite mandats (2018), contrôle total société (crédit social, Xinjiang, Hong Kong 2020).
    \item \textbf{B. Défis économiques structurels :} Vieillissement (politique enfant unique), bulle immobilière (Evergrande), endettement local, piège du revenu moyen, guerre techno US (puces).
    \item \textbf{C. Modèle à l'exportation ?} Nouvelles Routes de la Soie (BRI) : projection du modèle (infrastructures + influence politique) en Afrique/Cameroun. Dette, dépendance, non-transparence.
\end{itemize}

\textbf{CONCLUSION}
\textit{Bilan :} L'émergence chinoise \textbf{repose sur un modèle politique unique} : État fort, planification stratégique, marché encadré, nationalisme. Succès historique (sortie pauvreté 800M personnes).
\textit{Ouverture :} Ce modèle entre en zone de turbulence (ralentissement croissance \textasciitilde \num{5}\%, tensions Taïwan, découplage US). La question n'est plus « émergence » mais « transformation » : la Chine peut-elle innover sans liberté ? Le Cameroun doit-il copier le modèle ou négocier son autonomie ?

\textbf{Phrase de conclusion type :} \emph{« L'émergence chinoise valide l'hypothèse d'un développement autoritaire efficace à court terme, mais expose désormais ses limites structurelles : l'absence de correction démocratique rend la gestion des crises (démographique, financière, géopolitique) plus risquée. »}
\end{exoresolu}

\begin{methode}[Exploiter des données statistiques (tableaux, graphiques) pour mesurer l'émergence]
\textbf{Objectif :} Transformer des chiffres bruts en arguments historiques (parts de marché, taux de croissance, structure du PIB).
\begin{enumerate}
    \item \textbf{Lire les métadonnées :} Source (Banque Mondiale, FMI, Douanes chinoises, NBS), Unités (Mrd \$, \%, milliards \$), Années, Périmètre (PPI vs PPA).
    \item \textbf{Calculer les indicateurs clés :} 
    \begin{itemize}
        \item Taux de croissance : $\frac{N - N-1}{N-1} \times 100$.
        \item Part dans le total : $\frac{Valeur_{\text{Chine}}}{Valeur_{\text{Monde}}} \times 100$.
        \item Structure : $\frac{Secteur}{PIB} \times 100$ (Primaire / Secondaire / Tertiaire).
    \end{itemize}
    \item \textbf{Comparer dans l'espace et le temps :} Chine vs USA / UE / Afrique ; 1978 vs 2001 vs 2020.
    \item \textbf{Interpréter :} Une hausse de la part des services = tertiarisation / maturité. Baisse part primaire = modernisation. Excédent commercial = modèle export-led.
    \item \textbf{Formuler une phrase-argument :} « Le graphique 2 montre que la part de l'industrie dans le PIB passe de 40\% (1990) à 30\% (2020), signe d'une \cle{tertiairisation} tardive mais rapide, caractéristique du passage au stade d'économie émergente avancée. »
\end{enumerate}
\cle{Attention :} Distinguer PIB nominal (\$) et PIB PPA (Parité Pouvoir Achat). La Chine est 1ère en PPA depuis 2014, 2e en nominal.
\end{methode}

\begin{exoresolu}[Analyse statistique : Structure du commerce extérieur chinois (2000 vs 2020)]
\textbf{Énoncé :} À partir du tableau ci-dessous, répondre aux questions.
\begin{center}
\begin{tabularx}{\linewidth}{|l|X|X|X|X|}\hline
\textbf{Indicateur} & \textbf{2000} & \textbf{2010} & \textbf{2020} & \textbf{Monde 2020} \\ \hline
Exportations (Mrd \$) & 249 & 1 578 & 2 591 & 17 500 \\ \hline
Importations (Mrd \$) & 225 & 1 396 & 2 056 & 17 300 \\ \hline
Solde commercial (Mrd \$) & +24 & +182 & +535 & — \\ \hline
Part exportations haute technologie (\% total) & 21\% & 30\% & 26\% & 18\% \\ \hline
\end{tabularx}
\end{center}
\begin{enumerate}
    \item Calculer le taux de croissance des exportations entre 2000 et 2010 (annuel moyen) et entre 2010 et 2020 (annuel moyen). Comparer.
    \item Calculer la part de la Chine dans les exportations mondiales en 2020.
    \item Interpréter l'évolution de la part des haute technologies (2000-2020) au regard de la stratégie « Made in China 2025 ».
\end{enumerate}

\tcblower
\textbf{Solution détaillée}

\textbf{1. Taux de croissance annuel moyen (TCAM) des exportations}
\textbf{Période 2000-2010 (10 ans) :}
$V_i = 249$, $V_f = 1578$.
Ratio $= 1578 / 249 = \num{6,337}$.
$TCAM = \num{6,337}^{1/10} - 1 = e^{\ln(\num{6,337})/10} - 1 = e^{\num{1,846}/10} - 1 = e^{\num{0,1846}} - 1 \approx \num{1,2027} - 1 = \textbf{\num{20,3}\% par an}$.
\textit{Interprétation :} Croissance explosive post-OMC (2001), phase « atelier du monde » bas de gamme.

\textbf{Période 2010-2020 (10 ans) :}
$V_i = 1578$, $V_f = 2591$.
Ratio $= 2591 / 1578 = \num{1,642}$.
$TCAM = \num{1,642}^{1/10} - 1 = e^{\ln(\num{1,642})/10} - 1 = e^{\num{0,496}/10} - 1 = e^{\num{0,0496}} - 1 \approx \num{1,0508} - 1 = \textbf{\num{5,1}\% par an}$.
\textit{Interprétation :} Net ralentissement (crise 2008, hausse salaires, guerre commerciale US, Covid). Transition vers croissance plus qualitative.

\textbf{Comparaison :} Le rythme a été \textbf{divisé par 4} (\num{20}\% $\rightarrow$ \num{5}\%), marquant la fin de l'hyper-croissance extensive.

\textbf{2. Part de la Chine dans les exportations mondiales (2020)}
$\text{Part} = \frac{2591}{17500} \times 100 = \textbf{\num{14,8}\%}$.
\textit{Conclusion :} La Chine est le \textbf{premier exportateur mondial} (devant USA \textasciitilde \num{10}\%, Allemagne \textasciitilde \num{8}\%).

\textbf{3. Évolution part haute technologie (HT) et lien « Made in China 2025 »}
\begin{itemize}
    \item 2000 : \num{21}\% (Déjà significative : assemblage électronique, Dell/HP/Foxconn).
    \item 2010 : \num{30}\% (Pic : chaîne de valeur mondiale intégrée, mais \emph{assembly} bas de gamme, valeur ajoutée faible).
    \item 2020 : \num{26}\% (Légère baisse).
\end{itemize}
\textbf{Analyse :} La baisse 2010-2020 peut sembler contre-intuitive. Elle s'explique par :
\begin{enumerate}
    \item Hausse des exportations de biens intermédiaires/traditionnels (masques, acier, navires) en 2020 (Covid).
    \item Déplacement de l'assemblage final vers Vietnam/Inde (coûts).
    \item \textbf{Mais :} La stratégie \cle{Made in China 2025} (2015) vise à monter dans la chaîne de valeur (puces, IA, biotech, aéro). La part HT \emph{en valeur ajoutée nationale} augmente (réduction dépendance puces US), même si la part \emph{brute} à l'export stagne. Le chiffre 2020 (\num{26}\%) reste \textbf{supérieur à la moyenne mondiale (\num{18}\%)}, prouvant une spécialisation structurelle.
\end{enumerate}

\textbf{Phrase de conclusion :} \emph{« Les données confirment le basculement d'une croissance extensive (volume, bas de gamme) vers une croissance intensive (technologie, valeur ajoutée), objectif central du modèle chinois post-2010, même si les statistiques brutes d'exportation masquent la montée en gamme réelle de la production domestique. »}
\end{exoresolu}

\begin{methode}[Réaliser et commenter un croquis / une carte mentale sur l'insertion mondiale de la Chine]
\textbf{Objectif :} Visualiser les flux (matières, capitaux, biens) et la stratégie géoéconomique (BRI).
\begin{enumerate}
    \item \textbf{Choisir le fond de carte :} Planisphère simplifié (contours continents) ou Carte centrée Pacifique (vision chinoise).
    \item \textbf{Définir la légende (3 à 4 items max) :} 
    \begin{itemize}
        \item Flux commerciaux (flèches proportionnelles : Exportations Chine $\rightarrow$ Monde ; Importations Matières $\rightarrow$ Chine).
        \item IDE sortants (Stocks : traits pointillés épais vers Afrique, Asie du Sud-Est, Europe, USA).
        \item Projets BRI (Nouvelles Routes de la Soie) : Traits continus (Route terrestre), Traits pointillés (Route maritime), Points (Ports stratégiques : Gwadar, Djibouti, Le Pirée, Kribi).
        \item Zones d'influence / Conflits (Mer de Chine, Taïwan, Frontière Inde) : Hachures ou couleur distincte.
    \end{itemize}
    \item \textbf{Respecter la sémiologie graphique :} 
    \begin{itemize}
        \item Flèches : Épaisseur $\propto$ Volume.
        \item Couleurs : Rouge/Orange (Chine/Flux sortants), Bleu (Flux entrants/Mer), Vert (BRI Terre), Violet (BRI Mer).
        \item Villes/Ports : Carrés/Points proportionnels.
    \end{itemize}
    \item \textbf{Titrer et sourcer :} « La Chine, pivot des flux mondiaux (2023) » — Source : UNCTAD, Ministère Commerce Chine.
    \item \textbf{Commenter (3 phrases) :} 1. Description (Quoi ? Où ?). 2. Explication (Pourquoi ? Stratégie). 3. Enjeu/Limite (Dépendance, Endettement, Rivalité).
\end{enumerate}
\cle{Reflexe Bac :} Même sans dessiner, savoir \emph{décrire} le croquis attendu oralement ou en texte structuré (Légende organisée).
\end{methode}

\begin{exoresolu}[Description et analyse d'un croquis : « La Chine et l'Afrique : le corridor Cameroun »]
\textbf{Énoncé :} Vous devez présenter à l'oral (3 min) un croquis sur « La présence chinoise au Cameroun ». Structurez votre présentation en suivant la méthode. Indiquez les éléments de légende, le titre, et le commentaire type.

\tcblower
\textbf{Solution détaillée (Guide de présentation orale / rédaction)}

\textbf{1. Titre proposé :} \textbf{« Cameroun-Chine : Infrastructures, Ressources et Dette (2000-2024) »}

\textbf{2. Légende structurée (4 items)}
\begin{itemize}
    \item \textbf{Fonds de carte :} Carte du Cameroun (contour), capitales (Yaoundé, Douala), pays limitrophes, Golfe de Guinée.
    \item \textbf{Flux physiques (Flèches Bleues $\rightarrow$ Chine) :} 
    \begin{itemize}
        \item Épaisseur forte : Bois (Est/Sud $\rightarrow$ Douala), Coton (Nord), Pétrole (Kribi/Logbaba), Minerais (Fer Mbalam/Nabeba $\rightarrow$ Kribi - projet).
        \item Légende : « Exportations matières premières (Bois, Coton, Minerai, Pétrole) ».
    \end{itemize}
    \item \textbf{Flux financiers \& Projets (Flèches Rouges $\leftarrow$ Chine / Points Rouges) :}
    \begin{itemize}
        \item Points (Taille $\propto$ Montant) : Barrage Lom Pangar, Barrage Memve'ele, Autoroute Yaoundé-Douala (tronçons), Port en eau profonde de Kribi (CHEC), Stade Olembé, Hôpital de référence Yaoundé, Université de Ngaoundéré.
        \item Traits pleins rouges : Chemin de fer (Transcamerounais réhabilitation), Pipeline Tchad-Cameroun (partenaire chinois).
        \item Légende : « Projets d'infrastructures financés/construits par la Chine (Prêts concessionnels, FOCAC) ».
    \end{itemize}
    \item \textbf{Présence humaine / Institutions (Symboles) :}
    \begin{itemize}
        \item Ambassade (Yaoundé), Consulat (Douala).
        \item Instituts Confucius (Yaoundé, Douala, Maroua) — \emph{Carré violet}.
        \item Zones d'implantation entreprises (Douala, Kribi, Yaoundé) — \emph{Hachures rouges}.
    \end{itemize}
\end{itemize}

\textbf{3. Commentaire type (3 étapes)}

\textbf{Étape 1 : Description (Le « Quoi/Où »)}
\emph{« Le croquis met en évidence une asymétrie structurelle : le Cameroun exporte vers la Chine des matières premières brutes (bois, coton, minerais) depuis l'intérieur vers le port de Douala/Kribi (flèches bleues), tandis que la Chine finance et réalise les infrastructures lourdes (barrages, port, routes, stades) et forme les élites (Instituts Confucius), créant un maillage territorial dense (points/traits rouges). »}

\textbf{Étape 2 : Explication (Le « Pourquoi/Comment » - Cours)}
\emph{« Cette configuration illustre la stratégie chinoise « Ressources contre Infrastructures » (Angola Model) et l'axe « Route de la Soie Maritime » : la Chine sécurise son approvisionnement (bois, minerais critiques pour transition énergétique) et place ses excédents de capacité (BTP, acier) via des prêts d'État (Exim Bank) remboursables en ressources ou sur budget national. Le port de Kribi (CHEC) est la clé de voûte logistique. »}

\textbf{Étape 3 : Enjeux / Limites (L'esprit critique / « Savoir-faire : Justifier »)}
\emph{« Cependant, le croquis masque des tensions : 1) \textbf{Endettement} : Part chinoise \textasciitilde \num{40-50}\% dette extérieure publique (risque surendettement FMI). 2) \textbf{Contenu local faible} : Main-d'œuvre/entreprises chinoises majoritaires, transfert technologique limité. 3) \textbf{Gouvernance/Environnement} : Opacité contrats, déforestation (bois), impact social barrages. 4) \textbf{Dépendance stratégique} : Port de Kribi géré par Chinois sur 50 ans. Le Cameroun doit négocier l'industrialisation locale (transformation bois/coton sur place) pour rééquilibrer. »}

\textbf{Phrase de conclusion orale :} \emph{« En somme, la carte révèle une interdépendance asymétrique : la Chine est devenue le partenaire structurel incontournable du Cameroun, mais le modèle actuel reproduit le schéma extractif colonial, nécessitant une renégociation vers plus de valeur ajoutée locale. »}
\end{exoresolu}

\begin{methode}[Comparer deux périodes ou deux modèles à l'aide d'un tableau comparatif]
\textbf{Objectif :} Mettre en évidence ruptures et continuités (Mao vs Deng ; Chine vs Japon/Corée ; Plan vs Marché).
\begin{enumerate}
    \item \textbf{Définir les critères de comparaison (colonnes) :} Politique (Pouvoir, Idéologie), Économie (Propriété, Allocation ressources, Ouverture), Société (Ville/Campagne, Éducation, Contrôle), International (Alignement, Statut).
    \item \textbf{Remplir les lignes (périodes/modèles) avec des mots-clés précis :} Éviter « C'est mieux ». Préférer « Planification centrale rigide » vs « Économie de marché socialiste / Planification indicative ».
    \item \textbf{Utiliser des symboles visuels :} $\nearrow$ (hausse), $\searrow$ (baisse), $=$ (stabilité), $\neq$ (rupture), $\rightarrow$ (évolution).
    \item \textbf{Rédiger une synthèse (2-3 lignes) :} Identifier la \cle{grande rupture} et la \cle{grande continuité}.
\end{enumerate}
\cle{Grille type Bac :} 4 critères $\times$ 2 périodes = 8 cases à remplir minimum pour 4-5 points.
\end{methode}

\begin{exoresolu}[Tableau comparatif : La Chine de Mao (1949-1976) vs La Chine de Deng (1978-1997)]
\textbf{Énoncé :} Complétez le tableau comparatif ci-dessous pour opposer les deux ères. Chaque case doit contenir un terme précis ou une courte phrase (max 5 mots). Puis rédigez la synthèse.

\begin{center}
\begin{tabularx}{\linewidth}{|l|X|X|}\hline
\textbf{Critère} & \textbf{Ère Mao (1949-1976)} & \textbf{Ère Deng (1978-1997)} \\ \hline
\textbf{Idéologie directrice} & Pensée Mao Zedong (Lutte des classes continue) & Socialisme à la chinoise (Pragmatisme : « Chat noir/chat blanc ») \\ \hline
\textbf{Moteur de l'histoire} & Lutte des classes / Mobilisation politique & Développement forces productives / Croissance économique \\ \hline
\textbf{Organisation agricole} & Collectivisation forcée $\rightarrow$ Communes populaires (Grand Bond) $\rightarrow$ Responsabilité collective & \textbf{Système de responsabilité familiale} (Contrat ménage) $\rightarrow$ Désengagement État \\ \hline
\textbf{Entreprises d'État / Privé} & 100\% Public (Prop. d'État) / Privé interdit & \textbf{Dualisme} : SSE réformées + Entreprises collectives (TVE) + Privé légalisé (1988) + IDE (ZES) \\ \hline
\textbf{Ouverture au monde} & Autarcie / Rupture URSS (1960) / Rapprochement US (1972) diplomatique seulement & \textbf{« Réforme et Ouverture »} : ZES, OMC (préparée), IDE massifs, Technologie importée \\ \hline
\textbf{Place de l'armée / Police} & Armée politique (Lin Biao) / Gardes Rouges (chaos) / Milices & Armée professionnalisée (PLA) / Police armée (PAP) / Stabilité sociale prioritaire \\ \hline
\textbf{Objectif affiché} & Communisme (transition rapide) / Égalitarisme & « Xiaokang » (Aisance modérée) / Modernisation / Puissance nationale \\ \hline
\end{tabularx}
\end{center}

\textbf{Synthèse (Ruptures / Continuités)}

\textbf{Rupture majeure (Économie/Société) :} Passage d'une \textbf{économie administrée, fermée, égalitariste et politisée} (Mao) à une \textbf{économie de marché socialiste, ouverte, inégalitaire mais performante} (Deng). Le paysan devient acteur économique (responsabilité familiale) ; l'entrepreneur privé devient légitime ; le marché alloue les ressources.

\textbf{Continuité majeure (Politique) :} \textbf{Monopole absolu du PCC} sur le pouvoir. Refus du pluralisme, de la séparation des pouvoirs, de la liberté de presse. L'armée reste « l'armée du Parti ». La répression de 1989 (Deng) prouve que la réforme économique ne s'accompagne pas de réforme politique.

\textbf{Conclusion :} \emph{« Deng Xiaoping a « changé la base » (économie) pour « sauver le sommet » (pouvoir du PCC). Le « socialisme à la chinoise » est né de ce paradoxe : libérer les forces productives capitalistes pour renforcer un État-parti léniniste. »}
\end{exoresolu}

\begin{attention}[Les 5 erreurs qui coûtent des points au Bac (Histoire - Chine)]
\begin{enumerate}
    \item \textbf{Confondre « Révolution culturelle » et « Grande Saut en avant » :} La Révolution culturelle (1966-1976) est politico-culturelle (Gardes Rouges, lutte au sommet PCC). Le Grand Bond (1958-1961) est économique (communes, fours à acier, famine). \textbf{Sanction :} Hors-sujet / Incohérence chronologique majeure.
    \item \textbf{Dire « La Chine est devenue capitaliste » ou « La Chine a abandonné le communisme » :} \textbf{Erreur conceptuelle.} Le PCC revendique le « socialisme à la chinoise ». L'État garde le contrôle des secteurs stratégiques (banques, énergie, télécoms, médias) et le monopole politique. Dire « transition capitaliste » = 0 point sur l'analyse du modèle politique.
    \item \textbf{Oublier l'année 1989 (Tian'anmen) dans une dissertation sur 1978-2010 :} C'est le \textbf{tournant politique majeur} qui fige le modèle « Marché OUI / Démocratie NON ». L'ignorer montre une méconnaissance de la continuité autoritaire.
    \item \textbf{Utiliser des chiffres sans source ni année (ex: « La Chine a beaucoup grandi ») :} \textbf{Sanction :} Argument non étayé. Au Bac, tout chiffre (PIB, \% export, population urbaine) doit être \textbf{daté, sourcé (Banque Mondiale, NBS) et comparé}. Ex: « PIB $\times$ 40 entre 1978 et 2010 (BM) ».
    \item \textbf{Traiter les relations Chine-Afrique / Cameroun comme une « aide » ou « amitié » sans nuance :} \textbf{Manque d'esprit critique.} C'est un \textbf{partenariat d'intérêts} (Ressources/Infrastructures/Marchés/Dette). Dire « La Chine aide le Cameroun » = vision naïve. Dire « Néocolonialisme » sans preuves = vision idéologique. \textbf{Attendu :} « Relation asymétrique : infrastructures contre ressources, endettement risqué, faible transfert tech, enjeu géopolitique (Taiwan, ONU) ».
\end{enumerate}
\end{attention}

\newpage

\coursec{Exercices}

\courssub{Je vérifie mes connaissances}

\begin{exercice}[QCM : dates et acteurs]
Pour chaque question, entoure la bonne réponse.
\begin{enumerate}
\item La proclamation de la République populaire de Chine a lieu en : a) 1911 \quad b) 1949 \quad c) 1966 \quad d) 1978.
\item La Révolution culturelle est lancée par : a) Deng Xiaoping \quad b) Sun Yat-sen \quad c) Mao Zedong \quad d) Jiang Qing seule.
\item Les réformes économiques chinoises débutent en : a) 1966 \quad b) 1976 \quad c) 1978 \quad d) 2001.
\item L'entrée de la Chine dans l'Organisation mondiale du commerce (OMC) date de : a) 1978 \quad b) 1989 \quad c) 1997 \quad d) 2001.
\end{enumerate}
\end{exercice}

\begin{exercice}[Vrai ou faux, à justifier]
Réponds par vrai ou faux, puis justifie ta réponse en une ou deux phrases.
\begin{enumerate}
\item La Révolution culturelle (1966-1976) a renforcé l'économie chinoise.
\item Les \cle{Gardes rouges} étaient des soldats professionnels de l'armée chinoise.
\item Deng Xiaoping a lancé les \cle{Quatre Modernisations} (agriculture, industrie, défense, sciences et techniques).
\item Les \cle{ZES} (zones économiques spéciales) se situent principalement dans l'intérieur de la Chine.
\end{enumerate}
\end{exercice}

\begin{exercice}[Définitions]
Définis en deux ou trois lignes chacune des notions suivantes : \cle{Révolution culturelle}, \cle{ZES}, \cle{hukou}, \cle{« atelier du monde »}.
\end{exercice}

\begin{exercice}[Calcul direct : la croissance du PIB chinois]
Le PIB de la Chine (en milliards de dollars courants) est donné ci-dessous :
\begin{center}
\begin{tabularx}{\linewidth}{|l|X|X|X|X|X|}
\hline
\rowcolor{popblueL} Année & 1978 & 1990 & 2000 & 2010 & 2020 \\
\hline
PIB (milliards de \$) & 150 & 360 & 1\,210 & 6\,090 & 14\,720 \\
\hline
\end{tabularx}
\end{center}
\begin{enumerate}
\item Calcule le taux de variation du PIB chinois entre 1978 et 2020 (en \%, arrondi à l'unité).
\item Calcule la part de la croissance réalisée entre 2000 et 2020 dans la croissance totale de 1978 à 2020.
\end{enumerate}
\end{exercice}

\courssub{Je m'entraîne}

\begin{exercice}[Frise chronologique]
Construis une frise chronologique de 1949 à 2020 comportant au moins huit dates clés de l'histoire de la Chine (proclamation de la RPC, Grand Bond en avant, Révolution culturelle, mort de Mao, réformes de Deng, printemps de Tian'anmen, entrée à l'OMC, Jeux olympiques de Pékin). Pour chaque date, indique l'événement en une phrase.
\end{exercice}

\begin{exercice}[Construction graphique]
À partir du tableau de données de l'exercice « Calcul direct » (PIB chinois 1978-2020) :
\begin{enumerate}
\item Construis un diagramme en bâtons représentant l'évolution du PIB chinois (échelle : 1 cm pour 1\,500 milliards de dollars).
\item Rédige un commentaire de cinq lignes : que montre ce graphique sur l'émergence économique chinoise ? À quel tournant historique rattaches-tu le décollage observé ?
\end{enumerate}
\end{exercice}

\begin{exercice}[Les Quatre Modernisations]
\begin{enumerate}
\item Cite les quatre domaines des \cle{Quatre Modernisations} lancées par Deng Xiaoping à partir de 1978.
\item Pour chacun de ces domaines, donne un exemple concret de réforme ou de résultat observé en Chine.
\item Explique en quelques lignes pourquoi ces réformes ont été menées sans remettre en cause le pouvoir du Parti communiste chinois.
\end{enumerate}
\end{exercice}

\begin{exercice}[La Chine et le Cameroun]
\begin{enumerate}
\item Cite trois domaines de la présence chinoise au Cameroun (infrastructures, commerce, financement).
\item Donne deux exemples concrets de réalisations chinoises au Cameroun (routes, barrages, stades, marchés).
\item Rédige un paragraphe de huit à dix lignes expliquant pourquoi cette présence suscite à la fois des espoirs et des inquiétudes.
\end{enumerate}
\end{exercice}

\courssub{Je me prépare au Bac}

\begin{exercice}[Dissertation guidée : une puissance complète ?]
\textbf{Sujet :} « La Chine : une puissance économique complète ? »
\begin{enumerate}
\item Rédige une introduction (accroche, présentation du sujet, problématique, annonce du plan).
\item Première partie : montre les dynamismes de l'économie chinoise (réformes de 1978, ZES et littoral industrialisé, « atelier du monde », entrée à l'OMC en 2001, deuxième économie mondiale).
\item Deuxième partie : montre les limites de cette puissance (inégalités littoral/intérieur, \cle{hukou} et exode rural, pollution, dépendance technologique, absence de libéralisation politique).
\item Rédige une conclusion ouverte : la Chine peut-elle devenir la première puissance mondiale ?
\end{enumerate}
\emph{Barème indicatif : introduction 3 pts ; développement 12 pts ; conclusion 3 pts ; présentation 2 pts.}
\end{exercice}

\begin{exercice}[Étude de document : photographie de Gardes rouges (1966)]
\textbf{Document :} photographie prise à Pékin en 1966, place Tian'anmen. On y voit une foule immense de jeunes en uniforme, brandissant le \emph{Petit Livre rouge} de Mao Zedong, les bras levés, devant un portrait géant du Président Mao.
\begin{enumerate}
\item Décris précisément le document (nature, date, lieu, personnages, objets visibles).
\item Replace cette photographie dans son contexte historique : quel événement majeur débute en Chine en 1966 ?
\item Qui sont les jeunes photographiés ? Quel rôle Mao leur assigne-t-il ?
\item Explique les objectifs de la Révolution culturelle : quels « ennemis » veut-on combattre ?
\item Présente deux conséquences de la Révolution culturelle pour la société et l'économie chinoises (1966-1976).
\end{enumerate}
\emph{Barème indicatif : description 4 pts ; mise en contexte 4 pts ; explication 8 pts ; conséquences 4 pts.}
\end{exercice}

\begin{exercice}[Croquis de synthèse noté et débat argumenté]
\textbf{Partie A — Croquis (12 pts).} Réalise un croquis de synthèse intitulé « Les espaces de l'émergence chinoise ». Tu y feras figurer :
\begin{enumerate}
\item le littoral dynamique (ZES, grandes métropoles : Shanghai, Canton/Guangzhou, Shenzhen, Pékin) ;
\item les espaces intérieurs en retard (campagnes, zones d'exode rural) ;
\item les flux (migrations vers le littoral, exportations vers le monde) ;
\item une légende organisée (au moins quatre items) et une orientation.
\end{enumerate}
\textbf{Partie B — Débat (8 pts).} « La présence chinoise au Cameroun : chance ou risque ? » Après un débat en classe, rédige un paragraphe argumenté d'une quinzaine de lignes dans lequel tu défends ta position à l'aide d'au moins trois arguments étayés par des exemples concrets (infrastructures, endettement, concurrence commerciale, emplois).
\end{exercice}

\newpage

\coursec{Corrigés des exercices}

\coursec{La Chine : de la révolution culturelle à l'émergence économique}
\courssub{Exercices d'application et corrigés détaillés}

\begin{corrige}[Exercice 1 — QCM : dates et acteurs]
\begin{enumerate}
\item \textbf{Réponse b) 1949.} Le 1\ier{} octobre 1949, Mao Zedong proclame la République populaire de Chine (RPC) place Tian'anmen, après la victoire du Parti communiste sur les nationalistes du Guomindang.
\item \textbf{Réponse c) Mao Zedong.} Mao lance la Révolution culturelle en 1966 pour reprendre le pouvoir et éliminer ses rivaux au sein du Parti. Jiang Qing (son épouse) y joue un rôle, mais n'en est pas l'initiatrice.
\item \textbf{Réponse c) 1978.} En décembre 1978, au 3\ieme{} plénum du XI\ieme{} comité central, Deng Xiaoping lance la politique de réforme et d'ouverture (\cle{Quatre Modernisations}).
\item \textbf{Réponse d) 2001.} La Chine entre dans l'OMC le 11 décembre 2001, ce qui consacre son intégration dans l'économie mondiale et accélère ses exportations.
\end{enumerate}
\textbf{Points de vigilance :} ne pas confondre 1949 (proclamation de la RPC) et 1966 (Révolution culturelle) ; ne pas confondre 1976 (mort de Mao) et 1978 (début des réformes).
\end{corrige}

\begin{corrige}[Exercice 2 — Vrai ou faux, à justifier]
\begin{enumerate}
\item \textbf{Faux.} La Révolution culturelle (1966-1976) a désorganisé l'économie : fermeture d'usines et d'universités, envoi de cadres et d'étudiants à la campagne, arrêt de la production. C'est une décennie de chaos économique, pas de renforcement.
\item \textbf{Faux.} Les \cle{Gardes rouges} étaient des jeunes (lycéens et étudiants) mobilisés par Mao à partir de 1966 pour détruire les « Quatre Vieilleries » et dénoncer les « ennemis de la révolution ». Ce ne sont pas des soldats professionnels, même si l'armée (APL) les encadre parfois.
\item \textbf{Vrai.} À partir de 1978, Deng Xiaoping lance les \cle{Quatre Modernisations} : agriculture, industrie, défense, sciences et techniques. C'est le fondement de la politique de réforme et d'ouverture.
\item \textbf{Faux.} Les \cle{ZES} (Shenzhen, Zhuhai, Shantou, Xiamen, créées en 1980) se situent sur le \textbf{littoral} du sud-est de la Chine, face à Hong Kong et Taïwan, afin d'attirer les capitaux étrangers et de faciliter les exportations.
\end{enumerate}
\textbf{Points de vigilance :} la justification est obligatoire ; une réponse « vrai/faux » seule ne vaut aucun point.
\end{corrige}

\begin{corrige}[Exercice 3 — Définitions]
\begin{itemize}
\item \cle{Révolution culturelle} : mouvement politique lancé par Mao Zedong en 1966 pour éliminer ses rivaux et « purifier » la révolution. Menée par les Gardes rouges, elle plonge la Chine dans le chaos (persécutions, destructions, économie désorganisée) jusqu'à la mort de Mao en 1976.
\item \cle{ZES} (zone économique spéciale) : territoire du littoral chinois (Shenzhen, Zhuhai, Shantou, Xiamen, 1980) bénéficiant d'avantages fiscaux et douaniers pour attirer les investissements étrangers et développer les industries d'exportation.
\item \cle{hukou} : système d'enregistrement des ménages qui lie chaque Chinois à son lieu de résidence (rural ou urbain). Il limite l'accès des migrants ruraux aux services urbains (école, santé, logement) et crée une population de migrants sans droits complets dans les villes.
\item \cle{« atelier du monde »} : expression désignant la Chine des années 1990-2010, devenue le premier pays manufacturier grâce à sa main-d'œuvre abondante et bon marché, produisant massivement des biens (textile, électronique) exportés dans le monde entier.
\end{itemize}
\end{corrige}

\begin{corrige}[Exercice 4 — Calcul direct : la croissance du PIB chinois]
\begin{enumerate}
\item Taux de variation entre 1978 et 2020 :
\[ t = \frac{14\,720 - 150}{150} \times 100 = \frac{14\,570}{150} \times 100 \approx 9\,713\,\%. \]
Le PIB chinois a été multiplié par environ 98 (car $14\,720 / 150 \approx 98,1$), soit une hausse d'environ \textbf{9\,713 \%}.
\item Croissance totale 1978-2020 : $14\,720 - 150 = 14\,570$ milliards de dollars. Croissance 2000-2020 : $14\,720 - 1\,210 = 13\,510$ milliards de dollars.
\[ \text{Part} = \frac{13\,510}{14\,570} \times 100 \approx 92,7\,\%. \]
Environ \textbf{92,7 \%} de la croissance totale du PIB entre 1978 et 2020 a été réalisée au cours des seules vingt années 2000-2020, ce qui montre l'accélération spectaculaire de l'émergence après l'entrée dans l'OMC (2001).
\end{enumerate}
\textbf{Points de vigilance :} bien appliquer la formule du taux de variation $\frac{V_A - V_D}{V_D} \times 100$ ; ne pas confondre « multiplié par 98 » et « augmentation de 9\,800 \% » ; arrondir comme demandé.
\end{corrige}

\begin{corrige}[Exercice 5 — Frise chronologique]
Frise attendue (au moins huit dates, dans l'ordre chronologique) :
\begin{itemize}
\item \textbf{1949} : proclamation de la République populaire de Chine par Mao Zedong (1\ier{} octobre).
\item \textbf{1958-1961} : Grand Bond en avant, tentative de modernisation forcée qui provoque une terrible famine (des dizaines de millions de morts).
\item \textbf{1966} : lancement de la Révolution culturelle par Mao ; mobilisation des Gardes rouges.
\item \textbf{1976} : mort de Mao Zedong ; fin de la Révolution culturelle, arrestation de la « bande des Quatre ».
\item \textbf{1978} : Deng Xiaoping lance la politique de réforme et d'ouverture (Quatre Modernisations) ; création des premières ZES en 1980.
\item \textbf{1989} : printemps de Tian'anmen : répression sanglante du mouvement étudiant pro-démocratie (4 juin).
\item \textbf{2001} : entrée de la Chine dans l'OMC (11 décembre).
\item \textbf{2008} : Jeux olympiques de Pékin, vitrine de la puissance chinoise.
\item \textbf{2010} : la Chine devient la deuxième économie mondiale, devant le Japon.
\end{itemize}
\textbf{Barème indicatif :} 1 pt par date correctement placée et légendée (8 pts) ; lisibilité et ordre chronologique (2 pts).
\textbf{Points de vigilance :} respecter l'ordre chronologique ; une phrase par événement ; ne pas inverser 1976 (mort de Mao) et 1978 (réformes).
\end{corrige}

\begin{corrige}[Exercice 6 — Construction graphique]
\begin{enumerate}
\item Hauteurs des bâtons (échelle : 1 cm pour 1\,500 milliards de dollars) :
\begin{center}
\begin{tabularx}{\linewidth}{|l|X|X|X|X|X|}
\hline
\rowcolor{popblueL} Année & 1978 & 1990 & 2000 & 2010 & 2020 \\
\hline
PIB (milliards de \$) & 150 & 360 & 1\,210 & 6\,090 & 14\,720 \\
\hline
Hauteur (cm) & 0,1 & 0,24 & 0,81 & 4,06 & 9,81 \\
\hline
\end{tabularx}
\end{center}
Exemple de calcul : $6\,090 / 1\,500 = 4,06$ cm. Les bâtons sont espacés régulièrement, les axes titrés (années en abscisse, PIB en ordonnée).

\begin{popfigure}
\centering
\begin{tikzpicture}
\begin{axis}[
    ybar,
    bar width=12pt,
    width=\linewidth,
    height=0.6\linewidth,
    xlabel={Années},
    ylabel={PIB (milliards de \$ US)},
    xtick={1978,1990,2000,2010,2020},
    xticklabels={1978,1990,2000,2010,2020},
    ymin=0,
    ymax=16000,
    ymajorgrids=true,
    grid style=dashed,
    nodes near coords,
    nodes near coords align={vertical},
    every node near coord/.append style={font=\tiny, rotate=90, anchor=west, yshift=2pt},
    popcurve/.style={draw=popblue, fill=popblue!60},
]
\addplot[popcurve] coordinates {
    (1978,150)
    (1990,360)
    (2000,1210)
    (2010,6090)
    (2020,14720)
};
\end{axis}
\end{tikzpicture}
\legende{Évolution du PIB chinois (1978-2020) — Source : Banque mondiale, données nominales en milliards de dollars US courants.}
\end{popfigure}

\item \textbf{Commentaire (cinq lignes) :} le graphique montre une croissance d'abord lente (1978-1990), puis de plus en plus rapide, et enfin explosive après 2000 : le PIB est multiplié par 5 entre 2000 et 2010, puis encore par 2,4 entre 2010 et 2020. Cette envolée illustre l'émergence économique chinoise : la Chine devient l'« atelier du monde » puis la deuxième économie mondiale (2010). Le décollage se rattache au tournant des réformes de Deng Xiaoping (1978) : ouverture aux capitaux étrangers, ZES, décollectivisation de l'agriculture. L'accélération après 2000 s'explique par l'entrée dans l'OMC (2001), qui ouvre grand les marchés mondiaux aux exportations chinoises.
\end{enumerate}
\textbf{Points de vigilance :} respecter l'échelle annoncée, légender les axes, donner un titre au graphique ; rattacher explicitement le décollage au tournant de 1978 et à l'OMC.
\end{corrige}

\begin{corrige}[Exercice 7 — Les Quatre Modernisations]
\begin{enumerate}
\item Les quatre domaines : l'\textbf{agriculture}, l'\textbf{industrie}, la \textbf{défense}, les \textbf{sciences et techniques}.
\item Exemples concrets :
\begin{itemize}
\item \emph{Agriculture} : démantèlement des communes populaires et « système de responsabilité familiale » : les paysans louent la terre et vendent leurs surplus, ce qui fait bondir les rendements.
\item \emph{Industrie} : création des ZES (Shenzhen, 1980) attirant les usines étrangères ; développement des industries d'exportation (textile, électronique).
\item \emph{Défense} : modernisation de l'Armée populaire de libération (équipements, professionnalisation, réduction des effectifs).
\item \emph{Sciences et techniques} : réouverture des universités, envoi d'étudiants à l'étranger, programmes spatiaux et informatiques.
\end{itemize}
\item Les réformes sont menées \textbf{sans remettre en cause le pouvoir du PCC} car Deng Xiaoping distingue l'économie de la politique : il libéralise l'économie (« peu importe que le chat soit noir ou blanc, pourvu qu'il attrape les souris ») tout en maintenant le monopole politique du Parti. La répression de Tian'anmen (1989) montre que toute contestation politique est refusée : la modernisation économique doit légitimer le Parti, non l'affaiblir.
\end{enumerate}
\textbf{Points de vigilance :} citer les quatre domaines exactement ; bien distinguer libéralisation économique (réelle) et libéralisation politique (refusée).
\end{corrige}

\begin{corrige}[Exercice 8 — La Chine et le Cameroun]
\begin{enumerate}
\item Trois domaines : les \textbf{infrastructures} (routes, barrages, stades, ports), le \textbf{commerce} (importations de produits manufacturés chinois, marchés, boutiques), le \textbf{financement} (prêts et investissements chinois pour les grands projets).
\item Exemples concrets : le barrage hydroélectrique de Memve'ele (Sud), le port en eau profonde de Kribi, les stades de football (stade de Bafoussam, stade de Limbé), les marchés et commerces chinois à Douala et Yaoundé, la route Yaoundé-Douala (tronçons financés/réalisés avec des partenaires chinois).
\item \textbf{Paragraphe (8 à 10 lignes) :} la présence chinoise au Cameroun suscite d'abord de grands espoirs : elle finance et réalise rapidement des infrastructures (barrages, ports, routes, stades) que le pays n'aurait pas pu construire seul, créant des emplois et dynamisant l'économie ; les produits chinois bon marché améliorent le pouvoir d'achat des ménages modestes. Mais elle nourrit aussi des inquiétudes : l'endettement croissant du Cameroun envers la Chine menace la soutenabilité de la dette ; la concurrence des produits importés étouffe l'artisanat et l'industrie locaux ; les entreprises chinoises emploient souvent leur propre main-d'œuvre et respectent mal les normes sociales et environnementales. La présence chinoise est donc ambivalente : une chance si le Cameroun négocie des partenariats équilibrés, un risque de dépendance si les projets ne profitent pas durablement au développement national.
\end{enumerate}
\textbf{Points de vigilance :} citer des exemples précis et localisés ; équilibrer espoirs et inquiétudes ; conclure par une prise de position nuancée.
\end{corrige}

\begin{corrige}[Exercice 9 — Dissertation guidée : une puissance complète ?]
\textbf{Introduction (3 pts).} \emph{Accroche :} en 2010, la Chine devient la deuxième économie mondiale, soixante ans après la proclamation de la RPC. \emph{Présentation :} depuis les réformes de Deng Xiaoping (1978), la Chine connaît une croissance exceptionnelle qui en fait l'« atelier du monde ». \emph{Problématique :} cette puissance économique est-elle complète et équilibrée ? \emph{Annonce du plan :} nous montrerons d'abord les dynamismes de l'économie chinoise, puis ses limites.

\textbf{Première partie : les dynamismes (6 pts).} Les réformes de 1978 (Quatre Modernisations, décollectivisation agricole) libèrent les forces productives. Les ZES (Shenzhen, 1980) et le littoral industrialisé attirent les capitaux étrangers. La Chine devient l'« atelier du monde » : premier exportateur mondial de biens manufacturés. L'entrée dans l'OMC (2001) accélère les exportations ; le PIB passe de 150 milliards de dollars (1978) à 14\,720 milliards (2020). La Chine est la deuxième économie mondiale (2010), accumule des réserves de change et investit à l'étranger (routes de la soie, Afrique).

\textbf{Deuxième partie : les limites (6 pts).} Les inégalités s'aggravent entre le littoral riche et l'intérieur rural pauvre ; le \cle{hukou} prive les 250 millions de migrants ruraux de droits complets dans les villes. La pollution (air, eau, sols) atteint des niveaux record. La Chine reste dépendante des technologies étrangères (composants, logiciels) et de ses marchés d'exportation. Enfin, l'absence de libéralisation politique (parti unique, censure, répression de Tian'anmen en 1989) fragilise le modèle.

\textbf{Conclusion (3 pts).} La Chine est une puissance économique majeure mais incomplète : riche mais inégale, industrielle mais polluée et dépendante. Elle peut devenir la première puissance mondiale si elle résout ses contradictions sociales, environnementales et technologiques — question ouverte pour le XXI\ieme{} siècle.

\textbf{Points de vigilance :} problématique explicite ; plan en deux parties équilibrées ; dates et chiffres précis ; éviter le hors-sujet politique général.
\end{corrige}

\begin{corrige}[Exercice 10 — Étude de document : photographie de Gardes rouges (1966)]
\begin{enumerate}
\item \textbf{Description (4 pts) :} il s'agit d'une photographie prise à Pékin, place Tian'anmen, en 1966. On y voit une foule immense de jeunes en uniforme, bras levés, brandissant le \emph{Petit Livre rouge} (recueil de citations de Mao), devant un portrait géant du Président Mao Zedong.
\item \textbf{Contexte (4 pts) :} en 1966, Mao Zedong lance la \cle{Révolution culturelle} (« Grande Révolution culturelle prolétarienne ») pour reprendre le pouvoir perdu après l'échec du Grand Bond en avant (1958-1961) et éliminer ses rivaux au sein du Parti.
\item \textbf{Les jeunes photographiés (4 pts) :} ce sont les \cle{Gardes rouges}, lycéens et étudiants fanatisés. Mao leur assigne le rôle d'avant-garde de la révolution : dénoncer les « révisionnistes », attaquer les cadres du Parti jugés « bourgeois » et détruire les « Quatre Vieilleries » (vieilles idées, vieille culture, vieilles coutumes, vieilles habitudes).
\item \textbf{Objectifs (4 pts) :} la Révolution culturelle veut combattre les « ennemis de classe » : intellectuels, enseignants, cadres du Parti accusés de prendre la « voie capitaliste », ainsi que tout héritage traditionnel ou occidental. Il s'agit de « purifier » la révolution et de consolider le culte de la personnalité de Mao.
\item \textbf{Conséquences (4 pts) :} (a) \emph{Société} : des centaines de milliers de persécutés (humiliations, emprisonnements, morts), fermeture des écoles et universités, envoi de millions de jeunes « éduqués » à la campagne, destruction du patrimoine culturel. (b) \emph{Économie} : production désorganisée, usines perturbées, cadres compétents écartés : une décennie perdue pour le développement chinois (1966-1976).
\end{enumerate}
\textbf{Points de vigilance :} distinguer description (ce qu'on voit) et interprétation (ce qu'on sait) ; dater précisément ; citer les Gardes rouges et le \emph{Petit Livre rouge} ; donner exactement deux conséquences.
\end{corrige}

\begin{corrige}[Exercice 11 — Croquis de synthèse noté et débat argumenté]
\textbf{Partie A — Croquis (12 pts).} Attendus :
\begin{itemize}
\item \textbf{Titre} : « Les espaces de l'émergence chinoise » (1 pt).
\item \textbf{Littoral dynamique} (4 pts) : zone colorée le long des côtes est et sud ; localisation des ZES (Shenzhen) et des métropoles (Pékin, Shanghai, Canton/Guangzhou) par des points nommés.
\item \textbf{Espaces intérieurs} (2 pts) : zone différemment colorée à l'ouest et au centre (campagnes, zones d'exode rural).
\item \textbf{Flux} (2 pts) : flèches de migration de l'intérieur vers le littoral ; flèches d'exportation du littoral vers le Pacifique/le monde.
\item \textbf{Légende et orientation} (3 pts) : légende organisée (au moins quatre items : littoral industrialisé, intérieur rural, flux migratoires, flux d'exportation), flèche du nord, échelle indicative.
\end{itemize}

\begin{popfigure}
\centering
\begin{tikzpicture}[scale=0.85]
% Contour très simplifié de la Chine (polygone approximatif)
\draw[popdark, thick] (0,0) -- (1,3) -- (3,4) -- (5,3.5) -- (7,3) -- (8,1.5) -- (7.5,0) -- (6,-1) -- (4,-1.5) -- (2,-1) -- (0.5,-0.5) -- cycle;

% Littoral dynamique (Est/Sud-Est)
\fill[popblue!20] (3,4) -- (5,3.5) -- (7,3) -- (8,1.5) -- (7.5,0) -- (6,0.5) -- (4,1) -- (2,1.5) -- cycle;

% Intérieur rural (Ouest)
\fill[poporange!20] (0,0) -- (1,3) -- (3,4) -- (2,1.5) -- (4,1) -- (6,0.5) -- (7.5,0) -- (6,-1) -- (4,-1.5) -- (2,-1) -- (0.5,-0.5) -- cycle;

% Villes principales
\node[circle, fill=popblue, inner sep=1.5pt, label=above:\tiny Pékin] (P) at (2.5, 3.5) {};
\node[circle, fill=popblue, inner sep=1.5pt, label=right:\tiny Shanghai] (S) at (5.5, 2.8) {};
\node[circle, fill=popblue, inner sep=1.5pt, label=below right:\tiny Canton] (G) at (6.5, 1) {};
\node[circle, fill=poporange, inner sep=1.5pt, label=below:\tiny Shenzhen (ZES)] (Sz) at (7.2, 0.5) {};

% Flèches migration (Intérieur -> Littoral)
\draw[poporange, thick, -{Stealth[length=3mm]}] (1.5, 0.5) -- (3.5, 1.8);
\draw[poporange, thick, -{Stealth[length=3mm]}] (3, -0.5) -- (5, 1);
\draw[poporange, thick, -{Stealth[length=3mm]}] (5, -1) -- (6.5, 0.2);

% Flèches exportation (Littoral -> Mer/Monde)
\draw[popblue, thick, -{Stealth[length=3mm]}] (7.5, 2) -- (9, 2.5);
\draw[popblue, thick, -{Stealth[length=3mm]}] (7, 0.5) -- (9, 0);

% Légende
\begin{scope}[shift={(9.5, 3)}]
\node[anchor=west, font=\small\bfseries] at (0,0) {Légende};
\fill[popblue!20] (0,-0.4) rectangle (0.5,-0.1); \node[anchor=west, font=\small] at (0.6,-0.25) {Littoral industrialisé};
\fill[poporange!20] (0,-0.7) rectangle (0.5,-0.4); \node[anchor=west, font=\small] at (0.6,-0.55) {Intérieur rural};
\draw[poporange, thick, -{Stealth[length=2mm]}] (0,-1.0) -- (0.5,-1.0); \node[anchor=west, font=\small] at (0.6,-1.05) {Migrations internes};
\draw[popblue, thick, -{Stealth[length=2mm]}] (0,-1.3) -- (0.5,-1.3); \node[anchor=west, font=\small] at (0.6,-1.35) {Exportations};
\fill[popblue] (0,-1.6) circle (1.5pt); \node[anchor=west, font=\small] at (0.6,-1.65) {Métropoles / ZES};
\end{scope}

% Flèche Nord
\begin{scope}[shift={(0.5, 3.5)}]
\draw[popdark, thick, -{Stealth[length=3mm]}] (0,0) -- (0,0.5);
\node[font=\small\bfseries] at (0,0.7) {N};
\end{scope}

% Échelle indicative
\node[font=\tiny] at (0.5, -1.8) {Échelle indicative : $\approx$ 1 cm $\simeq$ 1000 km};
\end{tikzpicture}
\legende{Croquis de synthèse : Les espaces de l'émergence chinoise (modèle corrigé).}
\end{popfigure}

\textbf{Points de vigilance :} pas de carte muette reproduite : un croquis est un schéma simplifié et raisonné ; tout élément figuré doit apparaître dans la légende ; la Chine doit être reconnaissable (contour simplifié, littoral à l'est).

\textbf{Partie B — Débat (8 pts).} Exemple de paragraphe argumenté (position « chance, à condition de partenariats équilibrés ») : la présence chinoise au Cameroun est d'abord une chance. Premier argument : elle réalise des infrastructures majeures que le pays attendait depuis longtemps — le port en eau profonde de Kribi et le barrage de Memve'ele renforcent les capacités logistiques et énergétiques nationales. Deuxième argument : ces chantiers créent des emplois et dynamisent l'activité locale, tandis que les produits chinois bon marché améliorent le pouvoir d'achat des ménages. Troisième argument : la Chine offre un partenaire alternatif, ce qui diversifie les relations extérieures du Cameroun. Cependant, des risques existent : l'endettement envers la Chine s'alourdit, la concurrence des importations asphyxie l'artisanat local (textile, chaussures), et certaines entreprises chinoises emploient peu de main-d'œuvre camerounaise. En définitive, cette présence sera une chance si l'État camerounais négocie des contrats transparents, impose des transferts de compétences et maîtrise sa dette ; sinon, elle deviendra un risque de dépendance durable. (Barème : position claire 2 pts ; trois arguments étayés 4 pts ; exemples concrets et expression 2 pts.)
\end{corrige}

\newpage

\coursec{Fiche bilan}

\courssub{Carte mentale de synthèse}
\begin{popfigure}
\begin{tikzpicture}[
  mindmap,
  grow cyclic,
  every node/.style={concept, execute at begin node=\small\bfseries},
  concept color=popblue,
  level 1/.append style={level distance=4.5cm, sibling angle=72, font=\small\bfseries},
  level 2/.append style={level distance=3cm, sibling angle=40, font=\footnotesize},
  text width=2.8cm, align=center
]
\node[concept, minimum size=2.2cm, align=center]{Chine\\R\'evolution culturelle \rightarrow \'Emergence}
  child[concept color=popblue] { node[concept, align=center]{1. Mao \& R\'evolution\\culturelle (1966--1976)}
    child { node {Gardes rouges} }
    child { node {Purges \& chaos} }
    child { node {Culte de la personnalit\'e} }
  }
  child[concept color=poporange] { node[concept, align=center]{2. Deng Xiaoping\\R\'eformes (1978--1989)}
    child { node {Quatre modernisations} }
    child { node {ZES (Shenzhen)} }
    child { node {Responsabilit\'e familiale} }
    child { node {Tian'anmen 1989} }
  }
  child[concept color=popgreen] { node[concept, align=center]{3. \'Emergence\\``Atelier du monde'' (1990--2010)}
    child { node {Adh\'esion OMC 2001} }
    child { node {Croissance \textasciitilde 10\%/an} }
    child { node {Exportations massives} }
    child { node {R\'eserves de change} }
  }
  child[concept color=poppurple] { node[concept, align=center]{4. Limites \&\\Contradictions}
    child { node {In\'egalit\'es c\^oti\`eres/int\'erieur} }
    child { node {Pollution \& ressources} }
    child { node {Vieillissement d\'emographique} }
    child { node {Contr\^ole politique unique} }
  }
  child[concept color=poppink] { node[concept, align=center]{5. Chine -- Afrique\\Cameroun}
    child { node {Infrastructures (routes, barrages)} }
    child { node {Mati\`eres premi\`eres} }
    child { node {FOCAC / Nouvelles routes de la soie} }
    child { node {Dette \& souverainet\'e} }
  };
\end{tikzpicture}
\legende{Carte mentale : les 5 axes majeurs de la le\c{c}on sur la Chine contemporaine.}
\end{popfigure}

\begin{bilan}
\textbf{D\'efinitions cl\'es \`a conna\^itre par c\oe{}ur}
\begin{itemize}
  \item \cle{R\'evolution culturelle (1966-1976)} : Mouvement politique massif lanc\'e par Mao Zedong pour ``purger'' le Parti des ``r\'evisionnistes'' et r\'eaffirmer son autorit\'e ; mobilisation des \cle{Gardes rouges}, chaos social, destruction du patrimoine, arr\^et de l'\'economie.
  \item \cle{Quatre modernisations} (Zhou Enlai 1975, relanc\'ees par Deng 1978) : Agriculture, Industrie, D\'efense, Sciences/Technologies. Objectif : rattraper le retard \'economique.
  \item \cle{Zones \'economiques sp\'eciales (ZES)} : Territoires (Shenzhen, Zhuhai, Shantou, Xiamen) b\'en\'eficiant de lois d\'erogatoires (incitations fiscales, autonomie de gestion) pour attirer les \cle{IDE} (Investissements Directs \'Etrangers).
  \item \cle{\'Economie socialiste de march\'e} : Concept officiel depuis 1992 (14e Congr\`es PCC) : maintien du pouvoir politique du PCC + m\'ecanismes de march\'e (prix, concurrence, propri\'et\'e priv\'ee encadr\'ee).
  \item \cle{Atelier du monde} : Sp\'ecialisation de la Chine dans l'assemblage et l'exportation de biens manufactur\'es \`a bas co\^ut (main-d'\oe{}uvre abondante, infrastructure, taux de change avantageux).
  \item \cle{FOCAC} (Forum sur la Coop\'eration Sino-Africaine, depuis 2000) : Cadre institutionnel diplomatique, \'economique et financier structurant les relations Chine-Afrique (pr\^ets, dons, infrastructures, formation).
  \item \cle{BRI} (Belt and Road Initiative, Nouvelles Routes de la Soie, 2013) : Strat\'egie mondiale d'infrastructures et d'investissements (routes, ports, chemins de fer, corridors num\'eriques) pour \'etendre l'influence \'economique et g\'eopolitique chinoise.
\end{itemize}

\textbf{Rep\`eres chronologiques \& Acteurs majeurs}
\begin{center}
\begin{tabularx}{\linewidth}{|c|X|X|}\hline
\textbf{P\'eriode / Date} & \textbf{\'Ev\'enement / Acteur cl\'e} & \textbf{Signification} \\ \hline
1966--1976 & Mao Zedong / Gardes rouges & R\'evolution culturelle : chaos politique et social \\ \hline
1976 & Mort de Mao / Arrestation de la \cle{Bande des Quatre} & Fin de la R\'evolution culturelle \\ \hline
1978 & IIIe Pl\'enum du 11e Comit\'e central (Deng Xiaoping) & Tournant r\'eformateur : ``La pratique est le seul crit\`ere de la v\'erit\'e'' \\ \hline
1979 & Instauration politique de \cle{l'enfant unique} & Contr\^ole d\'emographique (assoupli 2016 : 2 enfants, 2021 : 3) \\ \hline
1980 & Cr\'eation des 4 premi\`eres \cle{ZES} (dont Shenzhen) & Ouverture aux IDE, laboratoire du capitalisme d'\'Etat \\ \hline
1989 & R\'epression de Tian'anmen & Maintien du monopole politique du PCC, arr\^et lib\'eralisation politique \\ \hline
1992 & Tourn\'ee du Sud de Deng / 14e Congr\`es PCC & ``\'Economie socialiste de march\'e'' officialis\'ee \\ \hline
2001 & Adh\'esion \`a l'\cle{OMC} & Int\'egration pleine dans la mondialisation, boom exportations \\ \hline
2000--auj. & Sommets \cle{FOCAC} (P\'ekin, Johannesburg, Dakar) & Partenariat strat\'egique Chine-Afrique (infrastructures, ressources) \\ \hline
2013 & Lancement \cle{Nouvelles Routes de la Soie} (BRI) & Projection g\'eopolitique et \'economique mondiale \\ \hline
\end{tabularx}
\end{center}

\textbf{M\'ethode express : Commentaire de document historique (Bac)}
\begin{enumerate}
  \item \textbf{Identifier} : Nature (discours, stats, caricature, photo), Auteur, Date, Contexte imm\'ediat.
  \item \textbf{Analyser} : Th\`ese centrale / message ; arguments ou donn\'ees chiffr\'ees ; vocabulaire cl\'e (id\'eologie, propagande, r\'eforme).
  \item \textbf{Contextualiser} : Rattacher au plan de cours (ex: ZES $\rightarrow$ Deng ; Gardes rouges $\rightarrow$ Mao ; FOCAC $\rightarrow$ \'emergence diplomatique).
  \item \textbf{Critiquer / Porter} : Fiabilit\'e (source partisane ?) ; Port\'ee (effets r\'eels, limites, cons\'equences \`a long terme).
  \item \textbf{Conclure} : R\'eponse synth\'etique \`a la probl\'ematique pos\'ee.
\end{enumerate}
\end{bilan}

\begin{aretenir}[Checklist avant le Bac]
$\square$ Je sais d\'efinir \cle{R\'evolution culturelle}, \cle{Quatre modernisations}, \cle{ZES}, \cle{\'Economie socialiste de march\'e}, \cle{Atelier du monde}, \cle{FOCAC}, \cle{BRI}.
$\square$ Je place les 5 grandes dates : 1966 (d\'ebut RC), 1976 (fin RC), 1978 (tournant Deng), 1989 (Tian'anmen), 2001 (OMC).
$\square$ J'explique les \textbf{causes} de la R\'evolution culturelle (lutte de pouvoir Mao vs Liu/Deng, peur r\'evisionnisme) et ses \textbf{cons\'equences} (g\'en\'eration perdue, \'economie paralys\'ee, discr\'edit mao\"\i{}sme).
$\square$ Je d\'etaille la strat\'egie de \cle{Deng Xiaoping} : pragmatisme (``peu importe la couleur du chat''), ouverture progressive (ZES $\rightarrow$ c\^oti\`ere $\rightarrow$ int\'erieure), primaut\'e de l'\'economie sur la politique.
$\square$ Je caract\'erise le mod\`ele de croissance 1990-2010 : investissement/exportation, main-d'\oe{}uvre migrante bas co\^ut, \'Etat strat\`ege (banques, SOE), adh\'esion OMC.
$\square$ Je cite 3 \textbf{limites structurelles} : in\'egalit\'es spatiales/sociales (coefficient de Gini), vieillissement/pyramide des \^ages, mod\`ele extensif non durable (environnement, dette locale).
$\square$ Je d\'ecris la relation \textbf{Chine-Cameroun} : infrastructures (barrages de M\'ekin et Memve'ele, routes, stade Olemb\'e), ressources (bois, coton, minerais), dette, formation, FOCAC/BRI.
$\square$ Je sais analyser un document type : discours de Deng (1978/1992), stats commerce/PIB, carte ZES/flux, photo Gardes rouges, caricature Tian'anmen, accord Chine-Cameroun.
$\square$ Je ma\^itrise le vocabulaire sp\'ecifique : \cle{PCC}, \cle{Comit\'e central}, \cle{Bureau politique}, \cle{SOE} (Entreprises d'\'Etat), \cle{IDE}, \cle{BRI}, \cle{Pouvoir doux}.
$\square$ Je sais r\'ediger une introduction (accroche, d\'efinition, probl\'ematique, plan) et une conclusion (bilan, ouverture) structur\'ees pour une dissertation ou un commentaire.
$\square$ Je distingue \textbf{continu\'it\'e} (PCC au pouvoir, planification indicative, contr\^ole id\'eologique) et \textbf{rupture} (march\'e, ouverture, in\'egalit\'es, soci\'et\'e de consommation) 1949 / 1978 / aujourd'hui.
\end{aretenir}

\findecours

\end{document}
